% Lesson 45 — Vocabulary: Words and expressions related to polls/surveys — Anglais Tle A — Cours signé OnBuch+
% Programme officiel MINESEC. Compiler avec : tectonic EN45-vocabulary-words-and-expressions-related-to-polls-surve.tex
\newcommand{\DOCMATIERE}{Anglais}
\newcommand{\DOCNIVEAU}{Tle A}
\newcommand{\DOCMODULE}{Chapter 4 : Citizenship and Human Rights}
\newcommand{\DOCLECON}{EN45}
\newcommand{\DOCTITRE}{Lesson 45 — Vocabulary: Words and expressions related to polls/surveys}
% OnBuch+ — préambule « cours signés » ENRICHI (compilé avec tectonic / XeLaTeX)
% Système visuel « Pop » d'OnBuch+ : fond crème (jamais blanc), bordures encre
% épaisses, ombres « dures » décalées, titres Archivo Black en capitales.
% Version MATHÉMATIQUES : graphes (pgfplots/tikz), boîtes pédagogiques
% (définition, propriété, méthode, attention, le savais-tu, exercice résolu,
% exercice, TP, bilan), page de garde et signature OnBuch+ ancrée en bas.
%
% Macros attendues dans main.tex AVANT \input{preamble} :
%   \DOCMATIERE  \DOCNIVEAU  \DOCMODULE  \DOCLECON (numéro)  \DOCTITRE
\documentclass[11pt]{article}

\usepackage{fontspec}
\usepackage[english,french]{babel} % français = langue principale ; l'anglais via \eng{..} / textebox
\usepackage[a4paper,margin=2cm,top=3.4cm,bottom=2.8cm,headheight=1.4cm,footskip=1.3cm]{geometry}
\usepackage{amsmath,amssymb,amsfonts}
\usepackage{mathtools} % \xmapsto, \coloneqq... souvent utilisés par les modèles
\usepackage{cancel} % simplifications barrées (bilans de forces)
% \abs{z} (module, valeur absolue) et \norm{u} (norme), utilisés par les modèles
\providecommand{\abs}{}\let\abs\relax\DeclarePairedDelimiter{\abs}{\lvert}{\rvert}
\providecommand{\norm}{}\let\norm\relax\DeclarePairedDelimiter{\norm}{\lVert}{\rVert}
\usepackage[table]{xcolor} % [table] : fournit aussi \rowcolors (alternance de lignes)
\usepackage{pagecolor}
\usepackage{tikz}
\usetikzlibrary{arrows.meta,positioning,calc,shapes.geometric,shapes.misc,shapes.arrows,decorations.pathmorphing,decorations.pathreplacing,decorations.markings,patterns,shadows,mindmap,trees,fit,backgrounds,matrix,intersections,angles,quotes}
\usepackage{pgfplots}
\usepackage[version=4]{mhchem} % équations nucléaires \ce{^{235}_{92}U}
\usepackage[european,siunitx]{circuitikz} % schémas électriques
\pgfplotsset{compat=1.18}
\usepgfplotslibrary{fillbetween,groupplots} % aires entre courbes (calcul intégral)
\usepackage{enumitem}
\usepackage{fancyhdr}
% [most] charge toutes les bibliothèques tcolorbox (listings, minted, raster,
% documentation...) et gonfle inutilement la mémoire TeX sur un document long
% (~300 boîtes) : on ne charge que ce qui est réellement utilisé.
\usepackage[skins,breakable]{tcolorbox}
\usepackage{graphicx}
\usepackage{colortbl}
\usepackage{booktabs}
\usepackage{tabularx}
\usepackage{diagbox} % case d'en-tête diagonale des tableaux à double entrée
% Colonnes extensibles centrée / gauche / droite (C, L, R), souvent utilisées par les modèles
\newcolumntype{C}{>{\centering\arraybackslash}X}
\newcolumntype{L}{>{\raggedright\arraybackslash}X}
\newcolumntype{R}{>{\raggedleft\arraybackslash}X}
\usepackage{array}
\usepackage{multicol}
\usepackage{multirow}
\usepackage{siunitx}
% parse-numbers=false : accepte aussi \SI{2,5 \times 10^{-3}}{...}
\sisetup{locale=FR,output-decimal-marker={,},group-minimum-digits=4,parse-numbers=false}
\DeclareSIUnit\ohm{\ensuremath{\symup{\Omega}}} % U+2126 absent de la police
\DeclareSIUnit\division{div} % réglages d'oscilloscope (ms/div, V/div)
\DeclareSIUnit\tr{tr} % tours (vitesse de rotation, ex. \SI{1500}{\tr\per\minute})
\DeclareSIUnit\molar{M} % concentration molaire (ex. \SI{3}{\molar}, \SI{1}{\micro\molar})
\DeclareSIUnit\an{an} % année (le modèle confond parfois avec \year, un compteur TeX) — ex. \SI{5}{\kilo\meter\per\an}
\DeclareSIUnit\FCFA{FCFA} % franc CFA (ex. \SI{500}{\FCFA\per\kilo\gram})
\DeclareSIUnit\degreeBrix{\textdegree Brix} % teneur en sucre d'un sirop/jus (réfractométrie)
\DeclareSIUnit\month{mois} % le modèle utilise parfois \month, un compteur TeX, comme unité de durée
\usepackage{qrcode}
% \degree : commande fréquemment utilisée par les modèles pour un angle ou une
% température en dehors de \SI{}{} (non fournie par les paquets chargés).
\providecommand{\degree}{^{\circ}}
% \qed (fin de démonstration), utilisé par les modèles sans amsthm
\providecommand{\qed}{\hfill$\square$}
% \barre : double barre des valeurs interdites dans un tableau de variations (tkz-tab)
\providecommand{\barre}{\ensuremath{\|}}
% environnement proof (amsthm non chargé) utilisé par les modèles
\ifdefined\proof\else\newenvironment{proof}[1][Démonstration]{\par\noindent\textit{#1.}\ }{\qed\par}\fi
\usepackage{lastpage}
\usepackage{hyperref}
\hypersetup{hidelinks}

% ============================================================
% Palette « Pop » OnBuch+ — mêmes hex que la classe P (plus_ui.dart)
% ============================================================
\definecolor{poporange}{HTML}{F59321}
\definecolor{popdark}{HTML}{E07A0C}
\definecolor{popcream}{HTML}{FFF6E8}
\definecolor{popink}{HTML}{1C1714}
\definecolor{poppurple}{HTML}{7A5AE0}
\definecolor{popgreen}{HTML}{2E9E6A}
\definecolor{poppink}{HTML}{E0457B}
\definecolor{popblue}{HTML}{2F7DE1}
\definecolor{popgold}{HTML}{C98A12}
\definecolor{popmuted}{HTML}{8A7F76}
\definecolor{popline}{HTML}{EADFD2}
\definecolor{popgray}{HTML}{8A7F76}
\colorlet{popgrey}{popgray}
% Alias tolérants (le modèle nomme parfois les couleurs différemment)
\colorlet{popwhite}{white}
\colorlet{popblack}{popink}
\colorlet{popred}{poppink}
\colorlet{popyellow}{popgold}
% Teintes claires (fonds de tableaux/encadrés) — le modèle nomme parfois
% les couleurs avec un suffixe L (ex. popblueL) sans les avoir définies.
\colorlet{poporangeL}{poporange!20}
\colorlet{popdarkL}{popdark!20}
\colorlet{poppurpleL}{poppurple!20}
\colorlet{popgreenL}{popgreen!20}
\colorlet{poppinkL}{poppink!20}
\colorlet{popblueL}{popblue!20}
\colorlet{popgoldL}{popgold!20}
\colorlet{popredL}{popred!20}
\colorlet{popgrayL}{popgray!20}
% Teintes claires pour fonds de tableaux / zones
\colorlet{popblueL}{popblue!10!white}
\colorlet{poporangeL}{poporange!14!white}
\colorlet{popgreenL}{popgreen!12!white}
\colorlet{poppurpleL}{poppurple!12!white}
\colorlet{poppinkL}{poppink!12!white}
\colorlet{popgoldL}{popgold!14!white}

\pagecolor{popcream}

% ============================================================
% Typographie
% ============================================================
\defaultfontfeatures{Ligatures=TeX}
\setmainfont{PlusJakartaSans-Medium.ttf}[Path=../../../fonts/,BoldFont=PlusJakartaSans-Bold.ttf]
% Symboles, API et emojis absents de la police principale : repli sur DejaVu Sans (caractères actifs)
\newfontface\symfont{DejaVuSans.ttf}[Path=../../../fonts/]
\makeatletter
\newcommand\symactive[1]{\catcode#1=13 \begingroup\lccode`\~=#1 \lowercase{\endgroup\def~}{{\symfont\char#1}}}
\newcommand\symblank[1]{\catcode#1=13 \begingroup\lccode`\~=#1 \lowercase{\endgroup\def~}{}}
\newcommand\symhyph[1]{\catcode#1=13 \begingroup\lccode`\~=#1 \lowercase{\endgroup\def~}{\mbox{-}}}
\newcount\symc
\newcommand\symrange[2]{\symc=#1 \@whilenum\symc<#2 \do{\expandafter\symactive\expandafter{\the\symc}\advance\symc by 1 }}
\newcommand\symblankrange[2]{\symc=#1 \@whilenum\symc<#2 \do{\expandafter\symblank\expandafter{\the\symc}\advance\symc by 1 }}
\makeatother
\symrange{"0250}{"02FF}\symactive{"2713}\symactive{"2714}\symactive{"2717}\symactive{"2718}\symactive{"2610}\symactive{"2611}\symactive{"2612}
\symhyph{"2011}\symhyph{"2E1A}\symblankrange{"1F300}{"1FAFF}\symblank{"FE0F}
% Maths en sans-serif (Fira Math) avec chiffres et lettres droites de Plus
% Jakarta, pour que les formules se fondent dans le texte.
\usepackage[math-style=ISO,bold-style=ISO]{unicode-math}
\setmathfont{FiraMath-Regular.otf}[Path=../../../fonts/]
\setmathfont{PlusJakartaSans-Medium.ttf}[Path=../../../fonts/,range={up/num,up/{Latin,latin},\mathup}]
\usepackage{icomma} % virgule décimale sans espace en mode math
% Caractères mathématiques Unicode absents des polices de texte (Plus Jakarta,
% Archivo Black) : rendus par leur équivalent mathématique, en texte comme en
% formule — sinon ils s'affichent en carré vide (ex. « Arithmétique dans ℤ »).
\usepackage{newunicodechar}
\newunicodechar{ℝ}{\ensuremath{\mathbb{R}}}
\newunicodechar{ℤ}{\ensuremath{\mathbb{Z}}}
\newunicodechar{ℂ}{\ensuremath{\mathbb{C}}}
\newunicodechar{ℕ}{\ensuremath{\mathbb{N}}}
\newunicodechar{ℚ}{\ensuremath{\mathbb{Q}}}
\newunicodechar{𝔻}{\ensuremath{\mathbb{D}}}
\newunicodechar{α}{\ensuremath{\alpha}}
\newunicodechar{β}{\ensuremath{\beta}}
\newunicodechar{γ}{\ensuremath{\gamma}}
\newunicodechar{δ}{\ensuremath{\delta}}
\newunicodechar{ε}{\ensuremath{\varepsilon}}
\newunicodechar{θ}{\ensuremath{\theta}}
\newunicodechar{λ}{\ensuremath{\lambda}}
\newunicodechar{μ}{\ensuremath{\mu}}
\newunicodechar{σ}{\ensuremath{\sigma}}
\newunicodechar{τ}{\ensuremath{\tau}}
\newunicodechar{φ}{\ensuremath{\varphi}}
\newunicodechar{ψ}{\ensuremath{\psi}}
\newunicodechar{ω}{\ensuremath{\omega}}
\newunicodechar{Σ}{\ensuremath{\Sigma}}
% majuscules produites par \MakeUppercase dans les titres (α-aminés -> Α)
\newunicodechar{Α}{\ensuremath{\alpha}}
\newunicodechar{Β}{\ensuremath{\beta}}
\newunicodechar{Γ}{\ensuremath{\Gamma}}
\newunicodechar{Δ}{\ensuremath{\Delta}}
\newunicodechar{Ε}{\ensuremath{\varepsilon}}
\newunicodechar{Θ}{\ensuremath{\Theta}}
\newunicodechar{Λ}{\ensuremath{\Lambda}}
\newunicodechar{Μ}{\ensuremath{\mu}}
\newunicodechar{Τ}{\ensuremath{\tau}}
\newunicodechar{Φ}{\ensuremath{\Phi}}
\newunicodechar{Ψ}{\ensuremath{\Psi}}
\newunicodechar{Ω}{\ensuremath{\Omega}}
\newunicodechar{↦}{\ensuremath{\mapsto}}
\newunicodechar{∈}{\ensuremath{\in}}
\newunicodechar{∉}{\ensuremath{\notin}}
\newunicodechar{∧}{\ensuremath{\wedge}}
\newunicodechar{⇔}{\ensuremath{\Leftrightarrow}}
\newunicodechar{⟺}{\ensuremath{\Longleftrightarrow}}
\newunicodechar{⇒}{\ensuremath{\Rightarrow}}
\newunicodechar{⟹}{\ensuremath{\Longrightarrow}}
\newunicodechar{∘}{\ensuremath{\circ}}
\newunicodechar{∪}{\ensuremath{\cup}}
\newunicodechar{∩}{\ensuremath{\cap}}
\newunicodechar{≡}{\ensuremath{\equiv}}
\newunicodechar{⊂}{\ensuremath{\subset}}
\newunicodechar{⊥}{\ensuremath{\perp}}
\newunicodechar{∃}{\ensuremath{\exists}}
\newunicodechar{∀}{\ensuremath{\forall}}
\newunicodechar{∛}{\ensuremath{\sqrt[3]{\,}}}
\newunicodechar{∜}{\ensuremath{\sqrt[4]{\,}}}
\newunicodechar{⃗}{\ensuremath{{}^{\to}}}
\newunicodechar{ⁿ}{\ifmmode^{n}\else\textsuperscript{\ensuremath{n}}\fi}
\newunicodechar{ˣ}{\ifmmode^{x}\else\textsuperscript{\ensuremath{x}}\fi}
\newunicodechar{ᵇ}{\ifmmode^{b}\else\textsuperscript{\ensuremath{b}}\fi}
\newunicodechar{ʳ}{\ifmmode^{r}\else\textsuperscript{\ensuremath{r}}\fi}
\newunicodechar{ᵘ}{\ifmmode^{u}\else\textsuperscript{\ensuremath{u}}\fi}
\newunicodechar{ʸ}{\ifmmode^{y}\else\textsuperscript{\ensuremath{y}}\fi}
\newunicodechar{ᵃ}{\ifmmode^{a}\else\textsuperscript{\ensuremath{a}}\fi}
\newunicodechar{ᵉ}{\ifmmode^{e}\else\textsuperscript{\ensuremath{e}}\fi}
\newunicodechar{ᵏ}{\ifmmode^{k}\else\textsuperscript{\ensuremath{k}}\fi}
\newunicodechar{ᵖ}{\ifmmode^{p}\else\textsuperscript{\ensuremath{p}}\fi}
\newunicodechar{ᴺ}{\ifmmode^{N}\else\textsuperscript{\ensuremath{N}}\fi}
\newunicodechar{ᶜ}{\ifmmode^{c}\else\textsuperscript{\ensuremath{c}}\fi}
\newunicodechar{ʰ}{\ifmmode^{h}\else\textsuperscript{\ensuremath{h}}\fi}
\newunicodechar{ᵠ}{\ifmmode^{q}\else\textsuperscript{\ensuremath{q}}\fi}
\newunicodechar{ₙ}{\ifmmode_{n}\else\textsubscript{\ensuremath{n}}\fi}
\newunicodechar{ₐ}{\ifmmode_{a}\else\textsubscript{\ensuremath{a}}\fi}
\newunicodechar{ᵢ}{\ifmmode_{i}\else\textsubscript{\ensuremath{i}}\fi}
\newunicodechar{ᵣ}{\ifmmode_{r}\else\textsubscript{\ensuremath{r}}\fi}
\newunicodechar{ₖ}{\ifmmode_{k}\else\textsubscript{\ensuremath{k}}\fi}
\newunicodechar{ᵦ}{\ifmmode_{\beta}\else\textsubscript{\ensuremath{\beta}}\fi}
\newunicodechar{ₓ}{\ifmmode_{x}\else\textsubscript{\ensuremath{x}}\fi}
\newfontfamily\popdisplay{ArchivoBlack-Regular.ttf}[Path=../../../fonts/]
\newfontfamily\poplogo{SpaceGrotesk-Bold.ttf}[Path=../../../fonts/]
\newfontfamily\popsemi{PlusJakartaSans-SemiBold.ttf}[Path=../../../fonts/]
\newfontfamily\popxbold{PlusJakartaSans-ExtraBold.ttf}[Path=../../../fonts/]
\newfontfamily\popxboldls{PlusJakartaSans-ExtraBold.ttf}[Path=../../../fonts/,LetterSpace=3.0]

\setlist[enumerate]{leftmargin=1.7em,itemsep=5pt,topsep=4pt}
\setlist[itemize]{leftmargin=1.7em,itemsep=4pt,topsep=4pt,label={\color{poporange}\textbullet}}
\setlength{\parindent}{0pt}
\setlength{\parskip}{5pt}
\renewcommand{\arraystretch}{1.25}

% pgfplots : style Pop par défaut
\pgfplotsset{
  every axis/.append style={axis line style={popink,line width=1pt},tick style={popink},
    grid style={popline,line width=0.5pt},label style={font=\small},tick label style={font=\footnotesize}},
  popcurve/.style={poporange,line width=1.8pt,no markers,smooth},
}
% Même style disponible dans un \draw TikZ simple (hors pgfplots)
\tikzset{popcurve/.style={poporange,line width=1.8pt}}
\tikzset{popgrid/.style={popline,very thin}}
\colorlet{popcurve}{poporange} % le nom du style est parfois utilisé comme couleur

% ============================================================
% Pill / squiggle
% ============================================================
\newcommand{\poppill}[3]{%
  \tikz[baseline=-0.55ex]\node[draw=popink,line width=1.2pt,rounded corners=2.6pt,%
    fill=#2,inner xsep=6.5pt,inner ysep=3pt]{%
    \popxboldls\fontsize{7}{7}\selectfont\color{#3}\MakeUppercase{#1}};%
}
\newcommand{\popsquiggle}[1]{%
  \tikz[baseline=-0.4ex]{
    \draw[#1,line width=2.6pt,line cap=round]
      plot [smooth] coordinates {(0,0.09) (0.34,0.01) (0.68,0.21) (1.02,0.01) (1.36,0.10)};
  }%
}
% Mot-clé surligné (marqueur orange sous le texte)
\newcommand{\cle}[1]{{\popxbold\color{popdark}#1}}

% ============================================================
% En-tête / pied
% ============================================================
\pagestyle{fancy}
\fancyhf{}
\renewcommand{\headrulewidth}{0pt}
\renewcommand{\footrulewidth}{0pt}
\lhead{\raisebox{-0.15ex}{\poplogo\fontsize{13}{13}\selectfont\color{popink}OnBuch\color{poporange}+}}
\rhead{\poppill{\DOCMATIERE\ \textbullet\ \DOCNIVEAU\ \textbullet\ Leçon \DOCLECON}{white}{popink}}
\lfoot{\popxbold\fontsize{7.6}{7.6}\selectfont\color{popmuted}\MakeUppercase{Cours signé OnBuch+}\ \textbullet\ {\popsemi\DOCTITRE}}
\rfoot{\tikz[baseline=-0.6ex]\node[rounded corners=8.5pt,fill=popink,minimum height=17pt,minimum width=17pt,inner xsep=5pt,inner sep=0]{\popxbold\fontsize{8}{8}\selectfont\color{popcream}\thepage\,/\,\pageref*{LastPage}};}

% ============================================================
% Titres
% ============================================================
\newcommand{\chaptitre}[2]{%
  \begin{tcolorbox}[enhanced,colback=poporange,colframe=popink,boxrule=2.6pt,arc=11pt,
      drop shadow=popink,left=15pt,right=15pt,top=12pt,bottom=11pt,before skip=3pt,after skip=18pt]
    \poppill{#2}{popink}{popcream}\par\vspace{9pt}
    {\raggedright\hyphenpenalty=10000\exhyphenpenalty=10000\popdisplay\fontsize{22}{24}\selectfont\color{popcream}\MakeUppercase{#1}\par}
  \end{tcolorbox}
}
% Section numérotée : pastille orange avec le numéro + titre Archivo
\newcounter{popsec}
\newcommand{\coursec}[1]{%
  \stepcounter{popsec}\setcounter{popsub}{0}%
  \par\vspace{16pt}\noindent
  \begin{minipage}{\linewidth}
  \tikz[baseline=-0.7ex]\node[draw=popink,line width=1.6pt,fill=poporange,rounded corners=4pt,minimum size=22pt,inner sep=0]{\popdisplay\fontsize{12}{12}\selectfont\color{popcream}\thepopsec};%
  \hspace{8pt}{\popdisplay\fontsize{15}{17}\selectfont\color{popink}\MakeUppercase{#1}}\par
  \vspace{4pt}\noindent{\color{poporange}\rule{36pt}{3.6pt}}
  \end{minipage}\par\nopagebreak\vspace{8pt}
}
\newcounter{popsub}
\newcommand{\courssub}[1]{%
  \stepcounter{popsub}%
  \par\vspace{9pt}\noindent{\popxbold\fontsize{11.5}{13}\selectfont\color{popink}{\color{poporange}\thepopsec.\thepopsub}\hspace{6pt}#1}\par\nopagebreak\vspace{4pt}
}

% ============================================================
% Boîtes pédagogiques — même recette : fond blanc, bordure épaisse colorée,
% ombre dure encre, badge plein. Titre optionnel [#1].
% ============================================================
\tcbset{popbase/.style={breakable,enhanced,colback=white,boxrule=2.3pt,arc=9pt,
  shadow={3pt}{-3pt}{0pt}{popink},left=14pt,right=14pt,top=11pt,bottom=11pt,before skip=11pt,after skip=15pt,
  fontupper=\popsemi\fontsize{10.6}{14.5}\selectfont\color{popink},
  fontlower=\popsemi\fontsize{10.6}{14.5}\selectfont\color{popink}}}
% Coche dessinée (le glyphe ✓ n'existe pas dans les polices du document)
\newcommand{\popcheck}[1]{\tikz[baseline=-0.5ex]\draw[#1,line width=1.6pt,line cap=round,line join=round](-0.13,0.0)--(-0.04,-0.09)--(0.13,0.1);}
\providecommand{\coche}{\popcheck{popgreen}}
\providecommand{\resultat}[1]{\par\smallskip\noindent\fcolorbox{popgreen}{popgreenL}{\parbox{\dimexpr\linewidth-2\fboxsep-2\fboxrule\relax}{\textbf{Résultat.} #1}}\par\smallskip}
\newcommand{\popboxhead}[3]{\poppill{#1}{#2}{white}\ifx\relax#3\relax\else\hspace{6pt}{\popxbold\fontsize{10.6}{12}\selectfont\color{popink}#3}\fi\par\vspace{7pt}}

\newtcolorbox{aretenir}[1][]{popbase,colframe=popgreen,before upper={\popboxhead{À retenir}{popgreen}{#1}}}
% Texte en anglais (document de lecture, dialogue, extrait) : cadre
% d'étude. Un titre optionnel (source, auteur) s'affiche dans l'en-tête.
\newtcolorbox{textebox}[1][]{popbase,colframe=popblue,before upper={\popboxhead{Text}{popblue}{#1}\begin{otherlanguage}{english}},after upper={\end{otherlanguage}}}
% Marques ✓ / ✗ (forme correcte / fautive) souvent utilisées par les modèles
\providecommand{\cross}{\textcolor{poppink}{\ensuremath{\times}}}
\providecommand{\xmark}{\cross}\providecommand{\crossmark}{\cross}\providecommand{\cmark}{\ensuremath{\checkmark}}
% Ligne de réponse / justification à compléter (inventée par les modèles)
\providecommand{\justification}[1]{\par\noindent\textit{Justification~:} \dotfill\par}
% \textipa (package tipa non chargé) : la police ne couvre pas l'API -> simple passage
\providecommand{\textipa}[1]{#1}
% Soulignement (ulem non chargé) : \uline, \ul
\providecommand{\uline}[1]{\underline{#1}}\providecommand{\ul}[1]{\underline{#1}}
% \glossaire{\item ...} inventée par les modèles : liste de glossaire
\providecommand{\glossaire}[1]{\begin{itemize}#1\end{itemize}}
% \alert (beamer) inventée par les modèles : texte mis en évidence
\providecommand{\alert}[1]{\textbf{\textcolor{poppink}{#1}}}
% variantes ASCII d'ordinaux écrites en macro
\providecommand{\iere}{\textsuperscript{re}}\providecommand{\ieme}{\textsuperscript{e}}\providecommand{\ere}{\textsuperscript{re}}\providecommand{\eme}{\textsuperscript{e}}\providecommand{\er}{\textsuperscript{er}}\providecommand{\ie}{\textsuperscript{e}}
% Ordinaux français écrits en macro par les modèles : 1\ière, 3\ième
\newcommand{\ière}{\textsuperscript{re}}\newcommand{\ième}{\textsuperscript{e}}
% \exemple (inventée par les modèles) : étiquette « Exemple : »
\providecommand{\exemple}{\textbf{Exemple~:}\ }
% \square utilisable aussi hors mode math (cases à cocher)
\AtBeginDocument{\let\squareMath\square\renewcommand{\square}{\ensuremath{\squareMath}}}
% Mot ou phrase en anglais à l'intérieur d'un paragraphe français
\newcommand{\eng}[1]{\ifmmode\text{\textit{#1}}\else\begingroup\selectlanguage{english}\itshape #1\endgroup\fi}
\let\esp\eng % alias toléré
% flèches d'intonation inventées par les modèles
\providecommand{\textsearrow}{}\renewcommand{\textsearrow}{\ensuremath{\searrow}}\providecommand{\textnearrow}{}\renewcommand{\textnearrow}{\ensuremath{\nearrow}}
% \fr{..} : mot français (inventé par les modèles)
\providecommand{\fr}[1]{\textit{#1}}
\newtcolorbox{exemplebox}[1][]{popbase,colframe=poporange,before upper={\popboxhead{Exemple}{poporange}{#1}}}
\newtcolorbox{definition}[1][]{popbase,colframe=popblue,before upper={\popboxhead{Définition}{popblue}{#1}}}
\newtcolorbox{propriete}[1][]{popbase,colframe=poppurple,before upper={\popboxhead{Propriété}{poppurple}{#1}}}
\newtcolorbox{methode}[1][]{popbase,colframe=popgold,before upper={\popboxhead{Méthode}{popgold}{#1}}}
\newtcolorbox{attention}[1][]{popbase,colframe=poppink,before upper={\popboxhead{Attention}{poppink}{#1}}}
\newtcolorbox{experience}[1][]{popbase,colframe=popblue,colback=popblueL,before upper={\popboxhead{Expérience}{popblue}{#1}}}
% « Le savais-tu ? » : carte violette pleine (contexte, culture, Cameroun)
\newtcolorbox{savaistu}[1][]{popbase,colback=poppurple,colframe=popink,
  fontupper=\popsemi\fontsize{10.4}{14.2}\selectfont\color{white},
  before upper={\poppill{Le savais-tu\,?}{white}{poppurple}\ifx\relax#1\relax\else\hspace{6pt}{\popxbold\color{white}#1}\fi\par\vspace{7pt}}}
% Exercice résolu : énoncé en haut, \tcblower puis la solution sur fond crème
\newcounter{popexo}\newcounter{popexr}
\newtcolorbox{exoresolu}[1][]{popbase,colframe=poporange,colbacklower=popcream,
  segmentation style={popink,line width=1.2pt,dashed},
  before upper={\stepcounter{popexr}\popboxhead{Exercice résolu \thepopexr}{poporange}{#1}},
  before lower={\poppill{Solution}{popink}{popcream}\par\vspace{6pt}}}
% Exercice (non résolu en place) — la correction est dans la partie Corrigés
\newtcolorbox{exercice}[1][]{popbase,colframe=popink,
  before upper={\stepcounter{popexo}\popboxhead{Exercice \thepopexo}{popink}{#1}}}
% Corrigé
\newtcolorbox{corrige}[1][]{popbase,colframe=popgreen,colback=popgreenL,
  before upper={\popboxhead{Corrigé}{popgreen}{#1}}}
% Objectifs / prérequis / bilan
\newtcolorbox{objectifs}{popbase,colframe=popink,colback=poporangeL,
  before upper={\popboxhead{Objectifs de la leçon}{popink}{}}}
\newtcolorbox{prerequis}{popbase,colframe=popink,colback=popblueL,
  before upper={\popboxhead{Prérequis}{popblue}{}}}
\newtcolorbox{bilan}{popbase,colframe=popink,colback=white,boxrule=2.8pt,
  before upper={\popboxhead{Fiche bilan}{poporange}{L'essentiel à connaître pour l'examen}}}
% Figure encadrée avec légende
\newtcolorbox{popfigure}[1][]{enhanced,breakable=false,colback=white,colframe=popline,boxrule=1.4pt,arc=8pt,
  left=10pt,right=10pt,top=10pt,bottom=8pt,before skip=10pt,after skip=12pt,halign=center,
  lower separated=false,halign lower=center,
  fontlower=\popsemi\fontsize{9}{11.5}\selectfont\color{popmuted},
  #1}
\newcommand{\legende}[1]{\par\vspace{6pt}{\popsemi\fontsize{9}{11.5}\selectfont\color{popmuted}#1}}

% ============================================================
% Signature OnBuch+
% ============================================================
\newcommand{\poplogoinline}[1]{{\poplogo\fontsize{#1}{#1}\selectfont\color{popink}OnBuch\color{poporange}+}}
\newcommand{\signatureonbuch}{%
  \begin{tcolorbox}[enhanced,colback=popink,colframe=popink,boxrule=2.2pt,arc=10pt,
      drop shadow=poporange,left=14pt,right=14pt,top=10pt,bottom=10pt,before skip=0pt,after skip=0pt]
    \tikz[baseline=-0.7ex]\node[circle,fill=popgreen,draw=popcream,line width=1.3pt,minimum size=22pt,inner sep=0]{\popcheck{white}};%
    \hspace{9pt}\begin{minipage}[c]{0.62\linewidth}
      {\popxbold\fontsize{12}{13}\selectfont\color{popcream}Signé\ \textbullet\ {\poplogo OnBuch\color{poporange}+}}\par\vspace{2pt}
      {\raggedright\popsemi\fontsize{8}{10.5}\selectfont\color{popline}\MakeUppercase{Équipe pédagogique OnBuch+}\\\MakeUppercase{Conforme au programme officiel MINESEC}\par}
    \end{minipage}\hfill
    {\poplogo\fontsize{22}{22}\selectfont\color{popcream}OnBuch\color{poporange}+}
  \end{tcolorbox}%
}

% ============================================================
% Page de garde
% #1 titre · #2 matière · #3 niveau · #4 module · #5 n° leçon · #6 durée ·
% #7 description / objectifs (texte ou itemize)
% ============================================================
\newcommand{\pagegarde}[7]{%
  \thispagestyle{empty}
  \newgeometry{a4paper,margin=1.7cm,top=1.6cm,bottom=1.4cm}%
  \begin{tikzpicture}[remember picture,overlay]
    \fill[poporange!18] ([xshift=-3.2cm,yshift=-3.6cm]current page.north east) circle (4.6cm);
    \fill[poppurple!14] ([xshift=2.4cm,yshift=7.6cm]current page.south west) circle (3.8cm);
  \end{tikzpicture}%
  {\poplogo\fontsize{30}{30}\selectfont\color{popink}OnBuch\color{poporange}+}\hfill
  \raisebox{0.6ex}{\poppill{Cours signé \textbullet\ Premium}{popink}{popcream}}\par
  \vspace{1pt}\popsquiggle{poppurple}\par
  \vspace{8pt}
  \begin{tcolorbox}[enhanced,colback=poporange,colframe=popink,boxrule=2.8pt,arc=13pt,
      drop shadow=popink,left=18pt,right=18pt,top=12pt,bottom=13pt,after skip=10pt]
    \poppill{#2 \textbullet\ #3}{popink}{popcream}\hspace{5pt}\poppill{#4}{white}{popink}\par\vspace{8pt}
    {\popdisplay\fontsize{14}{14}\selectfont\color{popink}LEÇON #5}\par\vspace{4pt}
    {\raggedright\hyphenpenalty=10000\exhyphenpenalty=10000\popdisplay\fontsize{25}{27}\selectfont\color{popcream}\MakeUppercase{#1}\par}
    \vspace{8pt}
    {\popxbold\fontsize{9.5}{9.5}\selectfont\color{popink}\MakeUppercase{Durée conseillée : #6}}
  \end{tcolorbox}
  \begin{tcolorbox}[enhanced,colback=white,colframe=popink,boxrule=2.2pt,arc=10pt,
      drop shadow=popink,left=16pt,right=16pt,top=10pt,bottom=10pt,after skip=8pt]
    \poppill{Dans cette leçon}{poppurple}{white}\par\vspace{5pt}
    \popsemi\fontsize{9.3}{12.6}\selectfont\color{popink}#7
  \end{tcolorbox}
  \noindent\begin{minipage}[t]{0.487\linewidth}
    \begin{tcolorbox}[enhanced,colback=popgreen,colframe=popink,boxrule=2.2pt,arc=10pt,
        drop shadow=popink,left=11pt,right=11pt,top=10pt,bottom=10pt,height=3.1cm,valign=center]
      \tcbox[colback=white,colframe=popink,boxrule=1.3pt,arc=5pt,boxsep=0pt,nobeforeafter,
        left=3pt,right=3pt,top=3pt,bottom=3pt,valign=center]{\qrcode[height=1.9cm]{https://play.google.com/store/apps/details?id=com.luvvix.onbuch&hl=fr}}%
      \hspace{8pt}\begin{minipage}[c]{0.5\linewidth}
        {\popdisplay\fontsize{11}{12.5}\selectfont\color{white}\MakeUppercase{Télécharge OnBuch}}\par\vspace{4pt}
        {\raggedright\popsemi\fontsize{8.4}{11}\selectfont\color{white}Gratuit \textbullet\ Google Play\\Tuteur IA, annales, concours, résultats\par}
      \end{minipage}
    \end{tcolorbox}
  \end{minipage}\hfill
  \begin{minipage}[t]{0.487\linewidth}
    \begin{tcolorbox}[enhanced,colback=poppurple,colframe=popink,boxrule=2.2pt,arc=10pt,
        drop shadow=popink,left=11pt,right=11pt,top=10pt,bottom=10pt,height=3.1cm,valign=center]
      \tikz[baseline=-0.7ex]\node[circle,fill=white,draw=popink,line width=1.3pt,minimum size=1.6cm,inner sep=0]{\popdisplay\fontsize{24}{24}\selectfont\color{poppurple}+};%
      \hspace{8pt}\begin{minipage}[c]{0.6\linewidth}
        {\popdisplay\fontsize{11}{12.5}\selectfont\color{white}\MakeUppercase{Passe à OnBuch+}}\par\vspace{4pt}
        {\raggedright\popsemi\fontsize{8.4}{11}\selectfont\color{white}Tous les cours signés, Tuteur IA illimité, fiches PDF — onglet OnBuch+ de l'app\par}
      \end{minipage}
    \end{tcolorbox}
  \end{minipage}
  \vspace{8pt}
  \popcontenu
  \vfill
  \signatureonbuch
  \restoregeometry
  \newpage
}

% Rangée « Ce cours contient » (page de garde)
\newcommand{\poptile}[3]{% #1 couleur #2 gros texte #3 légende
  \begin{tcolorbox}[enhanced,colback=#1,colframe=popink,boxrule=2pt,arc=9pt,shadow={2.5pt}{-2.5pt}{0pt}{popink},
      left=6pt,right=6pt,top=9pt,bottom=9pt,halign=center,valign=center,height=1.75cm,width=\linewidth,nobeforeafter]
    {\popdisplay\fontsize{12.5}{13}\selectfont\color{white}\MakeUppercase{#2}}\par\vspace{3pt}
    {\centering\popsemi\fontsize{7.8}{9.5}\selectfont\color{white}#3\par}
  \end{tcolorbox}}
\newcommand{\popcontenu}{%
  \par\noindent{\popxboldls\fontsize{7.5}{7.5}\selectfont\color{popmuted}CE COURS SIGNÉ CONTIENT}\par\vspace{7pt}
  \noindent\begin{minipage}[t]{0.235\linewidth}\poptile{popblue}{Cours}{complet, illustré, pas à pas}\end{minipage}\hfill
  \begin{minipage}[t]{0.235\linewidth}\poptile{poporange}{Méthodes}{et exercices résolus}\end{minipage}\hfill
  \begin{minipage}[t]{0.235\linewidth}\poptile{poppink}{Exercices}{type examen + corrigés}\end{minipage}\hfill
  \begin{minipage}[t]{0.235\linewidth}\poptile{popgreen}{Fiche bilan}{l'essentiel à retenir}\end{minipage}\par}

% Sommaire
\newtcolorbox{sommaire}{enhanced,colback=white,colframe=popink,boxrule=2.2pt,arc=10pt,
  drop shadow=popink,left=18pt,right=18pt,top=13pt,bottom=13pt,before skip=6pt,after skip=20pt,
  before upper={\poppill{Sommaire}{poppurple}{white}\par\vspace{9pt}\popsemi\fontsize{10.6}{14.5}\selectfont\color{popink}}}

% Signature de fin de cours — poussée tout en bas de la dernière page
\newcommand{\findecours}{%
  \par\vfill\nopagebreak
  \begin{center}\popsquiggle{poporange}\end{center}\vspace{4pt}
  \signatureonbuch
}
\AtBeginDocument{\def\llbracket{\mathopen{[\mkern-3.5mu[}}\def\rrbracket{\mathclose{]\mkern-3.5mu]}}} % intervalles d'entiers (glyphe ⟦ absent de la police)
% Glyphes absents de Fira Math : redéfinis à partir de glyphes disponibles
\AtBeginDocument{%
  \def\setminus{\mathbin{\backslash}}\let\smallsetminus\setminus
  \def\vdots{\mathord{\vbox{\baselineskip3.2pt\lineskiplimit0pt\kern4pt\hbox{.}\hbox{.}\hbox{.}}}}%
  \def\ddots{\mathinner{\mkern1mu\raise6.5pt\hbox{.}\mkern2mu\raise3.6pt\hbox{.}\mkern2mu\raise0.7pt\hbox{.}\mkern1mu}}%
  \def\triangle{\mathord{\scalebox{0.9}{$\symup{\Delta}$}}}%
  \def\nabla{\mathord{\raisebox{\depth}{\scalebox{1}[-1]{$\symup{\Delta}$}}}}%
  \def\checkmark{\popcheck{popdark}}%
  \def\star{\mathbin{\ast}}\def\bigstar{\mathord{\ast}}%
  \def\diamond{\mathbin{\scalebox{0.6}{$\Diamond$}}}%
  \def\bigcup{\mathop{\mathchoice{\raisebox{-0.3em}{\scalebox{1.7}{$\cup$}}}{\scalebox{1.25}{$\cup$}}{\cup}{\cup}}\displaylimits}%
  \def\bigcap{\mathop{\mathchoice{\raisebox{-0.3em}{\scalebox{1.7}{$\cap$}}}{\scalebox{1.25}{$\cap$}}{\cap}{\cap}}\displaylimits}%
  \def\lesseqqgtr{\mathrel{\lessgtr}}\def\bigoplus{\mathop{\oplus}\displaylimits}%
}
\providecommand{\ssi}{si et seulement si} % abréviation fréquente
\AtBeginDocument{\providecommand{\wideparen}{\overparen}} % arc de cercle

\begin{document}

\pagegarde{Lesson 45 — Vocabulary: Words and expressions related to polls/surveys}{Anglais}{Tle A}{Chapter 4 : Citizenship and Human Rights}{EN45}{2 h}{Cette leçon de vocabulaire présente les mots et expressions liés aux sondages et enquêtes d'opinion (polls/surveys). Les élèves apprennent le lexique thématique, les collocations, les faux amis et la famille des mots, puis réemploient ce lexique dans des phrases, dialogues et courts textes en lien avec la citoyenneté.
\vspace{4pt}
\begin{multicols}{2}\small
\begin{itemize}
\item À la fin de cette leçon, je sais nommer et définir les termes clés d'un sondage (poll, survey, sample, respondent, margin of error...)
\item je sais distinguer poll, survey, census et questionnaire
\item je sais utiliser les collocations propres au thème (conduct a survey, carry out a poll, a representative sample...)
\item je sais employer les verbes et expressions pour commenter des résultats (to lead, to rank, to account for, according to...)
\item je sais éviter les faux amis et calques du français (to assist ≠ assister à, actually ≠ actuellement, sensible ≠ sensible...)
\item je sais reconnaître et former les mots de la même famille (respond → respondent → response → responsive)
\end{itemize}\end{multicols}}

\begin{sommaire}
\begin{multicols}{2}
\begin{enumerate}
\item Mise en route : découvrir le monde des sondages
\item Champ lexical 1 : les acteurs et les outils du sondage
\item Champ lexical 2 : les actions et les résultats
\item Faux amis et pièges des francophones
\item Familles de mots, prononciation et orthographe
\item Emploi en contexte : lire, dire et écrire un sondage
\item Activité d'intégration
\item Méthodes et exercices résolus
\item Exercices
\item Corrigés des exercices
\item Fiche bilan
\end{enumerate}
\end{multicols}
\end{sommaire}

\begin{objectifs}
\begin{itemize}
\item À la fin de cette leçon, je sais nommer et définir les termes clés d'un sondage (poll, survey, sample, respondent, margin of error...)
\item je sais distinguer poll, survey, census et questionnaire
\item je sais utiliser les collocations propres au thème (conduct a survey, carry out a poll, a representative sample...)
\item je sais employer les verbes et expressions pour commenter des résultats (to lead, to rank, to account for, according to...)
\item je sais éviter les faux amis et calques du français (to assist ≠ assister à, actually ≠ actuellement, sensible ≠ sensible...)
\item je sais reconnaître et former les mots de la même famille (respond → respondent → response → responsive)
\item je sais prononcer correctement les mots difficiles du thème (survey, questionnaire, data, percentage)
\item je sais réemployer ce lexique à l'oral et à l'écrit dans des phrases, dialogues et courts textes sur la citoyenneté
\end{itemize}
\end{objectifs}

\begin{prerequis}
\begin{itemize}
\item Vocabulaire général de la citoyenneté et des droits de l'homme (citizen, vote, election, rights)
\item Les pourcentages, fractions et nombres en anglais (fifty per cent, two thirds, a quarter)
\item Le présent simple et le present perfect pour rapporter des faits et des résultats
\item Les structures de comparaison (more than, fewer than, the majority of)
\item La lecture de tableaux et de graphiques simples (bar chart, figure, rate)
\end{itemize}
\end{prerequis}

\newpage

\coursec{Mise en route : découvrir le monde des sondages}

\courssub{Brainstorming : qu'est-ce qu'un sondage ?}

Ouvrons cette leçon par une situation que tu as déjà vécue, à la télévision, à la radio ou sur ton téléphone :

\begin{textebox}[Situation d'accroche — In the news]
Imagine: a few months before elections in Cameroon, a research institute phones 1,200 people in Yaoundé, Douala and Bamenda and announces that 62\% of respondents trust the electoral process. In the USA or the UK, such opinion polls make headlines every week. But what exactly is a “sample”, a “margin of error”, or a “respondent”? Understanding this vocabulary is essential to read the news critically — and to shine in the Bac exam, where texts about surveys and citizenship are frequent.
\end{textebox}

« Imagine : quelques mois avant des élections au Cameroun, un institut de recherche téléphone à 1 200 personnes à Yaoundé, Douala et Bamenda et annonce que 62 \% des personnes interrogées font confiance au processus électoral. Aux États-Unis ou au Royaume-Uni, de tels sondages d'opinion font la une des journaux chaque semaine. Mais qu'est-ce exactement qu'un « échantillon », une « marge d'erreur » ou une « personne interrogée » ? Maîtriser ce vocabulaire est indispensable pour lire l'actualité avec un œil critique — et pour briller au baccalauréat, où les textes sur les sondages et la citoyenneté sont fréquents. »

\textbf{Problématique de la leçon :} comment lire, comprendre et raconter un sondage en anglais correct ? Pour cela, il faut d'abord posséder le vocabulaire précis du domaine : les acteurs, les outils, les actions et les résultats.

Avant toute explication, réponds mentalement à ces questions de brainstorming :

\begin{itemize}
\item Où as-tu déjà rencontré un sondage ? (élections, études de marché dans les supermarchés de Douala, questions sur les réseaux sociaux, enquêtes de satisfaction à l'école)
\item Qui pose les questions ? Qui y répond ?
\item Que fait-on des réponses ensuite ?
\end{itemize}

\begin{definition}[Qu'est-ce qu'un sondage ?]
Un \cle{sondage} est une enquête menée auprès d'une petite partie d'une population pour connaître les opinions ou les comportements de l'ensemble. En anglais, on dit principalement \eng{a poll} (sondage d'opinion, souvent politique) ou \eng{a survey} (enquête plus large). On interroge un \eng{sample} (échantillon) de personnes, appelées \eng{respondents} (personnes interrogées), et on publie les \eng{results} (résultats) avec une \eng{margin of error} (marge d'erreur).
\end{definition}

\courssub{Lecture guidée : Opinion Polls in Modern Democracies}

Lisons maintenant un court article de presse. Première lecture : ne t'arrête pas sur les mots inconnus, cherche seulement l'idée générale.

\begin{textebox}[Opinion Polls in Modern Democracies]
Last month, a polling institute conducted a survey among 1,500 respondents in three regions. According to the results, 58\% of those interviewed said they trusted the electoral commission, while 30\% disagreed and 12\% were undecided. The margin of error was plus or minus 3\%. The sample was chosen at random to be representative of the whole population.

Opinion polls play a key role in modern democracies. Before elections, they help journalists and citizens follow the popularity of candidates. After elections, they measure public opinion on new laws and government decisions. Market research companies also use surveys to find out what consumers want to buy.

However, polls are not perfect. If the sample is too small or badly chosen, the results can be misleading. That is why serious institutes always publish their methods: how many people were interviewed, when, and how the questions were worded. A well-conducted poll is a useful tool for democracy; a badly-conducted one can misinform millions of citizens.
\end{textebox}

\textbf{Glossaire :}

\begin{center}
\begin{tabularx}{\linewidth}{|l|l|X|}
\hline
\rowcolor{popblueL} \textbf{Mot anglais} & \textbf{Nature} & \textbf{Traduction} \\
\hline
to conduct & verbe & mener, réaliser (une enquête) \\
\hline
respondent & nom & personne interrogée, répondant \\
\hline
undecided & adjectif & indécis \\
\hline
margin of error & nom & marge d'erreur \\
\hline
at random & locution & au hasard \\
\hline
representative & adjectif & représentatif \\
\hline
misleading & adjectif & trompeur \\
\hline
to word (a question) & verbe & formuler (une question) \\
\hline
\end{tabularx}
\end{center}

\textbf{Questions de compréhension :}

\begin{enumerate}
\item How many respondents took part in the survey?
\item What percentage of people were undecided?
\item Why must a sample be chosen at random?
\item According to the text, what makes a poll “misleading”?
\end{enumerate}

\courssub{Deviner le sens d'un mot inconnu à partir du contexte}

Au baccalauréat, tu rencontreras toujours des mots inconnus dans les textes. La mauvaise méthode consiste à paniquer ou à traduire mot à mot. La bonne méthode consiste à faire des \cle{hypothèses de sens} à partir du contexte, puis à les vérifier.

Prenons un exemple dans notre texte : \eng{“The sample was chosen at random to be representative of the whole population.”} Supposons que tu ne connaisses pas \eng{sample}. Observons : c'est un nom (précédé de \eng{the}), il est « choisi », il doit être « représentatif de toute la population ». Hypothèse : il s'agit d'une partie de la population, d'un échantillon. Le contexte global (on interroge 1 500 personnes pour connaître l'opinion de tous) confirme l'hypothèse.

\begin{methode}[Deviner le sens d'un mot inconnu en 4 étapes]
\begin{enumerate}
\item \textbf{Identifier la nature du mot} : est-ce un nom, un verbe, un adjectif ? Regarde sa place dans la phrase et sa terminaison (\eng{-ed}, \eng{-ing}, \eng{-tion}…).
\item \textbf{Observer les mots voisins} : quels mots l'entourent ? Un article, une préposition, un adjectif qui le qualifie ?
\item \textbf{Chercher un mot de la même famille déjà connu} : \eng{respondent} ressemble à \eng{to respond} (répondre) ; \eng{representative} ressemble à \eng{to represent}.
\item \textbf{Vérifier avec le contexte global} : ton hypothèse a-t-elle un sens dans la phrase et dans le texte entier ? Seulement ensuite, vérifie au glossaire ou au dictionnaire.
\end{enumerate}
\end{methode}

\begin{exoresolu}[Deviner le sens de “misleading”]
Énoncé : dans la phrase \eng{“If the sample is too small or badly chosen, the results can be misleading.”}, devine le sens de l'adjectif \eng{misleading} sans dictionnaire, puis justifie ta réponse.
\tcblower
Solution détaillée :
\begin{enumerate}
\item \textbf{Nature} : \eng{misleading} suit le verbe \eng{can be} et qualifie \eng{the results} ; c'est un adjectif.
\item \textbf{Mots voisins} : la phrase parle d'un échantillon « trop petit ou mal choisi » — donc d'un problème.
\item \textbf{Famille de mots} : on reconnaît \eng{to lead} (mener, conduire) et le préfixe négatif \eng{mis-} (mal), comme dans \eng{misunderstand} (mal comprendre). Donc \eng{misleading} = « qui mène mal », qui induit en erreur.
\item \textbf{Contexte global} : le paragraphe explique que les sondages ne sont pas parfaits. Un résultat obtenu avec un mauvais échantillon peut donc être \textbf{trompeur}. Hypothèse confirmée par le glossaire : \eng{misleading} = trompeur.
\end{enumerate}
Traduction de la phrase : « Si l'échantillon est trop petit ou mal choisi, les résultats peuvent être trompeurs. »
\end{exoresolu}

\courssub{Première distinction : poll, survey, census, questionnaire, interview}

Ces cinq mots sont souvent confondus par les francophones, car le français utilise « sondage » et « enquête » de façon assez floue. L'anglais est plus précis. Observons d'abord des exemples :

\begin{exemplebox}[Cinq mots, cinq réalités]
\eng{The latest opinion poll gives the candidate 45\% of the votes.} « Le dernier sondage d'opinion donne 45 \% des voix au candidat. »

\eng{The ministry carried out a survey on youth unemployment.} « Le ministère a mené une enquête sur le chômage des jeunes. »

\eng{A national census counts every person living in the country.} « Un recensement national compte chaque personne vivant dans le pays. »

\eng{Please fill in this questionnaire about your eating habits.} « Veuillez remplir ce questionnaire sur vos habitudes alimentaires. »

\eng{The journalist had an interview with the mayor of Buea.} « Le journaliste a eu un entretien avec le maire de Buea. »
\end{exemplebox}

\begin{center}
\begin{tabularx}{\linewidth}{|l|X|X|}
\hline
\rowcolor{popblueL} \textbf{English} & \textbf{Français} & \textbf{Remarque} \\
\hline
a poll / an opinion poll & un sondage (d'opinion) & court, souvent politique ; quelques questions seulement \\
\hline
a survey & une enquête, une étude & plus large et plus détaillé qu'un poll \\
\hline
a census & un recensement & on interroge TOUTE la population, pas un échantillon \\
\hline
a questionnaire & un questionnaire & la feuille de questions écrites à remplir \\
\hline
an interview & une interview, un entretien & conversation directe, en face-à-face ou au téléphone \\
\hline
\end{tabularx}
\end{center}

\begin{propriete}[Règle d'emploi]
\begin{itemize}
\item On \textbf{mène} un sondage : \eng{to conduct / to carry out / to do a survey (or a poll)}. Jamais \eng{to make a survey} au sens de « mener ».
\item On \textbf{répond} à un questionnaire : \eng{to answer / to fill in a questionnaire}.
\item On \textbf{publie} les résultats : \eng{to publish / to release the results}.
\item On \textbf{interroge} les personnes : \eng{to interview people} ; les personnes interrogées sont \eng{the respondents} ou \eng{the people interviewed}.
\end{itemize}
\end{propriete}

\begin{attention}[Prononciation piège : SURvey ou surVEY ?]
Le mot \eng{survey} change d'accent selon sa nature :
\begin{itemize}
\item le \textbf{nom} s'accentue sur la première syllabe : \eng{a SURvey} (« SSEUR-vé ») ;
\item le \textbf{verbe} s'accentue sur la deuxième syllabe : \eng{to surVEY} (« seur-VÉ »).
\end{itemize}
C'est un phénomène fréquent en anglais : \eng{a REcord} / \eng{to reCORD}, \eng{an INcrease} / \eng{to inCREASE}. Retiens : nom = accent devant, verbe = accent derrière.
\end{attention}

\begin{savaistu}[Pourquoi “poll” ?]
Le mot \eng{poll} vient d'un vieux mot anglais signifiant « tête ». Au Moyen Âge, pour connaître l'opinion d'une assemblée, on comptait littéralement les têtes : compter les têtes, c'était compter les votes ! C'est pourquoi on dit encore aujourd'hui \eng{to go to the polls} pour « aller voter », et \eng{polling station} pour « bureau de vote ». Le mot \eng{poll tax}, dans l'histoire britannique, désignait un impôt « par tête », c'est-à-dire par personne.
\end{savaistu}

\courssub{Carte mentale : le monde des sondages en un coup d'œil}

Voici la carte mentale que nous allons remplir tout au long de cette leçon. Pour l'instant, retiens les cinq grandes familles du vocabulaire des sondages.

\begin{popfigure}
\begin{center}
\begin{tikzpicture}[mindmap, grow cyclic,
  every node/.style={concept, align=center, font=\small},
  root concept/.append style={concept color=popblue, font=\bfseries},
  level 1/.append style={level distance=3.2cm, sibling angle=72},
  level 1 concept/.append style={concept color=poporange}]
\node[root concept]{POLLS \& SURVEYS}
  child { node[align=center]{people\\ \emph{respondent,}\\\emph{interviewer,}\\\emph{pollster}} }
  child { node[align=center]{tools\\ \emph{questionnaire,}\\\emph{sample,}\\\emph{interview}} }
  child { node[align=center]{actions\\ \emph{conduct,}\\\emph{interview,}\\\emph{publish}} }
  child { node[align=center]{results\\ \emph{margin of error,}\\\emph{percentage,}\\\emph{findings}} }
  child { node[align=center]{opinions\\ \emph{trust,}\\\emph{disagree,}\\\emph{undecided}} };
\end{tikzpicture}
\end{center}
\legende{Carte mentale du champ lexical des sondages : cinq familles de mots à explorer.}
\end{popfigure}

\begin{exercice}[À toi de jouer]
\begin{enumerate}
\item Classe ces mots dans les cinq branches de la carte mentale : \eng{respondent, to carry out, questionnaire, undecided, margin of error, pollster, to publish, findings, sample, to trust}.
\item Complète avec \eng{poll}, \eng{survey}, \eng{census}, \eng{questionnaire} ou \eng{interview} :
\begin{enumerate}
\item The government organises a national \ldots{} every ten years to count the population.
\item Students had to fill in a two-page \ldots{} about their study habits.
\item According to a recent opinion \ldots{}, most young Cameroonians use social media every day.
\item The radio journalist asked for an \ldots{} with the new minister.
\end{enumerate}
\end{enumerate}
\end{exercice}

\begin{corrige}[Exercice : À toi de jouer]
\textbf{1.} People : \eng{respondent, pollster}. Tools : \eng{questionnaire, sample}. Actions : \eng{to carry out, to publish}. Results : \eng{margin of error, findings}. Opinions : \eng{undecided, to trust}.

\textbf{2.} a) \eng{census} (on compte toute la population) ; b) \eng{questionnaire} (une feuille de questions à remplir) ; c) \eng{poll} (sondage d'opinion rapide) ; d) \eng{interview} (entretien direct avec une personnalité).
\end{corrige}

\begin{aretenir}[L'essentiel de cette mise en route]
\begin{itemize}
\item Un sondage interroge un \eng{sample} (échantillon) représentatif, choisi \eng{at random} (au hasard), et publie ses résultats avec une \eng{margin of error}.
\item \eng{Poll} = sondage court ; \eng{survey} = enquête détaillée ; \eng{census} = recensement de toute la population ; \eng{questionnaire} = document de questions ; \eng{interview} = entretien direct.
\item Pour deviner un mot inconnu : nature, mots voisins, famille de mots, contexte global — le dictionnaire vient en dernier.
\item Attention à la prononciation : \eng{a SURvey} mais \eng{to surVEY}.
\end{itemize}
\end{aretenir}

\coursec{Champ lexical 1 : les acteurs et les outils du sondage}

Dans la section précédente, nous avons découvert le monde des sondages : pourquoi on sonde, à quoi servent les enquêtes d'opinion dans une démocratie, et comment les résultats circulent dans les médias. Mais qui fabrique un sondage ? Avec quels outils ? Avant d'interroger, de calculer et de publier, il faut des \cle{personnes}, des \cle{instruments} et des \cle{institutions}. C'est ce premier champ lexical que nous allons construire pas à pas, en partant d'un texte authentique.

\courssub{Observation : un texte pour découvrir le vocabulaire}

\begin{textebox}[How a poll is made — adapted from a press article]
Last month, the National Institute of Statistics in Yaoundé asked a polling institute to measure public opinion on the new transport plan. The institute first defined the population to be studied: all adults living in the city. As it was impossible to question everyone, the researchers selected a sample of 1,500 people, chosen to represent the whole population in terms of age, sex and neighbourhood.

Then the institute trained a team of interviewers. Each interviewer received a questionnaire containing twenty questions. Some were multiple-choice questions, where respondents simply ticked a box; others were open questions, where people could answer freely in their own words. One question used a scale from 1 to 5: “How satisfied are you with the new bus lines?”

The interviewers worked in markets, schools and motor parks. Most respondents were cooperative, but about a quarter refused to answer. The completed forms were sent back to the research centre, where statisticians entered the data into computers. Three weeks later, the polling institute published its results: 62 \% of the respondents supported the transport plan.

A journalist from a research centre in Douala commented: “A poll is only as good as its sample and its questionnaire. If the questions are badly written, the answers mean nothing.”
\end{textebox}

\textbf{Glossaire}

\begin{center}
\begin{tabularx}{\linewidth}{|l|l|X|}
\hline
\rowcolor{popblueL} \textbf{Mot anglais} & \textbf{Nature} & \textbf{Traduction française} \\
\hline
polling institute & nom & institut de sondage \\
\hline
sample & nom & échantillon \\
\hline
interviewer & nom & enquêteur, enquêtrice (celui qui pose les questions) \\
\hline
questionnaire & nom & questionnaire \\
\hline
multiple-choice question & nom & question à choix multiple (QCM) \\
\hline
open question & nom & question ouverte \\
\hline
respondent & nom & personne interrogée, répondant \\
\hline
scale from 1 to 5 & nom & échelle de 1 à 5 \\
\hline
form & nom & formulaire \\
\hline
research centre & nom & centre de recherche \\
\hline
data & nom (pluriel) & données \\
\hline
statistics office & nom & office / institut de statistiques \\
\hline
\end{tabularx}
\end{center}

\textbf{Questions de compréhension}
\begin{enumerate}
\item Who selected the sample of 1,500 people?
\item What is the difference between a multiple-choice question and an open question, according to the text?
\item Where did the interviewers work?
\item What percentage of the respondents supported the transport plan?
\end{enumerate}

\courssub{Les personnes : qui participe à un sondage ?}

Observons d'abord les noms de personnes du texte : \eng{polling institute}, \eng{interviewer}, \eng{respondent}, \eng{sample}, \eng{population}. Chacun désigne un rôle précis dans la chaîne du sondage.

\begin{definition}[Les acteurs du sondage]
\begin{itemize}
\item \cle{pollster} : le sondeur, la personne (ou l'organisation) qui conçoit et réalise le sondage. \eng{The pollster asked ten questions.} « Le sondeur a posé dix questions. »
\item \cle{interviewer} : l'enquêteur, la personne qui pose les questions, souvent face à face ou au téléphone. \eng{The interviewer visited fifty households in Bamenda.} « L'enquêteur a visité cinquante ménages à Bamenda. »
\item \cle{respondent} (ou \cle{interviewee}) : la personne qui répond aux questions. \eng{Each respondent answered for about ten minutes.} « Chaque personne interrogée a répondu pendant environ dix minutes. »
\item \cle{sample} : l'échantillon, le petit groupe choisi pour représenter l'ensemble. \eng{A sample of 2,000 voters was selected.} « Un échantillon de 2 000 électeurs a été sélectionné. »
\item \cle{population} : la population étudiée, c'est-à-dire l'ensemble des personnes concernées par l'enquête (et non pas forcément tous les habitants d'un pays). \eng{The population studied was young people aged 18 to 25.} « La population étudiée était les jeunes de 18 à 25 ans. »
\item \cle{panel} : un panel, un groupe de personnes interrogées régulièrement sur une longue période. \eng{The panel meets every six months to discuss consumer habits.} « Le panel se réunit tous les six mois pour discuter des habitudes de consommation. »
\end{itemize}
\end{definition}

\begin{propriete}[Former les noms de personnes avec -er et -ee]
L'anglais forme souvent des paires de noms d'acteurs avec les suffixes \cle{-er} (celui qui fait l'action) et \cle{-ee} (celui qui la subit) :
\begin{itemize}
\item \eng{interviewer} (celui qui interroge) / \eng{interviewee} (celui qui est interrogé) ;
\item \eng{employer} (employeur) / \eng{employee} (employé) ;
\item \eng{trainer} (formateur) / \eng{trainee} (stagiaire).
\end{itemize}
Attention à la prononciation : \eng{interviewee} se prononce « in-tèr-vyou-ï », avec l'accent sur la dernière syllabe, comme en français « interviewé ».
\end{propriete}

\begin{attention}[Respondent ou correspondent ?]
Un \eng{respondent} est la personne qui \textbf{répond} à un questionnaire. Ne pas confondre avec \eng{correspondent}, qui désigne un journaliste envoyé à l'étranger : \eng{our correspondent in London} = « notre correspondant à Londres ». De même, \eng{respondent} n'a rien à voir avec un « répondant » au sens de « garant ».
\end{attention}

\courssub{Les outils : avec quoi sonde-t-on ?}

\begin{definition}[Les instruments du sondage]
\begin{itemize}
\item \cle{questionnaire} : le questionnaire, l'ensemble écrit des questions. \eng{The questionnaire contains twenty questions.} « Le questionnaire contient vingt questions. »
\item \cle{question} : la question. Attention à la prononciation : « kouèss-tchène ».
\item \cle{multiple-choice question} (souvent abrégé \eng{MCQ}) : question à choix multiple, où l'on coche une réponse. \eng{In a multiple-choice question, you tick only one box.} « Dans une question à choix multiple, on ne coche qu'une seule case. »
\item \cle{open question} : question ouverte, où le répondant répond librement. \eng{Open questions give richer answers but take longer to analyse.} « Les questions ouvertes donnent des réponses plus riches mais prennent plus de temps à analyser. »
\item \cle{form} : le formulaire, le document à remplir. \eng{Please fill in the form and return it before Friday.} « Veuillez remplir le formulaire et le renvoyer avant vendredi. »
\item \cle{scale (from 1 to 5)} : l'échelle (de 1 à 5), utilisée pour mesurer un degré d'accord ou de satisfaction. \eng{On a scale from 1 to 5, how do you rate the service?} « Sur une échelle de 1 à 5, comment évaluez-vous le service ? »
\end{itemize}
\end{definition}

\begin{exemplebox}[Collocations utiles avec les outils du sondage]
\begin{itemize}
\item \eng{to fill in a form} = remplir un formulaire (jamais \eng{to fill a form} seul : il faut la particule \eng{in}).
\item \eng{to answer a question} = répondre à une question (pas de préposition en anglais, contrairement au français « répondre à »).
\item \eng{to tick a box} = cocher une case.
\item \eng{to rate something on a scale} = évaluer quelque chose sur une échelle.
\item \eng{to design a questionnaire} = concevoir un questionnaire.
\end{itemize}
\end{exemplebox}

\courssub{Les institutions : qui organise les sondages ?}

\begin{definition}[Les institutions de l'enquête]
\begin{itemize}
\item \cle{polling institute} : institut de sondage (organisme privé ou public qui réalise les enquêtes). \eng{A well-known polling institute published the results.} « Un institut de sondage réputé a publié les résultats. »
\item \cle{research centre} : centre de recherche (souvent universitaire). \eng{The research centre studies voting behaviour.} « Le centre de recherche étudie le comportement électoral. »
\item \cle{statistics office} : office de statistiques (organisme officiel de l'État). \eng{The National Statistics Office counts the population every ten years.} « L'office national de statistiques recense la population tous les dix ans. »
\end{itemize}
\end{definition}

\begin{savaistu}[Gallup, YouGov et les panels]
Aux États-Unis et au Royaume-Uni, des instituts célèbres comme \eng{Gallup} ou \eng{YouGov} interrogent des \eng{panels}, c'est-à-dire des groupes de personnes suivies régulièrement, parfois pendant des années. Interroger toujours les mêmes personnes permet de mesurer l'\emph{évolution} des opinions, et pas seulement une opinion à un moment donné. Au Cameroun, l'institut national de la statistique réalise de grandes enquêtes auprès des ménages, par exemple sur l'emploi ou la pauvreté.
\end{savaistu}

\courssub{La chaîne du sondage : qui fait quoi ?}

La figure suivante résume le parcours d'un sondage, de l'institut jusqu'aux données.

\begin{popfigure}
\begin{center}
\begin{tikzpicture}[
  node distance=0.55cm and 0.9cm,
  box/.style={draw=popblue, fill=popblueL, rounded corners=2pt, align=center, minimum height=0.9cm, inner sep=4pt, font=\small},
  arr/.style={-{Stealth[length=2.5mm]}, thick, poporange}
]
\node[box, align=center] (inst){polling institute\\ institut de sondage};
\node[box, right=of inst, align=center] (inter){interviewer\\ enquêteur};
\node[box, right=of inter, align=center] (quest){questionnaire\\ questionnaire};
\node[box, right=of quest, align=center] (resp){respondent\\ répondant};
\node[box, right=of resp, align=center] (data){data\\ données};
\draw[arr] (inst) -- (inter);
\draw[arr] (inter) -- (quest);
\draw[arr] (quest) -- (resp);
\draw[arr] (resp) -- (data);
\end{tikzpicture}
\end{center}
\legende{La chaîne du sondage : l'institut forme l'enquêteur, qui administre le questionnaire au répondant ; les réponses deviennent des données.}
\end{popfigure}

\courssub{Tableau récapitulatif du vocabulaire}

\begin{center}
\begin{tabularx}{\linewidth}{|l|X|X|}
\hline
\rowcolor{popblueL} \textbf{English} & \textbf{Français} & \textbf{Exemple en contexte} \\
\hline
pollster & sondeur & \eng{The pollster asked ten questions.} \\
\hline
interviewer & enquêteur & \eng{The interviewer knocked at our door.} \\
\hline
respondent / interviewee & personne interrogée & \eng{Only 40 \% of the respondents returned the form.} \\
\hline
sample & échantillon & \eng{A sample of 2,000 voters was selected.} \\
\hline
population & population étudiée & \eng{The population is all adults in Douala.} \\
\hline
panel & panel & \eng{The panel is interviewed every year.} \\
\hline
questionnaire & questionnaire & \eng{The questionnaire has twenty items.} \\
\hline
multiple-choice question & question à choix multiple & \eng{Tick one box for each multiple-choice question.} \\
\hline
open question & question ouverte & \eng{The last open question asks for suggestions.} \\
\hline
form & formulaire & \eng{Fill in the form carefully.} \\
\hline
scale from 1 to 5 & échelle de 1 à 5 & \eng{Rate the service on a scale from 1 to 5.} \\
\hline
polling institute & institut de sondage & \eng{The polling institute published the results.} \\
\hline
research centre & centre de recherche & \eng{A research centre analysed the data.} \\
\hline
statistics office & office de statistiques & \eng{The statistics office released the figures.} \\
\hline
\end{tabularx}
\end{center}

\courssub{Exercice résolu : associer le mot à sa définition}

\begin{exoresolu}[Match the word with its definition]
Reliez chaque mot (1--6) à sa définition (a--f).

1. sample \quad 2. interviewer \quad 3. questionnaire \quad 4. panel \quad 5. scale \quad 6. polling institute

\begin{itemize}
\item[a.] a set of written questions used to collect information
\item[b.] a small group of people chosen to represent a larger population
\item[c.] an organisation that carries out surveys
\item[d.] the person who asks the questions during a survey
\item[e.] a group of people who are interviewed regularly over time
\item[f.] a range of numbers used to measure opinion, for example from 1 to 5
\end{itemize}
\tcblower
\textbf{Solution détaillée}

1 -- b : un \eng{sample} est un échantillon, un petit groupe représentatif d'une population plus large.

2 -- d : l'\eng{interviewer} est celui qui pose les questions (suffixe \eng{-er} = celui qui fait l'action).

3 -- a : le \eng{questionnaire} est l'ensemble des questions écrites.

4 -- e : le \eng{panel} est un groupe suivi régulièrement dans le temps.

5 -- f : la \eng{scale} est l'échelle de mesure, par exemple de 1 à 5.

6 -- c : le \eng{polling institute} est l'organisme qui réalise les sondages.
\end{exoresolu}

\courssub{Pièges fréquents des francophones}

\begin{attention}[Erreurs à éviter]
\begin{itemize}
\item Ne dites pas \eng{an inquiry} pour « une enquête d'opinion » : utilisez \eng{a survey} ou \eng{a poll}. \eng{Inquiry} désigne plutôt une enquête officielle ou judiciaire.
\item \eng{Actually} ne signifie pas « actuellement » mais « en fait » : \eng{The pollster actually works for a research centre.} « Le sondeur travaille en fait pour un centre de recherche. »
\item On dit \eng{to answer a question} sans préposition : jamais \eng{to answer to a question}.
\item \eng{Data} est grammaticalement le pluriel de \eng{datum} ; dans l'usage courant on dit souvent \eng{the data is}, mais à l'écrit soigné préférez \eng{the data are analysed}.
\item Ne traduisez pas « un formulaire » par \eng{a formula} : \eng{formula} signifie « formule » (mathématiques, chimie). Le bon mot est \eng{form}.
\end{itemize}
\end{attention}

\begin{exercice}[À vous de jouer]
Complétez les phrases avec le mot qui convient : \eng{sample}, \eng{interviewer}, \eng{questionnaire}, \eng{respondents}, \eng{scale}, \eng{polling institute}.

\begin{enumerate}
\item The \dots\dots\dots asked each person the same ten questions.
\item A \dots\dots\dots of 800 students was chosen from three schools in Buea.
\item Please rate your satisfaction on a \dots\dots\dots from 1 to 5.
\item The \dots\dots\dots published its results in the newspaper.
\item Half of the \dots\dots\dots did not answer the last open question.
\item The \dots\dots\dots contains both multiple-choice and open questions.
\end{enumerate}
\end{exercice}

\begin{corrige}[Exercice « À vous de jouer »]
1. interviewer — 2. sample — 3. scale — 4. polling institute — 5. respondents — 6. questionnaire.
\end{corrige}

Nous savons désormais qui intervient dans un sondage et avec quels instruments. Dans la section suivante, nous passerons au second champ lexical : les \emph{actions} du sondage (to conduct, to carry out, to survey) et le vocabulaire des \emph{résultats} (percentage, majority, trend), qui nous permettra de raconter une enquête du début à la fin.

\coursec{Champ lexical 2 : les actions et les résultats}

Dans la section précédente, nous avons rencontré les \emph{acteurs} du sondage (\eng{pollster}, \eng{respondent}, \eng{interviewer}, \eng{sample}) et ses \emph{outils} (\eng{questionnaire}, \eng{opinion poll}, \eng{survey}). Mais un sondage, ce n'est pas seulement des gens et des fiches : c'est d'abord une \emph{chaîne d'actions} — on mène l'enquête, on interroge, on recueille les données, on les analyse — puis des \emph{résultats} que l'on publie et que l'on commente. C'est précisément ce vocabulaire-là, celui du verbe et du chiffre, que l'examinateur du Baccalauréat attend dans vos essais sur la citoyenneté. Observons-le d'abord en situation réelle.

\courssub{Observer : un article de presse sur un sondage camerounais}

\begin{textebox}[Survey reveals growing interest in voter registration — The Cameroon Herald]
A survey carried out last month by the Yaoundé-based institute Civic Voice has revealed a growing interest in voter registration among young Cameroonians. The pollsters interviewed 2,400 respondents in Douala, Yaoundé, Buea and Bamenda, and collected data on their political attitudes.

According to the findings, 68\% of respondents said that they intended to register on the electoral roll before the next election, a figure significantly higher than in the previous survey, conducted two years ago. Women account for 52\% of the sample, and they appear slightly more motivated than men: 71\% of female respondents plan to register, against 65\% of male respondents.

The results also show a clear regional trend. In the South-West, registration intentions rank first, with roughly three quarters of the people interviewed declaring themselves ready to vote. On average, more than half of the respondents under 25 said they trusted the electoral process, although 18\% remained undecided.

“We analysed the data carefully before releasing the figures,” explained the director of Civic Voice. “The data show that civic education campaigns are beginning to bear fruit.” The institute plans to conduct another survey next year to confirm this trend.
\end{textebox}

\textbf{Glossaire}

\begin{center}
\begin{tabularx}{\linewidth}{|l|l|X|}
\hline
\rowcolor{popblueL} English & Nature & Français \\
\hline
to carry out a survey & verbe & mener / réaliser une enquête \\
\hline
findings & nom pluriel & résultats, conclusions (d'une enquête) \\
\hline
electoral roll & nom & liste électorale \\
\hline
to account for & verbe & représenter (une proportion) \\
\hline
trend & nom & tendance \\
\hline
to rank first & verbe & arriver en tête, se classer premier \\
\hline
roughly & adverbe & environ, à peu près \\
\hline
undecided & adjectif & indécis \\
\hline
to release the figures & verbe & publier les chiffres \\
\hline
to bear fruit & expression & porter ses fruits \\
\hline
\end{tabularx}
\end{center}

\textbf{Questions de compréhension}
\begin{enumerate}
\item How many people were interviewed, and in which towns?
\item What percentage of respondents intend to register to vote?
\item Which region ranks first for registration intentions?
\item What does the director say about the data?
\end{enumerate}

\begin{corrige}[Questions de compréhension]
1. \eng{The pollsters interviewed 2,400 respondents in Douala, Yaoundé, Buea and Bamenda.} 2. \eng{68\% of respondents intend to register to vote.} 3. \eng{The South-West ranks first for registration intentions.} 4. \eng{He says that the data were analysed carefully and that they show civic education campaigns are beginning to bear fruit.}
\end{corrige}

\courssub{Les actions : les verbes du sondage}

Lisez à nouveau le texte et soulignez les verbes : \eng{carried out}, \eng{interviewed}, \eng{collected data}, \eng{analysed}, \eng{releasing}, \eng{to conduct}. Vous remarquez qu'aucun de ces verbes n'est \eng{to make} ! C'est la grande leçon de cette section : chaque nom du champ du sondage s'associe à un verbe précis. Ce couple verbe + nom s'appelle une \cle{collocation}.

\begin{propriete}[Les collocations essentielles du sondage]
En anglais, on ne « fait » pas une enquête : on la \emph{mène}. Retenez ces couples verbe + nom, incontournables au Bac :
\begin{itemize}
\item \eng{to conduct a survey} / \eng{to carry out a survey} = mener, réaliser une enquête (registre soutenu, le plus sûr à l'écrit) ;
\item \eng{to run a poll} = organiser / réaliser un sondage (registre journalistique) ;
\item \eng{to interview respondents} / \eng{to ask questions} = interroger les personnes sondées ;
\item \eng{to answer (a questionnaire)} = répondre (à un questionnaire) — sans préposition : on dit \eng{answer the questions}, jamais \emph{answer to the questions} ;
\item \eng{to collect data} = recueillir les données ;
\item \eng{to analyse the results} = analyser les résultats ;
\item \eng{to publish / to release the findings} = publier les résultats ;
\item \eng{to take part in a survey} = participer à une enquête (côté répondant).
\end{itemize}
\end{propriete}

\begin{attention}[Le calque « to make a survey »]
Le français « faire une enquête » pousse les candidats francophones à écrire \emph{to make a survey}. Cette forme, parfois tolérée dans l'usage relâché, est à \textbf{éviter au Baccalauréat} : préférez toujours \eng{to conduct} ou \eng{to carry out a survey}. De même, on ne « pose » pas une question avec \emph{to pose} : on dit \eng{to ask a question}. Enfin, « répondre à » se traduit par \eng{to answer} suivi directement du complément : \eng{She answered my letter.} et non \emph{She answered to my letter.}
\end{attention}

\begin{exemplebox}[Les verbes en contexte camerounais]
\eng{The institute conducted a survey on citizens' trust in the media.} « L'institut a mené une enquête sur la confiance des citoyens dans les médias. »

\eng{A private firm ran a poll in Douala before the municipal elections.} « Un cabinet privé a réalisé un sondage à Douala avant les élections municipales. »

\eng{Interviewers asked the respondents ten questions about public services.} « Les enquêteurs ont posé dix questions aux personnes sondées sur les services publics. »

\eng{Only 60\% of the people contacted agreed to take part in the survey.} « Seulement 60 \% des personnes contactées ont accepté de participer à l'enquête. »

\eng{The team collected the data in three weeks and analysed the results in March.} « L'équipe a recueilli les données en trois semaines et a analysé les résultats en mars. »
\end{exemplebox}

\begin{popfigure}
\begin{center}
\begin{tikzpicture}[
etape/.style={rectangle, rounded corners=4pt, draw=popblue, fill=popblueL, align=center, minimum width=2.6cm, minimum height=1.1cm, font=\small},
fleche/.style={-{Stealth[length=2.5mm]}, thick, poporange}]
\node[etape, align=center] (e1) at (0,0){choose\\a sample};
\node[etape, right=0.7cm of e1, align=center] (e2){ask\\questions};
\node[etape, right=0.7cm of e2, align=center] (e3){collect\\data};
\node[etape, right=0.7cm of e3, align=center] (e4){analyse\\the results};
\node[etape, right=0.7cm of e4, align=center] (e5){publish\\the findings};
\draw[fleche] (e1) -- (e2);
\draw[fleche] (e2) -- (e3);
\draw[fleche] (e3) -- (e4);
\draw[fleche] (e4) -- (e5);
\end{tikzpicture}
\end{center}
\legende{Les cinq étapes d'un sondage : choisir un échantillon, poser les questions, recueillir les données, analyser les résultats, publier les conclusions.}
\end{popfigure}

\courssub{Les résultats : les noms du chiffre}

Une fois les données analysées, il faut les \emph{nommer}. Voici les noms qui reviennent dans tous les articles de presse anglophones :

\begin{center}
\begin{tabularx}{\linewidth}{|l|X|X|}
\hline
\rowcolor{popblueL} English & Français & Exemple \\
\hline
findings (pluriel) & résultats, conclusions & \eng{The findings were published on Monday.} \\
\hline
results & résultats & \eng{The results confirm earlier polls.} \\
\hline
figures & chiffres & \eng{These figures are encouraging.} \\
\hline
percentage & pourcentage & \eng{A high percentage of voters agree.} \\
\hline
rate & taux & \eng{The response rate was 80\%.} \\
\hline
majority & majorité & \eng{A majority of respondents said yes.} \\
\hline
minority & minorité & \eng{Only a minority opposed the reform.} \\
\hline
trend & tendance & \eng{The trend is clearly upward.} \\
\hline
\end{tabularx}
\end{center}

\begin{attention}[« Data » : singulier ou pluriel ?]
\eng{Data} est à l'origine le pluriel du latin \emph{datum}. Dans l'anglais soutenu, on écrit donc \eng{the data \textbf{show} that...} (pluriel). Dans l'usage courant, \eng{the data \textbf{is}...} est de plus en plus accepté. Au Bac, jouez la sécurité : \eng{the data show that...}. Autre piège : \eng{information} est, lui, toujours \textbf{indénombrable} — jamais \emph{informations} avec un « s » ! On dira \eng{a piece of information} pour « une information ».
\end{attention}

\courssub{Commenter les résultats : les phrases du journaliste}

Savoir nommer un chiffre ne suffit pas : il faut savoir le \emph{présenter} et le \emph{commenter}. Les articles de presse utilisent des formules quasi fixes :

\begin{itemize}
\item \eng{According to the poll, ...} = Selon le sondage, ... (citer la source) ;
\item \eng{58\% of respondents said that...} = 58 \% des personnes interrogées ont déclaré que... ;
\item \eng{The incumbent leads by 12 points.} = Le sortant devance de 12 points ;
\item \eng{The South-West ranks first.} = Le Sud-Ouest arrive en tête ;
\item \eng{Women account for 52\% of the sample.} = Les femmes représentent 52 \% de l'échantillon ;
\item \eng{18\% of voters are still undecided.} = 18 \% des électeurs sont encore indécis.
\end{itemize}

\begin{exemplebox}[Deux modèles à imiter]
\eng{According to the poll, the incumbent leads by 12 points.} « Selon le sondage, le sortant devance de 12 points. »

\eng{Women account for 52\% of the sample.} « Les femmes représentent 52 \% de l'échantillon. »
\end{exemplebox}

Pour \emph{nuancer} — qualité très valorisée à l'écrit —, ajoutez les adverbes de mesure : \eng{slightly} (légèrement), \eng{significantly} (nettement, de façon significative), \eng{on average} (en moyenne), \eng{roughly} / \eng{about} (environ), \eng{more than half} (plus de la moitié). Exemple : \eng{Turnout rose slightly, by roughly two points, while more than half of young voters stayed at home.} « La participation a légèrement augmenté, d'environ deux points, tandis que plus de la moitié des jeunes électeurs sont restés chez eux. »

\begin{attention}[Prononciation et orthographe : trois pièges]
\begin{itemize}
\item \eng{survey} : le nom se prononce « seu-veï » (accent sur la première syllabe) ; le verbe \eng{to survey} se prononce « se-veï » (accent sur la deuxième syllabe). Même mot, deux accents !
\item \eng{data} se prononce « deï-te » ; \eng{percentage} « pe-sèn-tidj » ; \eng{majority} « me-djo-ri-ti ».
\item Orthographe : \eng{analyse} (verbe, orthographe britannique) mais \eng{analysis} (nom) ; au pluriel, \eng{analyses} « e-na-li-siz ». Ne confondez pas \eng{poll} (sondage) et \eng{pole} (pôle).
\end{itemize}
\end{attention}

\begin{savaistu}[D'où vient le mot « poll » ?]
En vieil anglais, \emph{poll} désignait... la « tête » ! Compter les têtes, c'était compter les votes : d'où \eng{poll} pour « scrutin » puis « sondage », et \eng{polling station} pour « bureau de vote ». Le mot \eng{survey}, lui, vient du vieux français \emph{surveoir} (« regarder par-dessus »), entré en anglais après la conquête normande de 1066 — un petit rappel que l'anglais et le français partagent des racines profondes.
\end{savaistu}

\begin{methode}[Commenter un chiffre en trois temps]
Pour transformer un chiffre brut en phrase de commentaire (exercice fréquent en \emph{essay writing}), procédez toujours ainsi :
\begin{enumerate}
\item \textbf{Source} : citez d'où vient le chiffre — \eng{According to the survey conducted by Civic Voice, ...}
\item \textbf{Chiffre} : donnez-le avec sa structure — \eng{... 68\% of respondents said that they intended to register ...}
\item \textbf{Interprétation} : tirez la leçon du chiffre — \eng{... which shows that civic education campaigns are bearing fruit.}
\end{enumerate}
Résultat : une phrase complète, précise et nuancée, exactement ce que le correcteur attend.
\end{methode}

\begin{exoresolu}[Traduire et commenter]
Énoncé — Traduisez en anglais : « Selon le sondage, environ 55 \% des personnes interrogées ont déclaré qu'elles faisaient confiance aux médias, ce qui montre une légère amélioration. »
\tcblower
Solution détaillée — On applique la méthode en trois temps. 1) Source : \eng{According to the poll,}. 2) Chiffre : « environ » se rend par \eng{roughly} ou \eng{about} placé devant le pourcentage ; attention au temps : après \eng{said that}, le français « faisaient » appelle le past simple (concordance des temps) — \eng{roughly 55\% of respondents said that they trusted the media}. 3) Interprétation : \eng{which shows a slight improvement}. Phrase complète : \eng{According to the poll, roughly 55\% of respondents said that they trusted the media, which shows a slight improvement.} Pièges évités : pas de \emph{make confidence} (calque de « faire confiance » : on dit \eng{to trust}), et \eng{slight} est l'adjectif correspondant à l'adverbe \eng{slightly}.
\end{exoresolu}

\begin{exercice}[À vous de jouer]
Complétez avec le mot ou l'expression qui convient : \emph{findings, account for, undecided, conducted, trend, ranks}.
\begin{enumerate}
\item The institute \_\_\_\_\_\_ a survey on youth unemployment in Yaoundé.
\item Young people \_\_\_\_\_\_ 60\% of the sample.
\item The \_\_\_\_\_\_ will be released next week.
\item The North-West \_\_\_\_\_\_ first for voter registration.
\item 12\% of voters are still \_\_\_\_\_\_.
\item The general \_\_\_\_\_\_ is towards greater civic engagement.
\end{enumerate}
\end{exercice}

\begin{corrige}[Exercice « À vous de jouer »]
1. conducted (collocation : \eng{to conduct a survey}). 2. account for (représenter une proportion). 3. findings (résultats d'une enquête, toujours au pluriel). 4. ranks (\eng{to rank first} = arriver en tête ; attention à la 3\textsuperscript{e} personne : \emph{ranks}). 5. undecided (indécis). 6. trend (tendance générale).
\end{corrige}

\begin{aretenir}[L'essentiel de la section]
Un sondage se raconte avec des verbes précis : \eng{conduct/carry out a survey}, \eng{run a poll}, \eng{interview}, \eng{collect data}, \eng{analyse}, \eng{publish the findings}. Ses résultats se nomment : \eng{findings, figures, percentage, rate, majority, minority, trend}. Ils se commentent en trois temps : source (\eng{according to}), chiffre (\eng{58\% of respondents said that...}), interprétation (\eng{this shows that...}), et se nuancent avec \eng{slightly, significantly, roughly, on average}. Bannissez \emph{to make a survey} et rappelez-vous : \eng{the data show that...}.
\end{aretenir}

\coursec{Faux amis et pièges des francophones}

Dans la section précédente, nous avons appris à décrire les \cle{actions} et les \cle{résultats} d'un sondage : \eng{to conduct a survey}, \eng{to poll respondents}, \eng{the results show that...}. Mais attention : le vocabulaire des enquêtes regorge de mots qui \emph{ressemblent} au français sans avoir le même sens. Ces \cle{faux amis} (\eng{false friends} ou \eng{false cognates}) sont la première source d'erreurs dans les copies des élèves francophones. Cette section vous apprend à les repérer et à les éviter définitivement.

\courssub{Observer avant de conclure : deux phrases pièges}

Lisez attentivement les deux phrases suivantes, extraites d'un article sur un sondage réalisé à Douala :

\begin{exemplebox}[Observation]
\eng{Actually, 60\% of the population attended the meeting.}\par
\eng{The interviewer demanded an answer, but the respondent asked a sensitive question in return.}
\end{exemplebox}

Si vous avez compris « Actuellement, 60 \% de la population a assisté à la réunion » et « L'enquêteur a demandé une réponse », vous êtes tombé dans \emph{deux} pièges classiques. Traduction correcte :

\begin{itemize}
\item \eng{Actually, 60\% of the population attended the meeting.} = « En fait, 60 \% de la population a assisté à la réunion. »
\item \eng{The interviewer demanded an answer.} = « L'enquêteur a exigé une réponse. »
\item \eng{The respondent asked a sensitive question in return.} = « Le répondant a posé une question délicate en retour. »
\end{itemize}

\begin{attention}[Le piège n° 1 : actually]
\eng{Actually, 60\% of the population attended the meeting.} signifie « En fait, 60 \% de la population a assisté à la réunion. » Ne traduisez \textbf{jamais} \eng{actually} par « actuellement » ! Pour dire « actuellement », utilisez \eng{currently}, \eng{at present} ou \eng{at the moment} : \eng{Currently, the survey is being conducted in Yaoundé.} = « Actuellement, le sondage est mené à Yaoundé. »
\end{attention}

\courssub{Les six faux amis à connaître par cœur}

\begin{propriete}[to assist $\neq$ assister à]
\eng{To assist somebody} signifie \cle{aider} quelqu'un. Pour dire « assister à » (être présent à), on utilise \eng{to attend} (+ nom, \textbf{sans} préposition).\par
\eng{Two students assisted the interviewer with the questionnaires.} = « Deux élèves ont aidé l'enquêteur avec les questionnaires. »\par
\eng{Five hundred people attended the public debate.} = « Cinq cents personnes ont assisté au débat public. »\par
\textbf{Piège} : on ne dit jamais \emph{attend to a meeting} ; on dit \eng{attend a meeting}. (\eng{To attend to} signifie « s'occuper de » : \eng{The nurse attended to the patient.})
\end{propriete}

\begin{propriete}[actually $\neq$ actuellement]
\eng{Actually} = \cle{en fait}, \cle{en réalité} (sert à corriger ou préciser une idée). « Actuellement » se dit \eng{currently}, \eng{at present}, \eng{at the moment}.\par
\eng{Many people think turnout was low, but actually it reached 70\%.} = « Beaucoup pensent que la participation était faible, mais en fait elle a atteint 70 \%. »\par
\eng{The poll is currently being analysed.} = « Le sondage est actuellement en cours d'analyse. »
\end{propriete}

\begin{propriete}[sensible $\neq$ sensible]
\eng{Sensible} = \cle{raisonnable}, \cle{sensé}. « Sensible » (émotif, fragile, délicat) se dit \eng{sensitive}.\par
\eng{It is sensible to test the questionnaire on a small group first.} = « Il est raisonnable de tester d'abord le questionnaire sur un petit groupe. »\par
\eng{She asked a sensitive question about income.} = « Elle a posé une question délicate sur les revenus. »\par
\textbf{Remarque} : \eng{a sensitive issue} = « un sujet sensible / délicat » ; \eng{a sensitive person} = « une personne sensible ».
\end{propriete}

\begin{propriete}[to demand $\neq$ demander]
\eng{To demand} = \cle{exiger} (sens fort, autoritaire). « Demander » (simplement) se dit \eng{to ask} ; « demander qqch à qqn » = \eng{to ask somebody for something} ou \eng{to ask somebody something}.\par
\eng{The interviewer demanded an answer.} = « L'enquêteur a exigé une réponse. »\par
\eng{The pollster asked the respondent for his opinion.} = « Le sondeur a demandé son avis au répondant. »\par
\textbf{Piège} : on ne dit jamais \emph{demand to somebody}. Structure : \eng{demand something} ou \eng{demand that + subjonctif} (\eng{They demanded that the results be published.}).
\end{propriete}

\begin{propriete}[a survey $\neq$ une survie]
\eng{A survey} = \cle{une enquête}, \cle{un sondage}, \cle{une étude}. « La survie » se dit \eng{survival} ; « survivre » = \eng{to survive} ; « un survivant » = \eng{a survivor}.\par
\eng{The institute carried out a survey on voting intentions.} = « L'institut a réalisé un sondage sur les intentions de vote. »\par
\eng{The survival of the species is at stake.} = « La survie de l'espèce est en jeu. »\par
\textbf{Orthographe} : attention à ne pas écrire \emph{survay} ; le verbe est \eng{to survey} (« sonder, examiner »), prononcé avec l'accent sur la deuxième syllabe : « seur-VÉ ».
\end{propriete}

\begin{propriete}[eventually $\neq$ éventuellement]
\eng{Eventually} = \cle{finalement}, \cle{à la fin}, \cle{en fin de compte} (après une attente ou des difficultés). « Éventuellement » (= peut-être) se dit \eng{possibly}, \eng{perhaps}, \eng{maybe}, \eng{if necessary}.\par
\eng{After three reminders, the respondents eventually returned the forms.} = « Après trois relances, les répondants ont finalement renvoyé les formulaires. »\par
\eng{We could possibly extend the survey to Buea.} = « Nous pourrions éventuellement étendre l'enquête à Buea. »
\end{propriete}

\begin{center}
\begin{tabularx}{\linewidth}{|X|X|X|}
\hline
\rowcolor{popblueL} \textbf{Faux ami (anglais)} & \textbf{Vrai sens} & \textbf{Pour dire le mot français...} \\
\hline
to assist & aider & assister à = to attend \\
\hline
actually & en fait, en réalité & actuellement = currently, at present \\
\hline
sensible & raisonnable, sensé & sensible = sensitive \\
\hline
to demand & exiger & demander = to ask \\
\hline
a survey & un sondage, une enquête & la survie = survival \\
\hline
eventually & finalement & éventuellement = possibly, maybe \\
\hline
\end{tabularx}
\end{center}

\begin{popfigure}
\begin{tikzpicture}[
  fa/.style={rectangle, rounded corners, draw=poporange, fill=poporangeL, align=center, minimum width=3.4cm, minimum height=0.9cm, font=\small},
  ve/.style={rectangle, rounded corners, draw=popgreen, fill=popgreenL, align=center, minimum width=3.4cm, minimum height=0.9cm, font=\small},
  arr/.style={-{Stealth[length=3mm]}, thick, popdark}
]
\node[fa] (a1) at (0,0) {\eng{actually}};
\node[ve] (b1) at (7,0) {en fait (\emph{pas} actuellement)};
\draw[arr] (a1) -- (b1);
\node[fa] (a2) at (0,-1.4) {\eng{to assist}};
\node[ve] (b2) at (7,-1.4) {aider (\emph{pas} assister à)};
\draw[arr] (a2) -- (b2);
\node[fa] (a3) at (0,-2.8) {\eng{sensible}};
\node[ve] (b3) at (7,-2.8) {raisonnable (\emph{pas} sensible)};
\draw[arr] (a3) -- (b3);
\node[fa] (a4) at (0,-4.2) {\eng{to demand}};
\node[ve] (b4) at (7,-4.2) {exiger (\emph{pas} demander)};
\draw[arr] (a4) -- (b4);
\node[fa] (a5) at (0,-5.6) {\eng{eventually}};
\node[ve] (b5) at (7,-5.6) {finalement (\emph{pas} éventuellement)};
\draw[arr] (a5) -- (b5);
\node[fa] (a6) at (0,-7) {\eng{a survey}};
\node[ve] (b6) at (7,-7) {un sondage (\emph{pas} une survie)};
\draw[arr] (a6) -- (b6);
\end{tikzpicture}
\legende{Faux ami (à gauche) et vrai équivalent (à droite) : la flèche indique la correction à faire.}
\end{popfigure}

\courssub{Méthode : comment éviter les calques du français}

\begin{methode}[Trois réflexes anti-faux-amis]
\begin{enumerate}
\item \textbf{Repérez la pensée française.} Avant d'écrire ou de parler, demandez-vous : « Est-ce que je pense en français ? » Si le mot anglais qui vient spontanément ressemble au mot français, méfiance : c'est probablement un faux ami.
\item \textbf{Vérifiez la collocation apprise.} Ne traduisez jamais mot à mot : utilisez les blocs appris en leçon : \eng{to conduct a survey} (mener un sondage), \eng{to attend a meeting} (assister à une réunion), \eng{to ask a question} (poser une question), \eng{to carry out a poll} (réaliser une enquête).
\item \textbf{Testez le sens fort ou faible.} \eng{To demand} est plus fort que « demander » ; \eng{eventually} implique une fin après attente, pas une hypothèse. Si le degré ne convient pas, choisissez le mot simple : \eng{ask}, \eng{possibly}.
\end{enumerate}
\end{methode}

\begin{exoresolu}[Choisir le bon mot]
Énoncé — Complétez chaque phrase avec le mot correct entre parenthèses.\par
1. \eng{............... (Actually / Currently), the team is interviewing voters in Bamenda.}\par
2. \eng{Over 300 people ............... (assisted / attended) the press conference.}\par
3. \eng{It is ............... (sensible / sensitive) to avoid questions about religion.}\par
4. \eng{The journalist ............... (demanded / asked) the minister a polite question.}\par
5. \eng{The institute ............... (eventually / possibly) published the results after two weeks.}\par
6. \eng{A recent ............... (survey / survival) shows that citizens trust local radio.}
\tcblower
Solution détaillée :\par
1. \eng{Currently} — on parle d'une action en cours, « actuellement ». \eng{Actually} voudrait dire « en fait ».\par
2. \eng{attended} — « assister à » un événement = \eng{attend}. \eng{Assisted} signifierait « ont aidé ».\par
3. \eng{sensible} — « il est raisonnable » d'éviter ces questions. \eng{Sensitive} = « sensible, délicat » (pour une question ou une personne).\par
4. \eng{asked} — une question polie, simple : \eng{ask}. \eng{Demanded} = « a exigé », trop agressif ici.\par
5. \eng{eventually} — « finalement », après deux semaines d'attente. \eng{Possibly} = « éventuellement, peut-être ».\par
6. \eng{survey} — « un sondage ». \eng{Survival} = « la survie », sens impossible ici.
\end{exoresolu}

\begin{exercice}[À vous de jouer]
Traduisez en anglais, en évitant tout calque :\par
1. Actuellement, 40 \% des répondants refusent de répondre.\par
2. L'enquêteur a demandé son avis au chef du village.\par
3. Elle a assisté à la réunion sur les droits de l'homme à Yaoundé.\par
4. Les résultats du sondage seront éventuellement publiés demain.\par
5. C'est une question sensible ; soyez raisonnable.
\end{exercice}

\begin{corrige}[Exercice « À vous de jouer »]
1. \eng{Currently, 40\% of the respondents refuse to answer.}\par
2. \eng{The interviewer asked the village chief for his opinion.}\par
3. \eng{She attended the meeting on human rights in Yaoundé.}\par
4. \eng{The survey results will possibly be published tomorrow.}\par
5. \eng{It is a sensitive question; be sensible.}
\end{corrige}

\begin{savaistu}[Le faux ami star du Bac]
« Éventuellement » est l'un des faux amis les plus fréquents dans les copies du Baccalauréat : chaque année, des candidats écrivent \eng{eventually} pour « éventuellement », alors que \eng{eventually} signifie « finalement » ! Retenez l'astuce : \eng{eventual\textbf{ly}} contient l'idée d'une \textbf{fin} (\eng{in the end}), tandis que « éventuel » vient du latin \emph{eventus}, « ce qui peut arriver » — donc une possibilité : \eng{possibly}. De même, dans les épreuves de compréhension, \eng{actually} piège un candidat sur trois : dès que vous le voyez dans un texte, traduisez mentalement « en fait ».
\end{savaistu}

\begin{aretenir}[L'essentiel de la section]
\begin{itemize}
\item \eng{to assist} = aider ; « assister à » = \eng{to attend} (sans préposition).
\item \eng{actually} = en fait ; « actuellement » = \eng{currently}.
\item \eng{sensible} = raisonnable ; « sensible » = \eng{sensitive}.
\item \eng{to demand} = exiger ; « demander » = \eng{to ask}.
\item \eng{a survey} = un sondage ; « la survie » = \eng{survival}.
\item \eng{eventually} = finalement ; « éventuellement » = \eng{possibly}.
\item Réflexe d'or : si le mot anglais ressemble au français, vérifiez toujours la collocation avant de l'employer.
\end{itemize}
\end{aretenir}

\coursec{Familles de mots, prononciation et orthographe}

Dans la section précédente, nous avons débusqué les \emph{faux amis} qui guettent le francophone dans le domaine des sondages. Nous allons maintenant adopter une stratégie plus positive : au lieu de craindre les ressemblances entre le français et l'anglais, nous allons les exploiter. En effet, une grande partie du vocabulaire des \eng{polls and surveys} vient du latin, comme le français : \eng{respond}, \eng{represent}, \eng{decide}, \eng{majority}\ldots{} Apprendre ces mots \emph{par familles} — verbe, nom de personne, nom d'action, adjectif — permet de multiplier par quatre votre vocabulaire actif avec un effort minimal. C'est l'objet de cette section.

\courssub{Observer d'abord : un texte pour découvrir les familles de mots}

Lisez ce court article de presse. Soulignez mentalement tous les mots construits sur les racines \emph{respond-}, \emph{survey-}, \emph{represent-} et \emph{decid-}.

\begin{textebox}[Reading text — “A Survey That Changed Minds”]
Last month, a team of \textbf{surveyors} carried out an opinion \textbf{survey} in three secondary schools in Yaoundé. They wanted to \textbf{survey} students' attitudes towards civic education. Out of 800 questionnaires distributed, 620 \textbf{respondents} sent back a completed form, which gives a \textbf{response} rate of about 77 per cent — an unusually high figure. The organisers were pleased: “When people \textbf{respond} so willingly, it means the topic really matters to them,” said the head of the team.

The results were then analysed to check whether the sample was truly \textbf{representative} of the whole school population. Girls and boys were equally \textbf{represented}, and every class was included. “Good \textbf{representation} is the key to a reliable poll,” the statistician explained.

Interestingly, 30 per cent of the students were still \textbf{undecided} about whether civic education should be compulsory. “We have not yet \textbf{decided} our final position,” one respondent wrote. The headteacher called the survey “a \textbf{decisive} moment for our school”, because the \textbf{decision} to introduce a new citizenship club was taken directly after the results were published. The school board found the students remarkably \textbf{responsive}: within a week, more than two hundred had joined the club.
\end{textebox}

\textbf{Glossaire}

\begin{center}
\begin{tabularx}{\linewidth}{|l|l|X|}
\hline
\rowcolor{popblueL} \textbf{Mot anglais} & \textbf{Nature} & \textbf{Traduction} \\
\hline
to carry out & verbe & mener, réaliser \\
\hline
response rate & nom & taux de réponse \\
\hline
sample & nom & échantillon \\
\hline
compulsory & adjectif & obligatoire \\
\hline
headteacher & nom & proviseur, directeur d'école \\
\hline
board & nom & conseil (d'administration) \\
\hline
\end{tabularx}
\end{center}

\textbf{Questions de compréhension}
\begin{enumerate}
\item How many respondents completed the questionnaire?
\item Why was the sample considered representative?
\item What decision was taken after the survey?
\end{enumerate}

\courssub{Quatre familles de mots à maîtriser}

\textbf{1. La famille de \cle{respond}}

\begin{center}
\begin{tabular}{|l|l|l|}
\hline
\rowcolor{popblueL} \textbf{Mot} & \textbf{Nature} & \textbf{Traduction} \\
\hline
to respond & verbe & répondre, réagir \\
\hline
respondent & nom (personne) & le répondant, la personne interrogée \\
\hline
response & nom (action/chose) & la réponse, la réaction \\
\hline
responsive & adjectif & réceptif, qui réagit bien \\
\hline
\end{tabular}
\end{center}

\begin{exemplebox}[La famille de respond en phrases]
\eng{Only 500 people responded to our questionnaire.} « Seulement 500 personnes ont répondu à notre questionnaire. »

\eng{Each respondent answered twelve questions.} « Chaque répondant a répondu à douze questions. »

\eng{The response rate was very low.} « Le taux de réponse était très faible. »

\eng{The audience was very responsive to the journalist's questions.} « Le public s'est montré très réceptif aux questions du journaliste. »
\end{exemplebox}

\textbf{2. La famille de \cle{survey}}

\begin{center}
\begin{tabular}{|l|l|l|}
\hline
\rowcolor{popblueL} \textbf{Mot} & \textbf{Nature} & \textbf{Traduction} \\
\hline
survey & nom & l'enquête, le sondage, l'étude \\
\hline
to survey & verbe & enquêter, sonder, examiner \\
\hline
surveyor & nom (personne) & l'enquêteur (terrain) ; le géomètre-expert \\
\hline
\end{tabular}
\end{center}

\begin{attention}[Survey : nom ou verbe ? La prononciation change !]
Comme beaucoup de mots anglais de deux syllabes (\eng{record}, \eng{present}, \eng{conduct}), \eng{survey} change d'accent tonique selon sa nature :
\begin{itemize}
\item \textbf{Nom} : l'accent tombe sur la première syllabe — “SUR-vey” (/ˈsɜːveɪ/) : \eng{They conducted a SURvey in Douala.}
\item \textbf{Verbe} : l'accent tombe sur la deuxième syllabe — “sur-VEY” (/sɜːˈveɪ/) : \eng{They surVEYED 1,000 voters.}
\end{itemize}
Retenez la règle générale : \cle{nom = accent devant ; verbe = accent derrière}.
\end{attention}

\begin{attention}[Surveyor vs Interviewer / Pollster]
Le mot \eng{surveyor} désigne traditionnellement un \emph{géomètre-expert} (arpenteur). Dans le contexte des sondages d'opinion, on préfère souvent \eng{interviewer} (celui qui pose les questions) ou \eng{fieldworker} (agent de terrain). Le terme \eng{pollster} (sondeur, spécialiste des sondages) s'applique à l'organisme ou à l'expert qui conçoit et analyse le sondage, pas seulement à l'enquêteur de terrain.
\end{attention}

\textbf{3. La famille de \cle{represent}}

\begin{center}
\begin{tabular}{|l|l|l|}
\hline
\rowcolor{popblueL} \textbf{Mot} & \textbf{Nature} & \textbf{Traduction} \\
\hline
to represent & verbe & représenter \\
\hline
representative & adjectif & représentatif \\
\hline
representative & nom (personne) & le représentant, le délégué \\
\hline
representation & nom (action) & la représentation \\
\hline
\end{tabular}
\end{center}

\begin{exemplebox}[Represent en contexte de sondage]
\eng{The sample must be representative of the whole population.} « L'échantillon doit être représentatif de l'ensemble de la population. »

\eng{Women represent 52 per cent of the respondents.} « Les femmes représentent 52 \% des répondants. »

\eng{A class representative collected the completed forms.} « Un délégué de classe a collecté les formulaires remplis. »
\end{exemplebox}

\textbf{4. La famille de \cle{decide}}

\begin{center}
\begin{tabular}{|l|l|l|}
\hline
\rowcolor{popblueL} \textbf{Mot} & \textbf{Nature} & \textbf{Traduction} \\
\hline
to decide & verbe & décider \\
\hline
decision & nom (action) & la décision \\
\hline
undecided & adjectif & indécis \\
\hline
decisive & adjectif & décisif \\
\hline
\end{tabular}
\end{center}

\begin{exemplebox}[Decide en contexte]
\eng{Twenty per cent of voters are still undecided.} « Vingt pour cent des électeurs sont encore indécis. »

\eng{The poll was decisive: the government changed its policy.} « Le sondage a été décisif : le gouvernement a changé sa politique. »

\eng{They made a decision after reading the results.} « Ils ont pris une décision après avoir lu les résultats. »
\end{exemplebox}

\begin{attention}[Faux ami : “décisif” vs “indécis”]
En français, « décisif » et « indécis » appartiennent à la même famille. En anglais, attention : \eng{decisive} signifie « décisif » (qui tranche), tandis que « indécis » (pour une personne qui n'a pas choisi) se dit \eng{undecided}. Ne dites jamais \eng{*indecisive people} pour parler des électeurs indécis dans un sondage : le terme consacré est \eng{undecided voters} ou \eng{the undecided}.
\end{attention}

\courssub{Les suffixes qui fabriquent les mots : -er/-or/-ist et -tion/-sion}

\begin{propriete}[Former les noms de personnes : -er, -or, -ist]
Pour nommer la \textbf{personne qui fait l'action}, l'anglais ajoute le plus souvent un suffixe au verbe ou au nom de base :
\begin{itemize}
\item \textbf{-er} (le plus fréquent) : \eng{interview} $\rightarrow$ \eng{interviewer} (l'intervieweur) ; \eng{teach} $\rightarrow$ \eng{teacher} ; \eng{employ} $\rightarrow$ \eng{employer}.
\item \textbf{-or} (souvent après des racines latines) : \eng{survey} $\rightarrow$ \eng{surveyor} (l'enquêteur/arpenteur) ; \eng{conduct} $\rightarrow$ \eng{conductor} ; \eng{investigate} $\rightarrow$ \eng{investigator}.
\item \textbf{-ist} (spécialiste) : \eng{poll} $\rightarrow$ \eng{pollster} (le sondeur) ; \eng{statistics} $\rightarrow$ \eng{statistician} (le statisticien) ; \eng{journal} $\rightarrow$ \eng{journalist}.
\end{itemize}
Cas particulier à mémoriser : \eng{respond} $\rightarrow$ \eng{respondent} (suffixe latin \textbf{-ent}, pas de -er).
\end{propriete}

\begin{propriete}[Former les noms d'action : -tion, -sion, -sation]
Pour nommer \textbf{l'action ou son résultat}, l'anglais utilise surtout \textbf{-tion} et \textbf{-sion}, exactement comme le français :
\begin{itemize}
\item \eng{decide} $\rightarrow$ \eng{decision} (décision) — le \emph{d} final devient \emph{s}.
\item \eng{represent} $\rightarrow$ \eng{representation} (représentation).
\item \eng{conclude} $\rightarrow$ \eng{conclusion} ; \eng{divide} $\rightarrow$ \eng{division} (même transformation \emph{d} $\rightarrow$ \emph{s}).
\item \eng{organise} $\rightarrow$ \eng{organisation} ; \eng{analyse} $\rightarrow$ \eng{analysis} (irrégulier !).
\end{itemize}
Ces noms se construisent souvent avec un verbe support : \eng{to make a decision} (prendre une décision), \eng{to draw a conclusion} (tirer une conclusion), \eng{to carry out a survey} (mener une enquête).
\end{propriete}

\begin{popfigure}
\begin{center}
\begin{tikzpicture}[
  grow=right,
  level distance=3.6cm,
  level 1/.style={sibling distance=2.6cm},
  level 2/.style={sibling distance=1.2cm},
  edge from parent/.style={draw=popmuted, thick},
  every node/.style={font=\small}
]
\node[draw=popblue, fill=popblueL, rounded corners, thick, align=center, font=\bfseries] {to respond\\ \emph{répondre}}
  child { node[draw=poppurple, fill=poppurpleL, rounded corners, align=center] {responsive\\ \emph{adj. : réceptif}} }
  child { node[draw=popgreen, fill=popgreenL, rounded corners, align=center] {response\\ \emph{nom : la réponse}} }
  child { node[draw=poporange, fill=poporangeL, rounded corners, align=center] {respondent\\ \emph{nom : le répondant}} };
\end{tikzpicture}
\end{center}
\legende{Arbre lexical de la racine \emph{respond} : une racine, trois dérivés à connaître par cœur.}
\end{popfigure}

\courssub{Prononciation et orthographe : les cinq pièges à corriger}

\begin{center}
\begin{tabularx}{\linewidth}{|l|X|X|}
\hline
\rowcolor{popblueL} \textbf{Mot} & \textbf{Prononciation (en lettres françaises)} & \textbf{Piège à éviter} \\
\hline
survey (nom) & “SEU-veu” (accent devant) & ne pas accentuer la fin pour le nom \\
\hline
to survey (verbe) & “seu-VEU” (accent derrière) & ne pas dire “SUR-vey” pour le verbe \\
\hline
questionnaire & “kess-tio-NEUR” & deux \textbf{n}, comme en français \\
\hline
data & “DÉI-ta” & ne pas dire “da-ta” à la française \\
\hline
percentage & “peur-SEN-tidj” & un seul mot, jamais “per centage” \\
\hline
majority & “ma-DJO-ri-ti” & accent sur la 2\textsuperscript{e} syllabe ; \emph{g} se prononce “dj” \\
\hline
\end{tabularx}
\end{center}

\begin{attention}[Orthographe : trois mots à surveiller]
\begin{enumerate}
\item \eng{questionnaire} s'écrit avec \textbf{deux n} en anglais comme en français — mais ne vous fiez pas à la prononciation française : dites “kess-tio-NEUR” (approximation : /ˌkwɛstʃəˈnɛə/).
\item \eng{respondent} se termine en \textbf{-ent} (et non \emph{-ant}) : c'est un héritage du latin \emph{respondens}. De même : \eng{correspondent}, \eng{president}.
\item \eng{percentage} s'écrit en \textbf{un seul mot}. En revanche, \eng{per cent} (deux mots) s'emploie après un chiffre : \eng{70 per cent of respondents} = 70 \% des répondants. Comparez : \eng{A large percentage of voters} / \eng{70 per cent of voters}.
\end{enumerate}
\end{attention}

\courssub{Méthode et entraînement}

\begin{methode}[Apprendre par familles, pas par listes]
\begin{enumerate}
\item Choisissez une \textbf{racine} (ex. \eng{respond}) et écrivez-la en haut d'une colonne verticale.
\item Listez en dessous toutes ses formes : verbe, nom de personne, nom d'action, adjectif.
\item Pour \textbf{chaque forme}, écrivez une phrase complète tirée du thème des sondages — jamais le mot isolé.
\item Testez-vous en cachant la colonne : donnez les quatre formes à partir de la racine, puis refaites vos phrases à l'oral.
\item Révisez en croisant les familles : \eng{The respondent's response was decisive.} — une seule phrase, deux familles révisées !
\end{enumerate}
\end{methode}

\begin{exoresolu}[Complétez avec la bonne forme de la famille]
Énoncé — Complétez chaque phrase avec la forme correcte du mot entre parenthèses :
\begin{enumerate}
\item Only 60 \dots{} out of 100 sent back the form. (respond)
\item The \dots{} rate of this poll was excellent. (respond)
\item The sample must be \dots{} of the whole country. (represent)
\item Many voters are still \dots{} about the reform. (decide)
\item The \dots{} interviewed people in the streets of Bamenda. (survey)
\end{enumerate}
\tcblower
Solution détaillée :
\begin{enumerate}
\item \eng{respondents} — il faut un \textbf{nom de personne} au pluriel (60 personnes). « Seulement 60 répondants sur 100 ont renvoyé le formulaire. »
\item \eng{response} — il faut le \textbf{nom d'action} : la collocation consacrée est \eng{response rate} (taux de réponse).
\item \eng{representative} — il faut l'\textbf{adjectif} après \eng{must be} : « représentatif ».
\item \eng{undecided} — il faut l'adjectif négatif : « indécis » ; le contexte (\eng{still}) confirme que le choix n'est pas fait.
\item \eng{surveyor} — il faut le \textbf{nom de personne} au singulier (sujet de \eng{interviewed}) : celui qui mène l'enquête.
\end{enumerate}
\end{exoresolu}

\begin{exercice}[À vous de jouer]
\begin{enumerate}
\item Construisez la famille complète (verbe, nom de personne, nom d'action, adjectifs) de \eng{decide} et de \eng{represent}, avec une phrase pour chaque forme.
\item Traduisez : « Le sondeur a mené une enquête auprès de 500 répondants, mais le taux de réponse était faible et l'échantillon n'était pas représentatif. »
\item Dites à voix haute, en plaçant correctement l'accent tonique : \eng{a survey} / \eng{to survey} / \eng{questionnaire} / \eng{majority} / \eng{percentage}.
\end{enumerate}
\end{exercice}

\begin{corrige}[Exercice — éléments de réponse]
\begin{enumerate}
\item \eng{decide} : to decide / decision / undecided / decisive. Exemple : \eng{The commission decided to publish the results.} ; \eng{represent} : to represent / representative (noun \& adj.) / representation. Exemple : \eng{The sample is not representative.}
\item \eng{The pollster carried out a survey of 500 respondents, but the response rate was low and the sample was not representative.}
\item “SEU-veu” (nom) / “seu-VEU” (verbe) / “kess-tio-NEUR” / “ma-DJO-ri-ti” / “peur-SEN-tidj”.
\end{enumerate}
\end{corrige}

\begin{experience}[Prendre la parole : présenter des résultats]
En binôme, simulez une restitution orale de sondage (2 minutes). L'élève A est le \eng{pollster}, l'élève B est le journaliste.
\begin{itemize}
\item \textbf{Pollster} : « \eng{We carried out a survey of 1,000 voters. The response rate was 80 per cent. The sample is representative. 45 per cent support the reform, 30 per cent oppose it, and 25 per cent are undecided. The results are decisive for the campaign.} »
\item \textbf{Journaliste} : Posez deux questions de relance (ex. \eng{What about the margin of error? / How did you select the respondents?}).
\end{itemize}
Changez de rôle. Utilisez impérativement : \eng{sample, representative, respondent, response rate, undecided, decisive}.
\end{experience}

\begin{savaistu}[Le sondage au Cameroun et dans le monde]
Au Cameroun, les instituts comme \eng{AFROBAROMETER} ou \eng{IPSOS} mènent des enquêtes (\eng{surveys}) sur la gouvernance, l'économie ou la santé. La méthodologie suit les standards internationaux : \eng{random sampling} (échantillonnage aléatoire), \eng{face-to-face interviews} (entretiens en face-à-face) dans les zones urbaines (Douala, Yaoundé) et rurales. Au Royaume-Uni, \eng{YouGov} et \eng{Ipsos MORI} sont célèbres ; aux USA, \eng{Gallup} et \eng{Pew Research Center}. Le terme \eng{pollster} (sondeur) est très utilisé dans la presse anglo-saxonne (\eng{The pollsters predicted a tight race}).
\end{savaistu}

\begin{aretenir}[L'essentiel de la section]
\begin{itemize}
\item Quatre familles à connaître par cœur : \eng{respond / respondent / response / responsive} ; \eng{survey / to survey / surveyor} ; \eng{represent / representative / representation} ; \eng{decide / decision / undecided / decisive}.
\item Noms de personnes en \textbf{-er / -or / -ist} : \eng{interviewer}, \eng{surveyor}, \eng{pollster}, \eng{statistician} ; noms d'action en \textbf{-tion / -sion} : \eng{decision}, \eng{representation}, \eng{conclusion}.
\item \eng{survey} : accent devant pour le nom (\eng{SURvey}), derrière pour le verbe (\eng{surVEY}).
\item Orthographe : \eng{questionnaire} (deux n), \eng{respondent} (-ent), \eng{percentage} (un seul mot) vs \eng{per cent} (deux mots après un chiffre).
\item Stratégie gagnante : apprendre les mots \textbf{en colonne, par famille}, avec une phrase d'exemple pour chaque forme.
\end{itemize}
\end{aretenir}

\coursec{Emploi en contexte : lire, dire et écrire un sondage}

Dans la section précédente, nous avons exploré les \cle{familles de mots} (\eng{to poll, a poll, a pollster} ; \eng{to survey, a survey}), la \cle{prononciation} et l'\cle{orthographe} du vocabulaire des enquêtes. Il est temps maintenant de \textbf{mettre tout ce lexique en action} : un mot n'est vraiment acquis que lorsqu'on sait le reconnaître dans un texte, l'employer dans une conversation et le réutiliser dans une production écrite. C'est exactement ce que nous allons faire ici, en trois temps : \textbf{lire} un article de presse, \textbf{dire} à travers un dialogue, puis \textbf{écrire} un commentaire de sondage.

\courssub{Lire : un article de presse avec résultats de sondage}

Commençons par observer comment la presse anglophone présente les résultats d'une enquête. Lisez l'article suivant, puis repérez les mots du champ lexical étudiés dans cette leçon.

\begin{textebox}[Youth and Citizenship: What the Poll Says — The Cameroon Tribune (fictional article)]
A new opinion poll published this week sheds light on how young Cameroonians view their role as citizens. The survey, conducted by the National Institute of Social Research, was carried out in October among a representative sample of 1,200 people aged 18 to 25, in Yaoundé, Douala, Buea and Bamenda.

Respondents were asked to fill in a questionnaire about voting, volunteering and trust in public institutions. According to the findings, 58\% of those interviewed intend to vote in the next elections, while 27\% say they are still undecided. Only 15\% declare they will not vote at all.

The results also reveal a growing interest in community life: nearly two thirds of the respondents have taken part in at least one volunteer activity, such as cleaning campaigns or blood donation drives. However, the poll shows that many young people feel their voices are not heard by decision-makers.

“The figures are encouraging, but there is still work to do,” explains Dr. Manka, the lead researcher. “With a margin of error of plus or minus three points, we can say that youth engagement is slowly rising.”

The survey was conducted face-to-face and by telephone. The response rate was high, which makes the data particularly reliable, according to the pollsters.
\end{textebox}

\textbf{Glossaire (mots difficiles)}

\begin{center}
\begin{tabularx}{\linewidth}{|l|l|X|}
\hline
\rowcolor{popblueL} Mot anglais & Nature & Traduction française \\
\hline
to shed light on & verbe & éclairer, mettre en lumière \\
\hline
to carry out & verbe & mener, réaliser (une enquête) \\
\hline
respondent & nom & personne interrogée, répondant \\
\hline
to fill in & verbe & remplir (un questionnaire) \\
\hline
undecided & adjectif & indécis \\
\hline
volunteer activity & nom & activité bénévole \\
\hline
margin of error & nom & marge d'erreur \\
\hline
response rate & nom & taux de réponse \\
\hline
reliable & adjectif & fiable, digne de confiance \\
\hline
\end{tabularx}
\end{center}

\textbf{Questions de compréhension}

\begin{enumerate}
\item Who conducted the survey, and when?
\item How large was the sample, and who were the respondents?
\item What percentage of young people intend to vote?
\item What is the margin of error mentioned by Dr. Manka?
\item Why is the data considered reliable?
\end{enumerate}

\begin{corrige}[Questions de compréhension]
\begin{enumerate}
\item The survey was conducted by the National Institute of Social Research, in October.
\item The sample consisted of 1,200 people aged 18 to 25, living in Yaoundé, Douala, Buea and Bamenda.
\item 58 per cent of young people intend to vote.
\item The margin of error is plus or minus three points.
\item The data is reliable because the response rate was high.
\end{enumerate}
\end{corrige}

\courssub{Dire : un dialogue modèle}

Passons à l'oral. Voici un dialogue type entre un journaliste et un élève qui a participé à une enquête scolaire. Lisez-le à voix haute, puis jouez-le à deux.

\begin{textebox}[Dialogue modèle — At the school radio]
Journalist: “Did you take part in the survey?”\\
Student: “Yes, I answered a questionnaire about voting.”\\
Journalist: “What were the main findings?”\\
Student: “According to the results, 70\% of students intend to vote, but many are still undecided.”\\
Journalist: “Interesting! And how was the poll conducted?”\\
Student: “The pollsters interviewed a sample of two hundred students from all classes.”\\
Journalist: “Thank you for sharing these figures with us.”
\end{textebox}

\begin{experience}[Jeu de rôle : l'interview du sondage]
Par groupes de deux : l'élève A est journaliste à la radio de l'école, l'élève B a participé à une enquête sur un sujet de votre choix (le sport, la cantine, le vote des délégués). Préparez cinq questions avec le lexique de la leçon (\eng{survey, questionnaire, sample, findings, respondents, per cent}), puis jouez la scène devant la classe. Changez ensuite les rôles.
\end{experience}

\begin{savaistu}[Les instituts de sondage au Cameroun et dans le monde]
Au Cameroun, l'\eng{National Institute of Statistics (NIS)} — Institut national de la statistique — réalise les grandes enquêtes démographiques et économiques (recensement, enquête sur l'emploi, etc.). Pour les sondages d'opinion politique ou sociale, on fait souvent appel à des cabinets privés ou à des ONG. Au Royaume-Uni, les noms célèbres sont \eng{YouGov, Ipsos MORI, Survation} ; aux États-Unis, \eng{Gallup, Pew Research Center}. Ces instituts utilisent des méthodes rigoureuses : \eng{random sampling} (échantillonnage aléatoire), \eng{weighting} (redressement) pour que l'échantillon soit \eng{representative} de la population.
\end{savaistu}

\courssub{Écrire : la méthode du commentaire de sondage}

Observer l'article ci-dessus montre qu'un commentaire de sondage suit toujours la même logique. Voici la démarche à suivre.

\begin{methode}[Rédiger un commentaire de sondage en trois étapes]
\begin{enumerate}
\item \textbf{Présenter la source et l'échantillon} : qui a mené l'enquête, quand, auprès de combien de personnes.\\
\eng{A survey conducted by... among a sample of 1,200 people shows that...}
\item \textbf{Donner deux ou trois chiffres clés} avec \eng{according to} et les pourcentages.\\
\eng{According to the findings, 58 per cent of respondents intend to vote, while 27 per cent are undecided.}
\item \textbf{Conclure par une tendance générale} : une phrase d'interprétation, sans chiffre.\\
\eng{This shows that young people are becoming more interested in public life.}
\end{enumerate}
\end{methode}

\begin{exemplebox}[Phrases modèles traduites]
\eng{The findings show that young people are increasingly interested in politics.} « Les résultats montrent que les jeunes s'intéressent de plus en plus à la politique. »

\eng{According to the poll, two thirds of the respondents trust their local leaders.} « Selon le sondage, les deux tiers des personnes interrogées font confiance à leurs responsables locaux. »

\eng{The survey was carried out among a representative sample of 500 students.} « L'enquête a été menée auprès d'un échantillon représentatif de 500 élèves. »

\eng{With a margin of error of three points, the results remain reliable.} « Avec une marge d'erreur de trois points, les résultats restent fiables. »
\end{exemplebox}

\textbf{Tableau récapitulatif : les phrases utiles pour commenter un sondage}

\begin{center}
\begin{tabularx}{\linewidth}{|X|X|}
\hline
\rowcolor{popblueL} Fonction & Expression anglaise \\
\hline
Présenter la source & \eng{A poll conducted by... / A survey carried out among...} \\
\hline
Citer les résultats & \eng{According to the findings/results, ...} \\
\hline
Donner un chiffre & \eng{58 per cent of respondents... / nearly two thirds...} \\
\hline
Comparer & \eng{while... / whereas... / only 15 per cent...} \\
\hline
Nuancer & \eng{with a margin of error of... / the figures suggest that...} \\
\hline
Conclure & \eng{This shows that... / This trend reveals that...} \\
\hline
\end{tabularx}
\end{center}

\begin{attention}[Prononciation et orthographe des pourcentages]
À l'oral, articulez bien les pourcentages : on dit \eng{fifty-eight per cent} (« fif-ti ète per sente ») et non \emph{fifty-eight persons} — \eng{per cent} ne veut pas dire « personnes » ! À l'écrit, attention : en anglais britannique, \cle{per cent} s'écrit en \textbf{deux mots} (en anglais américain, on trouve \eng{percent} en un seul mot). Le symbole \% s'écrit collé au chiffre : \eng{58\%}. Enfin, ne dites jamais \emph{the pourcentage} : le mot est \eng{percentage}.
\end{attention}

\courssub{Reformuler : transformer des chiffres en phrases}

Voici les résultats fictifs d'un sondage sur le vote des jeunes au Cameroun. Votre tâche : transformer ce graphique en phrases complètes.

\begin{popfigure}
\begin{center}
\begin{tikzpicture}[scale=0.9]
\draw[->, thick, popmuted] (0,0) -- (0,4.2) node[above, font=\small] {\%};
\draw[->, thick, popmuted] (0,0) -- (8.6,0);
\fill[popblue] (0.7,0) rectangle (2.1,3.19);
\node[font=\small, align=center] at (1.4,-0.5) {Will\\vote};
\node[font=\small\bfseries, popblue] at (1.4,3.45) {58\%};
\fill[poporange] (3.4,0) rectangle (4.8,1.485);
\node[font=\small, align=center] at (4.1,-0.5) {Undecided};
\node[font=\small\bfseries, poporange] at (4.1,1.75) {27\%};
\fill[poppurple] (6.1,0) rectangle (7.5,0.825);
\node[font=\small, align=center] at (6.8,-0.5) {Will not\\vote};
\node[font=\small\bfseries, poppurple] at (6.8,1.1) {15\%};
\foreach \y/\n in {0.55/10, 1.1/20, 1.65/30, 2.2/40, 2.75/50, 3.3/60}
  \draw[popline] (0,\y) -- (8.4,\y) node[pos=0, left, font=\scriptsize, popmuted] {\n};
\end{tikzpicture}
\end{center}
\legende{Résultats fictifs du sondage « Youth and Voting » (échantillon : 1,200 jeunes de 18 à 25 ans)}
\end{popfigure}

\begin{exoresolu}[Des chiffres aux phrases]
Énoncé : à l'aide du graphique ci-dessus, écrivez une phrase complète pour chaque option, en utilisant \eng{according to}, \eng{while} et \eng{only}.
\tcblower
Solution détaillée :
\begin{itemize}
\item \eng{According to the poll, 58 per cent of young people intend to vote in the next elections.} « Selon le sondage, 58 \% des jeunes ont l'intention de voter aux prochaines élections. »
\item \eng{27 per cent of the respondents are still undecided.} « 27 \% des personnes interrogées sont encore indécises. »
\item \eng{Only 15 per cent say they will not vote at all.} « Seulement 15 \% déclarent qu'ils ne voteront pas du tout. »
\end{itemize}
On peut ensuite relier les trois phrases : \eng{According to the poll, 58 per cent of young people intend to vote, while 27 per cent are undecided and only 15 per cent refuse to vote.}
\end{exoresolu}

\begin{exercice}[Phrases à compléter]
Complétez avec le mot juste : \eng{sample — margin of error — findings — questionnaire — respondents — conducted}.
\begin{enumerate}
\item The survey was \dots\ among 1,200 young people in four cities.
\item Each participant had to fill in a \dots\ of twenty questions.
\item The \dots\ of the poll were published in the newspaper.
\item The \dots\ were aged between 18 and 25.
\item With a \dots\ of three points, the results are reliable.
\item The researchers selected a representative \dots\ of the population.
\end{enumerate}
\end{exercice}

\begin{corrige}[Phrases à compléter]
\begin{enumerate}
\item \eng{conducted} — l'enquête a été \emph{menée}.
\item \eng{questionnaire} — remplir un \emph{questionnaire}.
\item \eng{findings} — les \emph{résultats} du sondage.
\item \eng{respondents} — les \emph{personnes interrogées}.
\item \eng{margin of error} — une \emph{marge d'erreur} de trois points.
\item \eng{sample} — un \emph{échantillon} représentatif.
\end{enumerate}
\end{corrige}

\courssub{Mini-production écrite guidée}

À vous maintenant ! Rédigez \textbf{cinq phrases} commentant les résultats fictifs du sondage « Youth and Voting » présenté dans le graphique. Suivez la méthode en trois étapes : source et échantillon, chiffres clés, tendance générale.

\begin{exemplebox}[Modèle de production (à ne pas recopier mot à mot)]
\eng{A survey on youth and voting was conducted among a sample of 1,200 young Cameroonians aged 18 to 25. According to the findings, 58 per cent of the respondents intend to vote in the next elections. However, 27 per cent are still undecided, and only 15 per cent say they will not vote. These figures suggest that most young people want to take part in the political life of their country. This trend shows that citizenship education in schools is bearing fruit.}
\end{exemplebox}

\begin{attention}[Dernières erreurs à éviter dans votre production]
\begin{itemize}
\item Ne traduisez pas « selon le sondage » par \emph{following the poll} : dites \eng{according to the poll}.
\item \eng{per cent} en deux mots (britannique) ; le symbole \% collé au chiffre.
\item Pas de \emph{s} à \eng{per cent} : \eng{58 per cent}, jamais \emph{58 per cents}.
\item Après \eng{58 per cent of the respondents}, le verbe se met au pluriel : \eng{58 per cent of respondents \textbf{intend} to vote}.
\end{itemize}
\end{attention}

Vous savez désormais \textbf{lire}, \textbf{dire} et \textbf{écrire} un sondage en anglais : compétence précieuse pour comprendre la presse anglophone et pour le baccalauréat, où les documents statistiques sont fréquents.

\newpage

\coursec{Activité d'intégration}

\courssub{Objectifs de l'activité}

À la fin de cette activité d'intégration, l'élève sera capable de :

\begin{itemize}
\item concevoir un \cle{questionnaire} court et correct en anglais (questions fermées à choix multiples et question ouverte) ;
\item mener un mini-sondage auprès d'un échantillon de camarades (\eng{to conduct a survey}, \eng{to interview respondents}) ;
\item collecter, compter et présenter des résultats en pourcentages (\eng{out of 20 respondents, 65\% said...}) ;
\item rédiger un court rapport de sondage de 6 à 8 phrases en réemployant le lexique de la leçon (\eng{according to our survey}, \eng{findings}, \eng{a majority of}, \eng{to be in favour of}) ;
\item présenter oralement les résultats à la classe et comparer les tendances (\eng{trends}) entre groupes ;
\item éviter les faux amis étudiés (\eng{actually} $\neq$ « actuellement », \eng{a survey} $\neq$ « une surveillance », \eng{to assist} $\neq$ « assister à »).
\end{itemize}

\courssub{Situation}

Le club d'anglais de ton lycée à Yaoundé prépare un numéro spécial du journal scolaire consacré à la citoyenneté des jeunes. Le rédacteur en chef a reçu la lettre suivante d'une ONG partenaire :

\begin{textebox}[Letter from the partner NGO]
Dear students,

Youth Voice Cameroon is preparing its annual report on young citizens' attitudes. We would like your school to contribute a short article based on a real opinion poll. Your task: design a five-question survey on citizenship, interview at least twenty respondents, and send us a written report of six to eight sentences presenting your main findings with percentages.

Suggested topics: voting at eighteen, trust in the media, community work, respect for human rights, national languages.

We look forward to reading your results.

Best regards,

The Youth Voice Team
\end{textebox}

\begin{definition}[Glossaire utile]
\begin{itemize}
\item \cle{annual report} (n.) : rapport annuel ; \cle{attitude} (n.) : attitude, opinion.
\item \cle{opinion poll} (n.) : sondage d'opinion ; \cle{respondent} (n.) : personne interrogée, répondant.
\item \cle{findings} (n. pl.) : résultats, conclusions d'une enquête.
\item \cle{community work} (n.) : travail communautaire, bénévolat.
\end{itemize}
\end{definition}

Ta classe relève le défi : par groupes de quatre, vous allez construire le sondage, interroger des élèves, analyser les données et rédiger l'article demandé.

\courssub{Consignes}

\textbf{Étape 1 — Comprendre la demande.} Lis la lettre de l'ONG et relève : (a) le nombre de questions attendu, (b) la taille minimale de l'échantillon (\eng{sample}), (c) la longueur du rapport, (d) les thèmes suggérés.

\textbf{Étape 2 — Concevoir le questionnaire.} En groupe, rédigez cinq questions sur la citoyenneté des jeunes :
\begin{enumerate}
\item trois questions fermées à choix multiples (\eng{multiple-choice questions}) avec les options \eng{yes / no / not sure} ou trois propositions ;
\item une question d'opinion avec échelle (\eng{Do you strongly agree, agree, disagree or strongly disagree?}) ;
\item une question ouverte (\eng{open question}) commençant par \eng{Why...?} ou \eng{What...?}.
\end{enumerate}
Exemples de questions possibles : \eng{Do you intend to vote when you turn eighteen?} — \eng{Do you trust the media?} — \eng{Would you do voluntary work for your community?} — \eng{Why do some young people avoid politics?}

\textbf{Étape 3 — Mener le sondage.} Interrogez au moins vingt élèves (d'une autre classe si possible). Utilisez les formules polies : \eng{Excuse me, we are conducting a survey on citizenship. Could you answer five questions? It will only take two minutes.} Notez chaque réponse sur une grille (\eng{answer sheet}).

\textbf{Étape 4 — Dépouiller les résultats.} Comptez les réponses (\eng{to count the answers}), calculez les pourcentages (\eng{to work out the percentages}) et remplissez un tableau de résultats au tableau :

\begin{center}
\begin{tabularx}{\linewidth}{|X|X|X|}\hline
\rowcolor{popblueL} Question & Yes & No / Not sure \\ \hline
1. Vote at 18? & 14 (70\%) & 6 (30\%) \\ \hline
2. Trust the media? & 8 (40\%) & 12 (60\%) \\ \hline
\end{tabularx}
\end{center}

\textbf{Étape 5 — Rédiger le rapport.} Écrivez six à huit phrases en respectant le plan : introduction du sondage (\eng{we conducted a survey among...}), présentation de l'échantillon, deux ou trois résultats chiffrés, une tendance générale (\eng{the findings show that...}), une conclusion. Attention aux prépositions : \eng{according to}, \eng{in favour of}, \eng{interested in}, \eng{to participate in}.

\textbf{Étape 6 — Présenter et comparer.} Chaque groupe présente oralement son rapport en deux minutes. La classe compare les \eng{findings} : quelles tendances se répètent ? Quels résultats surprennent (\eng{surprising results}) ?

\courssub{Ressources à mobiliser}

\begin{aretenir}[Boîte à outils du rapport de sondage]
\begin{itemize}
\item \textbf{Introduire} : \eng{We conducted a survey among twenty students of our school.} / \eng{Our sample consisted of...}
\item \textbf{Citer un résultat} : \eng{According to our survey, 70\% of the respondents said that...} / \eng{Out of twenty people interviewed, fourteen intend to vote.}
\item \textbf{Nuancer} : \eng{a majority of} (une majorité de), \eng{a minority of}, \eng{nearly half of}, \eng{only a few}.
\item \textbf{Conclure} : \eng{The findings show that young people are interested in...} / \eng{The main trend is...}
\item \textbf{Grammaire} : discours rapporté au passé (\eng{they said (that) they would vote}), pourcentages + \eng{of}, prépositions figées.
\end{itemize}
\end{aretenir}

\begin{attention}[Pièges à éviter dans le rapport]
\begin{itemize}
\item \eng{Actually} signifie « en fait », jamais « actuellement » : dis \eng{nowadays} ou \eng{at present}.
\item \eng{A survey} est un sondage ; « une surveillance » se dit \eng{surveillance}.
\item On dit \eng{65\% of the respondents} (avec \eng{of}), jamais \eng{the 65\% of respondents}.
\item \eng{To participate in}, \eng{to be interested in}, \eng{to be in favour of} : ne oublie pas la préposition.
\end{itemize}
\end{attention}

\courssub{Proposition de production et corrigé commenté}

\begin{exoresolu}[Modèle de rapport : « Young Citizens' Opinions »]
Rédige le rapport de 6 à 8 phrases demandé par l'ONG, à partir des résultats suivants : 20 élèves interrogés ; 70\% veulent voter à 18 ans ; 40\% font confiance aux médias ; 80\% accepteraient un travail bénévole ; raison principale du désintérêt politique : la méfiance.
\tcblower
\textbf{Production modèle :}

\eng{Last week, our group conducted a survey on citizenship among twenty students of our school. According to our findings, 70\% of the respondents said they intend to vote when they turn eighteen, which shows that most young people want to participate in elections. However, only 40\% of the people interviewed trust the media, and a majority believe that the news is not always reliable. Surprisingly, 80\% of our sample declared that they would do voluntary work for their community. When we asked why some young people avoid politics, most respondents answered that they do not trust politicians. In conclusion, the main trend is clear: young citizens are interested in public life, but they need leaders they can believe in.}

\textbf{Commentaire des choix de langue :}
\begin{itemize}
\item \eng{conducted a survey among twenty students} : collocation exacte (\eng{to conduct a survey}) et préposition \eng{among} pour l'échantillon.
\item \eng{70\% of the respondents said they intend to vote} : pourcentage + \eng{of}, discours rapporté sans \eng{that} obligatoire, verbe au présent car l'intention est encore vraie.
\item \eng{However, only 40\%...} : connecteur de contraste pour nuancer ; \eng{only} placé avant le chiffre.
\item \eng{a majority believe} : accord au pluriel après \eng{a majority of + nom pluriel}.
\item \eng{would do voluntary work} : conditionnel de discours rapporté ; \eng{voluntary work} est le terme correct pour « bénévolat ».
\item \eng{In conclusion, the main trend is...} : formule de clôture typique d'un rapport, avec le mot-clé \eng{trend} de la leçon.
\end{itemize}
\end{exoresolu}

\courssub{Bilan de l'activité}

Cette tâche a permis de réunir toutes les compétences de la leçon dans une situation authentique : comprendre une consigne réelle (la lettre de l'ONG), mobiliser le champ lexical des sondages (\eng{poll, sample, respondent, findings, trend, percentage}), appliquer la grammaire du discours rapporté et des pourcentages, puis produire à l'écrit et à l'oral un rapport structuré. Les erreurs les plus fréquentes observées en classe — faux amis (\eng{actually}, \eng{survey/surveillance}), oubli de \eng{of} après les pourcentages, prépositions manquantes — doivent être relevées collectivement après les présentations. Pour aller plus loin, le meilleur rapport peut être envoyé au journal scolaire ou affiché au tableau d'information du club d'anglais : le sondage devient alors un véritable acte de citoyenneté.

\newpage

\coursec{Méthodes et exercices résolus}

\coursec{Lesson 45 — Vocabulary: Words and expressions related to polls/surveys}

\courssub{Mise en route : découvrir le monde des sondages}

\begin{textebox}[Script d'écoute — Extrait d'un journal télévisé camerounais (CRTV News)]
“Good evening. A new opinion poll conducted by the National Institute of Statistics reveals that unemployment remains the top concern for young Cameroonians. The survey, carried out last month among a representative sample of two thousand respondents aged fifteen to thirty-five, shows that sixty-five per cent of them consider job creation a priority. The findings also highlight a growing interest in vocational training. The margin of error is three percentage points. More details after the break.”
\end{textebox}

\begin{experience}[Activité de classe : Sondage éclair]
En binôme, interrogez votre voisin sur ses habitudes de lecture (journaux, réseaux sociaux, livres). Utilisez les questions types : \eng{“How often do you…?”}, \eng{“Do you prefer…?”}, \eng{“What is your opinion on…?”}. Notez les réponses. Ensuite, présentez oralement les “findings” de votre mini-sondage à la classe en utilisant \eng{“According to my survey…”}, \eng{“The majority of respondents…”}, \eng{“One out of three…”}.
\end{experience}

\courssub{Champ lexical 1 : les acteurs et les outils du sondage}

\begin{propriete}[Acteurs (People) — Vocabulaire de base]
\begin{center}
\begin{tabularx}{\linewidth}{|X|X|X|}
\hline
\rowcolor{popblueL} \textbf{English} & \textbf{Français} & \textbf{Note / Collocation} \\
\hline
a \eng{pollster} & un sondeur, un expert en sondages & Spécialiste qui conçoit et analyse les sondages. \\
\hline
a \eng{respondent} & une personne interrogée, un répondant & Celui qui répond au questionnaire. \textbf{Pas} \eng{respondant}. \\
\hline
an \eng{interviewer} & un enquêteur, un intervieweur & Celui qui pose les questions (en face-à-face ou téléphone). \\
\hline
an \eng{interviewee} & l'interviewé, la personne interrogée & Suffixe \emph{-ee} = celui qui subit l'action. \\
\hline
a \eng{participant} / a \eng{taker} & un participant & \eng{to take part in a survey}, \eng{to participate in a poll}. \\
\hline
the \eng{target population} / the \eng{target audience} & la population cible & Groupe que l'on veut étudier. \\
\hline
a \eng{sample} & un échantillon & \eng{a representative sample}, \eng{a random sample} (aléatoire), \eng{a biased sample} (biaisé). \\
\hline
\end{tabularx}
\end{center}
\end{propriete}

\begin{propriete}[Outils et supports (Tools) — Vocabulaire de base]
\begin{center}
\begin{tabularx}{\linewidth}{|X|X|X|}
\hline
\rowcolor{popblueL} \textbf{English} & \textbf{Français} & \textbf{Note / Collocation} \\
\hline
a \eng{questionnaire} & un questionnaire & \textbf{Deux « n »}. \eng{to design a questionnaire}, \eng{to fill in/out a questionnaire}. \\
\hline
a \eng{survey} / a \eng{poll} & une enquête, un sondage & \eng{survey} plus large (souvent écrit) ; \eng{poll} souvent politique/opinion, rapide. \\
\hline
an \eng{opinion poll} & un sondage d'opinion & Expression figée. Attention : \eng{poll} $\neq$ \eng{pole} (poteau, pôle). \\
\hline
a \eng{ballot} / a \eng{ballot paper} & un bulletin de vote & Contexte électoral. \eng{secret ballot} (vote à bulletin secret). \\
\hline
a \eng{ballot box} & une urne & \\
\hline
a \eng{rating scale} / a \eng{Likert scale} & une échelle d'évaluation & Ex. : \eng{Strongly agree} à \eng{Strongly disagree}. \\
\hline
an \eng{online survey} / a \eng{poll} & un sondage en ligne & \eng{to conduct an online survey}. \\
\hline
\end{tabularx}
\end{center}
\end{propriete}

\begin{exemplebox}[Exemples traduits — Acteurs et outils]
\eng{The pollster designed a detailed questionnaire.} « Le sondeur a conçu un questionnaire détaillé. » \\
\eng{Two thousand respondents filled in the questionnaire online.} « Deux mille répondants ont rempli le questionnaire en ligne. » \\
\eng{The interviewer recorded the interviewee's answers carefully.} « L'enquêteur a noté soigneusement les réponses de l'interviewé. » \\
\eng{The sample was not representative of the whole population.} « L'échantillon n'était pas représentatif de l'ensemble de la population. »
\end{exemplebox}

\courssub{Champ lexical 2 : les actions et les résultats}

\begin{propriete}[Verbes d'action (Actions) — Collocations clés]
\begin{center}
\begin{tabularx}{\linewidth}{|X|X|X|}
\hline
\rowcolor{popblueL} \textbf{Collocation anglaise} & \textbf{Français} & \textbf{Structure / Piège} \\
\hline
\eng{to conduct a survey / a poll} & mener, réaliser une enquête & \textbf{Jamais} \eng{to make a survey} (style soutenu). \\
\hline
\eng{to carry out a survey / a poll} & effectuer, réaliser une enquête & Synonyme de \eng{conduct}. \\
\hline
\eng{to design a questionnaire} & élaborer un questionnaire & \\
\hline
\eng{to select a sample} & sélectionner un échantillon & \eng{a random sample}, \eng{a stratified sample}. \\
\hline
\eng{to interview respondents} & interroger des personnes & \\
\hline
\eng{to question / to ask someone} & questionner quelqu'un & \eng{to ask questions} (pas \eng{to question questions}). \\
\hline
\eng{to take part in} / \eng{to participate in} & participer à & \textbf{Préposition \eng{in} obligatoire}. \\
\hline
\eng{to answer / to respond to} & répondre à & \eng{to respond to a questionnaire}. \\
\hline
\eng{to release / to publish the findings} & publier les résultats & \eng{findings} toujours au pluriel. \\
\hline
\eng{to analyse / to interpret the data} & analyser / interpréter les données & \eng{data} pluriel (sing. \eng{datum}, rare). \\
\hline
\end{tabularx}
\end{center}
\end{propriete}

\begin{propriete}[Noms de résultats et d'analyse (Results) — Vocabulaire clé]
\begin{center}
\begin{tabularx}{\linewidth}{|X|X|X|}
\hline
\rowcolor{popblueL} \textbf{English} & \textbf{Français} & \textbf{Note / Collocation} \\
\hline
the \eng{findings} & les résultats, les conclusions & \textbf{Toujours au pluriel}. \eng{key findings}, \eng{preliminary findings}. \\
\hline
the \eng{results} / the \eng{figures} & les résultats, les chiffres & \eng{the results show that…}, \eng{according to the figures}. \\
\hline
a \eng{percentage} / \eng{per cent} & un pourcentage / pour cent & \eng{60 per cent} (deux mots, pas de « s » à \eng{per cent}). \eng{a high percentage of}. \\
\hline
a \eng{majority} / the \eng{majority of} & une majorité / la majorité de & \eng{the majority of voters are…} (verbe pluriel si nom pluriel après \eng{of}). \\
\hline
a \eng{minority} / the \eng{minority of} & une minorité / la minorité de & \\
\hline
the \eng{margin of error} & la marge d'erreur & \textbf{Jamais} \eng{error margin}. \eng{a margin of error of 3\%}. \\
\hline
a \eng{trend} & une tendance & \eng{an upward/downward trend}. \\
\hline
a \eng{turnout} & un taux de participation & \eng{voter turnout}, \eng{a low turnout}. \\
\hline
\eng{statistics} / \eng{stats} & des statistiques & Pluriel. \eng{official statistics}. \\
\hline
\end{tabularx}
\end{center}
\end{propriete}

\begin{exemplebox}[Exemples traduits — Actions et résultats]
\eng{The institute conducted a nationwide poll last week.} « L'institut a réalisé un sondage national la semaine dernière. » \\
\eng{Respondents were asked to rate the government's performance.} « Les personnes interrogées ont été invitées à noter la performance du gouvernement. » \\
\eng{The findings reveal a clear trend: youth unemployment is rising.} « Les résultats révèlent une tendance claire : le chômage des jeunes augmente. » \\
\eng{Sixty per cent of the sample agreed with the proposal.} « Soixante pour cent de l'échantillon étaient d'accord avec la proposition. » \\
\eng{The margin of error stands at plus or minus two percentage points.} « La marge d'erreur s'élève à plus ou moins deux points de pourcentage. »
\end{exemplebox}

\courssub{Faux amis et pièges des francophones}

\begin{attention}[Tableau des faux amis et calques fréquents — À mémoriser]
\begin{center}
\begin{tabularx}{\linewidth}{|X|X|X|X|}
\hline
\rowcolor{poporangeL} \textbf{Mot / Expression anglais} & \textbf{Sens réel (FR)} & \textbf{Piège (Faux ami / Calque)} & \textbf{Correction} \\
\hline
\eng{actually} & en fait, en réalité & « Actuellement » & \eng{currently}, \eng{at present} \\
\hline
\eng{a survey} & une enquête, un sondage & « Une survie » & \eng{survival} = survie \\
\hline
\eng{to demand} & exiger, réclamer fermement & « Demander » & \eng{to ask (for)}, \eng{to request} \\
\hline
\eng{interviewed} (adj/part.) & interrogé (enquête/emploi) & « Interrogé » (police) & \eng{interrogated} (contexte policier/judiciaire) \\
\hline
\eng{interrogated} & interrogé (par la police) & Utilisé pour un sondage & \eng{interviewed}, \eng{questioned} \\
\hline
\eng{according to} & selon, d'après & « Suivant / Après » & \eng{following} (temps), \eng{after} (temps) \\
\hline
\eng{in favour of} & en faveur de, pour & « En faveur pour » & \textbf{Jamais} \eng{in favour for} \\
\hline
\eng{opposed to} & opposé à, contre & « Opposé contre » & \textbf{Jamais} \eng{opposed against} \\
\hline
\eng{to ask for} & demander (qch) & « Demander » sans prép. & \eng{to ask} qn \eng{for} qch / \eng{to ask} qn \eng{to} V \\
\hline
\eng{information} (indénombrable) & une information, des informations & « Informations » (avec s) & \eng{pieces of information}, \eng{some information} \\
\hline
\eng{a questionnaire} & un questionnaire & Orthographe : un seul « n » & \textbf{Deux « n »} : \eng{questionn\underline{ai}re} \\
\hline
\eng{poll} vs \eng{pole} & sondage vs poteau/pôle & Homophones & \eng{opinion poll}, \eng{North Pole} \\
\hline
\end{tabularx}
\end{center}
\end{attention}

\begin{exoresolu}[Translation — Beware of false friends and calques]
\textbf{Énoncé} — Translate into English, avoiding the traps.
\begin{enumerate}
\item Actuellement, la majorité des électeurs est favorable à cette nouvelle loi.
\item Les personnes interrogées ont demandé plus d'informations sur le sondage.
\item Selon une enquête récente, un jeune sur trois participe aux élections locales.
\item Le sondeur a exigé que le questionnaire soit modifié.
\item Les résultats montrent que 40\% des répondants sont contre le projet.
\end{enumerate}
\tcblower
\textbf{Solution détaillée}
\begin{enumerate}
\item \eng{Currently} (ou \eng{At present}), \textbf{the majority of voters are in favour of} this new law. (\eng{actually} = en fait ; \eng{in favour of} + \eng{the majority of} + pluriel $\rightarrow$ \eng{are}).
\item \textbf{The respondents asked for more information about the poll.} (\eng{interrogated} = interrogatoire police ; \eng{demanded} = exigé ; \eng{information} indénombrable, pas de \eng{s}).
\item \textbf{According to a recent survey, one young person out of three takes part in local elections.} (\eng{following/after} = temps ; \eng{takes part in} / \eng{participates in} + \eng{in}).
\item \textbf{The pollster demanded that the questionnaire be modified.} (Ici \eng{demanded} = exigé, correct. Subjonctif \eng{be modified} après \eng{demand that}).
\item \textbf{The findings show that 40 per cent of the respondents are opposed to the project.} (\eng{results/findings} ; \eng{per cent} sans \eng{s} ; \eng{opposed to}, pas \eng{against}).
\end{enumerate}
\textbf{Conclusion} : Cinq pièges évités (\eng{currently}, \eng{respondents/asked for/information}, \eng{according to/takes part in}, \eng{demanded} utilisé correctement, \eng{findings/per cent/opposed to}).
\end{exoresolu}

\courssub{Familles de mots, prononciation et orthographe}

\begin{methode}[Comment construire les familles de mots (Word Formation)]
\begin{enumerate}
\item \textbf{Identifie la racine} (souvent le verbe) : \eng{respond}, \eng{interview}, \eng{elect}, \eng{represent}, \eng{analyse}.
\item \textbf{Applique les suffixes productifs} :
\begin{itemize}
\item Personne (agent) : \emph{-er}, \emph{-or}, \emph{-ist}, \emph{-ee} (celui qui subit). Ex. \eng{interviewer} / \eng{interviewee}, \eng{pollster} (irrégulier).
\item Action / Nom abstrait : \emph{-tion/-sion}, \emph{-ment}, \emph{-ance/-ence}, \emph{-al}. Ex. \eng{election}, \eng{analysis}, \eng{response}, \eng{approval}.
\item Adjectif : \emph{-ive}, \emph{-able}, \emph{-ent/-ant}, \emph{-al}. Ex. \eng{representative}, \eng{responsive}, \eng{analytical}.
\end{itemize}
\item \textbf{Surveille les changements orthographiques} : \eng{major $\rightarrow$ majority}, \eng{minor $\rightarrow$ minority}, \eng{popular $\rightarrow$ popularity}, \eng{refer $\rightarrow$ referendum} (pluriel \eng{referendums} ou \eng{referenda}).
\item \textbf{Note la prononciation} (lettres françaises) : \eng{questionnaire} « kès-tio-nèr » ; \eng{percentage} « per-sen-tidj » ; \eng{majority} « ma-djo-ri-ti » ; \eng{respondent} « ri-spon-dent » ; \eng{analyst} « a-na-liste ».
\end{enumerate}
\end{methode}

\begin{propriete}[Tableau de formation des mots — Familles du thème « Polls/Surveys »]
\begin{center}
\begin{tabular}{|l|l|l|l|}
\hline
\rowcolor{popblueL} \textbf{Verbe} & \textbf{Personne (Agent)} & \textbf{Action / Nom abstrait} & \textbf{Adjectif} \\
\hline
to \eng{vote} & a \eng{voter} & the \eng{vote}, the \eng{voting} & \eng{voting} (adj., ex. \eng{voting age}) \\
\hline
to \eng{interview} & an \eng{interviewer} / an \eng{interviewee} & an \eng{interview} & \eng{interview} (adj., ex. \eng{interview room}) \\
\hline
to \eng{elect} & an \eng{elector} & an \eng{election} & \eng{elected} (part. passé adj.), \eng{electoral} \\
\hline
to \eng{respond} & a \eng{respondent} & a \eng{response} & \eng{responsive}, \eng{responsible} (attention sens diff.) \\
\hline
to \eng{represent} & a \eng{representative} & \eng{representation} & \eng{representative} \\
\hline
to \eng{analyse} (UK) / \eng{analyze} (US) & an \eng{analyst} & \eng{analysis} (pl. \eng{analyses}) & \eng{analytical} \\
\hline
to \eng{participate} & a \eng{participant} & \eng{participation} & \eng{participatory} \\
\hline
to \eng{question} & a \eng{questioner} & a \eng{question}, a \eng{questionnaire} & \eng{questionable} \\
\hline
to \eng{poll} & a \eng{pollster} & a \eng{poll} & \eng{polling} (adj., ex. \eng{polling station}) \\
\hline
\end{tabular}
\end{center}
\end{propriete}

\begin{exoresolu}[Word formation — Complete the table and the sentences]
\textbf{Énoncé} — a) Complete the table with the missing forms (use British spelling).

\begin{center}
\begin{tabular}{|l|l|l|l|}
\hline
\rowcolor{popblueL} Verb & Person & Action / Abstract Noun & Adjective \\
\hline
to \eng{question} & \ldots\ldots & \ldots\ldots & \ldots\ldots \\
\hline
to \eng{analyse} & \ldots\ldots & \ldots\ldots & \ldots\ldots \\
\hline
to \eng{represent} & \ldots\ldots & \ldots\ldots & \ldots\ldots \\
\hline
\end{tabular}
\end{center}

b) Complete the sentences with the correct form of the word in brackets.
\begin{enumerate}
\item The \ldots\ldots (elect) of the new mayor was confirmed by an opinion \ldots\ldots (poll/pole).
\item We need a \ldots\ldots (represent) sample of the population for this \ldots\ldots (survey).
\item The data \ldots\ldots (analyse) by the expert revealed a surprising \ldots\ldots (tend).
\item Many \ldots\ldots (respond) complained about the length of the \ldots\ldots (question).
\end{enumerate}
\tcblower
\textbf{Solution détaillée}
\textbf{a)}
\begin{itemize}
\item \eng{to question} $\rightarrow$ Person: \eng{questioner} ; Action: \eng{question}, \eng{questionnaire} ; Adj: \eng{questionable}.
\item \eng{to analyse} $\rightarrow$ Person: \eng{analyst} ; Action: \eng{analysis} (pl. \eng{analyses}) ; Adj: \eng{analytical}.
\item \eng{to represent} $\rightarrow$ Person: \eng{representative} ; Action: \eng{representation} ; Adj: \eng{representative}.
\end{itemize}

\textbf{b)}
\begin{enumerate}
\item The \textbf{election} of the new mayor was confirmed by an opinion \textbf{poll}. (\eng{election} = nom action ; \eng{poll} = sondage, homophone de \eng{pole}).
\item We need a \textbf{representative} sample of the population for this \textbf{survey}. (Adj \eng{representative} ; nom \eng{survey}).
\item The data \textbf{was analysed} (ou \eng{were analysed}) by the expert \textbf{analyst} revealed a surprising \textbf{tendency} (ou \eng{trend}). (\eng{data} + verbe sing/plur ; \eng{analyst} nom personne ; \eng{tendency} nom, \eng{trend} plus courant pour sondage).
\item Many \textbf{respondents} complained about the length of the \textbf{questionnaire}. (Pluriel \eng{respondents} ; \eng{questionnaire} nom outil, pas \eng{question}).
\end{enumerate}
\textbf{Conclusion} : Maîtriser les suffixes (\emph{-er/-or}, \emph{-ee}, \emph{-tion/-sis}, \emph{-ive/-al}) et l'orthographe des homophones (\eng{poll/pole}) et mots complexes (\eng{questionnaire}, \eng{analysis}) est essentiel pour le Writing et le Use of English.
\end{exoresolu}

\courssub{Emploi en contexte : lire, dire et écrire un sondage}

\begin{methode}[Comment présenter oralement les résultats d'un sondage (Speaking)]
Pour décrire un sondage à l'oral (exposé, débat, interview), structure ton propos en 5 étapes :
\begin{enumerate}
\item \textbf{Introduis la source et l'échantillon} : \eng{According to a recent survey on youth unemployment carried out in Douala…} / \eng{A poll conducted last month among 1,000 voters shows that…}
\item \textbf{Donne les chiffres avec structures variées} : \eng{Sixty per cent of the respondents said that…} / \eng{One out of three students admitted that…} / \eng{The majority of voters are in favour of…} / \eng{A significant minority (15\%) opposed the measure.}
\item \textbf{Compare et nuance} : \eng{while}, \eng{whereas}, \eng{however}, \eng{on the other hand}, \eng{in contrast}. Ex. \eng{While 70\% of adults support the project, only 40\% of teenagers do.}
\item \textbf{Interprète} : \eng{This suggests that…} / \eng{These figures reveal that…} / \eng{It appears that…} / \eng{This could explain why…}
\item \textbf{Conclus} : \eng{To sum up, the survey clearly shows that…} / \eng{In conclusion, the data indicates…}
\item \textbf{Soigne la prononciation} : \eng{survey} (nom) « seur-vèï » (accent 1ère syll.) ; \eng{percentage} « per-sen-tidj » ; \eng{majority} « ma-djo-ri-ti » ; \eng{respondent} « ri-spon-dent » ; \eng{questionnaire} « kès-tio-nèr ».
\end{enumerate}
\end{methode}

\begin{exoresolu}[Speaking — Report the results of a class survey]
\textbf{Énoncé} — Your English club carried out a survey on mobile phone use among 100 students of your school. Results: 85 own a smartphone; 60 use it mainly for social media; 25 use it mainly for studies; 15 think phones should be banned in class. Report these findings orally in 5–6 sentences.
\tcblower
\textbf{Solution détaillée — Raisonnement} : (1) Source + échantillon. (2) Chiffres variés (\eng{per cent}, \eng{one out of four}, \eng{the majority}). (3) Comparaison (\eng{whereas}). (4) Interprétation. (5) Conclusion. Varier les verbes : \eng{declared}, \eng{admitted}, \eng{believed}.
\textbf{Production modèle} :
« Our English club recently \textbf{conducted a survey} on mobile phone use among \textbf{one hundred students} of our school. \textbf{According to the findings}, \textbf{eighty-five per cent} of the \textbf{respondents} \textbf{own} a smartphone. \textbf{The majority of them}, that is to say sixty students, \textbf{declared} that they use it mainly for social media, \textbf{whereas} only \textbf{one out of four} \textbf{uses} it mainly for studies. \textbf{Interestingly}, fifteen per cent of the students \textbf{believe} that phones should be \textbf{banned in class}. \textbf{These figures suggest that} mobile phones are mostly used for entertainment rather than for learning. \textbf{To sum up}, the survey shows that our school should teach students how to use their phones more responsibly. »
\textbf{Points vérifiés} : \eng{eighty-five per cent} (deux mots, pas de \eng{s}) ; \eng{one out of four} + \eng{uses} (singulier) ; \eng{banned in class} (préposition \eng{in}) ; \eng{rather than} (opposition) ; \eng{information} indénombrable (non utilisé ici mais rappel).
\end{exoresolu}

\begin{methode}[Comment rédiger un rapport écrit sur un sondage (Writing)]
\begin{enumerate}
\item \textbf{Phrase d'introduction (Topic sentence)} : Source + Sujet + Échantillon + Date/Lieu. \eng{A survey on reading habits was recently conducted among three hundred students in Bamenda.}
\item \textbf{2 à 3 phrases de résultats} : Varier les structures (\eng{per cent}, \eng{the majority of}, \eng{one out of five}, \eng{a minority of}) et les verbes de rapport (\eng{say, declare, admit, claim, believe, think}). Utiliser \eng{while / whereas} pour contraster.
\item \textbf{Une phrase d'interprétation} : \eng{These findings suggest that…} / \eng{This indicates a shift in…}
\item \textbf{Une phrase de conclusion ou recommandation} : \eng{Therefore, schools should…} / \eng{It is recommended that…}
\item \textbf{Relis (Check-list)} : Accords (\eng{the majority of students are}), Prépositions (\eng{survey on}, \eng{in favour of}, \eng{opposed to}, \eng{take part in}), Orthographe (\eng{questionnaire}, \eng{percentage}, \eng{respondent}), Indénombrables (\eng{information}, \eng{data}), Pas de calque (\eng{actually}, \eng{according to}).
\end{enumerate}
\end{methode}

\begin{exoresolu}[Writing — A short report (80–100 words)]
\textbf{Énoncé} — Write a short report on this survey: Topic: reading habits; Sample: 300 students in Bamenda; Results: 40\% read books every week; 35\% prefer reading on phones; 25\% never read outside school; Most say reading improves English.
\tcblower
\textbf{Solution détaillée — Raisonnement} : Plan en 5 étapes. Vérifier : \eng{was conducted} (passif), \eng{forty per cent}, \eng{thirty-five per cent}, \eng{one student out of four} / \eng{one out of four students}, \eng{the majority of students believe} (pluriel), \eng{improves} (gérondif sujet \eng{reading} $\rightarrow$ 3e pers. sg.), \eng{by creating} (moyen).
\textbf{Production modèle} :
« A survey on reading habits was recently conducted among three hundred students in Bamenda. According to the findings, forty per cent of the respondents read books every week, while thirty-five per cent declared that they prefer reading on their phones. Surprisingly, one student out of four admitted that he or she never reads outside school. However, the majority of the students believe that reading improves their English. These findings suggest that young people still enjoy reading, but that they are changing their habits. Therefore, schools should encourage digital reading by creating online libraries accessible on mobile phones. »
\textbf{Vérifications finales} : $\approx$ 95 mots. \eng{was conducted} (collocation) ; \eng{one student out of four} ; \eng{believe that reading improves} ; \eng{by creating}. Aucun calque.
\end{exoresolu}

\courssub{Activités de consolidation}

\begin{experience}[Jeu de rôle : « Le sondeur et le répondant »]
En groupes de 3 : 1 \eng{pollster} (note les réponses), 1 \eng{interviewer} (pose les questions), 1 \eng{respondent} (répond).
\textbf{Sujet} : « Les lycéens et l'orientation professionnelle ».
\textbf{Grille questions} (préparée par le \eng{pollster}) : 5 questions (choix multiples + une ouverte).
\textbf{Déroulement} : L'\eng{interviewer} interroge le \eng{respondent} (3 min). Le \eng{pollster} remplit la fiche. Rotation des rôles.
\textbf{Production orale finale} : Chaque \eng{pollster} présente les \eng{findings} de son « échantillon » de 2 personnes à la classe : \eng{“According to our mini-poll, 50\% of respondents want to study medicine, whereas 50\% prefer engineering. One respondent declared that…”}
\end{experience}

\begin{exercice}[Entraînement individuel — Vocabulary \& Grammar]
\textbf{Exercice 1 — Choisissez le bon mot.}
\begin{enumerate}
\item The institute \_\_\_\_\_\_ a nationwide poll last month. (conducted / made / did)
\item \_\_\_\_\_\_ to the survey, unemployment is the main worry. (According / Following / After)
\item 60\% of the respondents \_\_\_\_\_\_ in favour of the reform. (is / are / be)
\item The \_\_\_\_\_\_ of error is 3\%. (margin / error / mistake)
\item Please fill \_\_\_\_\_\_ this questionnaire. (in / on / up)
\item The \_\_\_\_\_\_ were interviewed by phone. (respondents / interrogated / askers)
\item He \_\_\_\_\_\_ more details about the poll. (asked for / demanded / required)
\item One student \_\_\_\_\_\_ four uses the library daily. (out of / on / for)
\end{enumerate}

\textbf{Exercice 2 — Complétez avec la forme correcte du mot entre parenthèses.}
\begin{enumerate}
\item The \_\_\_\_\_\_ (analyse) of the \_\_\_\_\_\_ (find) will take two weeks.
\item We need a \_\_\_\_\_\_ (represent) \_\_\_\_\_\_ (sample) for the \_\_\_\_\_\_ (survey).
\item The \_\_\_\_\_\_ (elect) results will be announced tomorrow.
\item Many young people are \_\_\_\_\_\_ (interest) \_\_\_\_\_\_ politics.
\item The \_\_\_\_\_\_ (question) was too difficult for the \_\_\_\_\_\_ (respond).
\end{enumerate}

\textbf{Exercice 3 — Traduisez en anglais.}
\begin{enumerate}
\item Selon un sondage récent, la majorité des Camerounais sont optimistes.
\item Les personnes interrogées ont exigé de meilleurs services publics.
\item Actuellement, un jeune sur cinq participe à des activités bénévoles.
\item La marge d'erreur de cette enquête est de 2\%.
\end{enumerate}
\end{exercice}

\begin{corrige}[Exercice 1, 2, 3]
\textbf{Exercice 1}
\begin{enumerate}
\item \textbf{conducted} (collocation \eng{conduct a poll}).
\item \textbf{According} (\eng{According to} = selon).
\item \textbf{are} (\eng{the majority of} + pluriel $\rightarrow$ verbe pluriel).
\item \textbf{margin} (\eng{margin of error} expression figée).
\item \textbf{in} (\eng{fill in a questionnaire} / \eng{fill out} US).
\item \textbf{respondents} (terme technique, pas \eng{interrogated}).
\item \textbf{asked for} (\eng{demand} = exiger ; \eng{ask for} = demander qch).
\item \textbf{out of} (\eng{one out of four} = un sur quatre).
\end{enumerate}

\textbf{Exercice 2}
\begin{enumerate}
\item \textbf{analysis} (nom, accent sur 2e syll.) / \textbf{findings} (toujours pluriel).
\item \textbf{representative} (adj.) / \textbf{sample} (nom) / \textbf{survey} (nom).
\item \textbf{election} (nom action) ou \textbf{electoral} (adj. si \eng{results} suit, mais ici \eng{The election results} ou \eng{The electoral results}). Attendu : \textbf{election} (nom composé) ou \textbf{electoral}. \eng{The election results} est plus courant.
\item \textbf{interested in} (préposition \eng{in} obligatoire).
\item \textbf{questionnaire} (nom outil, 2 \eng{n}) / \textbf{respondents} (pluriel personnes).
\end{enumerate}

\textbf{Exercice 3}
\begin{enumerate}
\item \eng{According to a recent poll, the majority of Cameroonians are optimistic.}
\item \eng{The respondents demanded better public services.} (ou \eng{asked for} si moins fort, mais \eng{exigé} = \eng{demanded}).
\item \eng{Currently, one young person out of five takes part in voluntary activities.} (\eng{currently} pas \eng{actually} ; \eng{takes part in}).
\item \eng{The margin of error of this survey is 2\%.} (\eng{margin of error} ; \eng{survey} ou \eng{poll}).
\end{enumerate}
\end{corrige}

\courssub{Pour réviser : les points clés}

\begin{aretenir}[Synthèse — Vocabulary: Polls and Surveys]
\textbf{1. Acteurs :} \eng{pollster}, \eng{respondent}, \eng{interviewer}, \eng{interviewee}, \eng{sample} (échantillon), \eng{target population}.
\textbf{2. Outils :} \eng{questionnaire} (2 \eng{n}), \eng{survey}, \eng{poll}, \eng{opinion poll}, \eng{ballot box}, \eng{rating scale}.
\textbf{3. Actions (Collocations) :} \eng{conduct/carry out a survey}, \eng{design a questionnaire}, \eng{select a sample}, \eng{interview respondents}, \eng{take part in/participate in}, \eng{release/publish findings}, \eng{analyse data}.
\textbf{4. Résultats :} \eng{findings} (pl.), \eng{figures}, \eng{percentage/per cent} (pas de \eng{s}), \eng{majority/minority of} + verbe accordé au complément, \eng{margin of error}, \eng{trend}, \eng{turnout}.
\textbf{5. Prépositions verrouillées :} \eng{according to}, \eng{in favour of}, \eng{opposed to}, \eng{interested in}, \eng{survey on}, \eng{ask for}, \eng{take part in}.
\textbf{6. Faux amis :} \eng{actually} (en fait) $\neq$ \eng{currently} (actuellement) ; \eng{demand} (exiger) $\neq$ \eng{ask} (demander) ; \eng{interviewed} (enquête) $\neq$ \eng{interrogated} (police) ; \eng{information} (indénombrable).
\textbf{7. Orthographe :} \eng{questionnaire}, \eng{percentage}, \eng{respondent}, \eng{analysis/analyses}, \eng{poll/pole}.
\textbf{8. Prononciation clé :} \eng{survey} « seur-vèï », \eng{questionnaire} « kès-tio-nèr », \eng{majority} « ma-djo-ri-ti », \eng{percentage} « per-sen-tidj ».
\end{aretenir}

\begin{savaistu}[Culture \& Contexte : Les sondages au Cameroun et dans le monde]
Au Cameroun, l'\eng{National Institute of Statistics (INS)} et des instituts privés (ex. \eng{GeoPoll}, \eng{Afrobarometer}) réalisent des enquêtes sur l'emploi, la santé, la gouvernance. La \eng{margin of error} y est cruciale car les \eng{samples} sont parfois difficiles à rendre \eng{representative} (zones rurales inaccessibles, barrières linguistiques).
Dans les pays anglophones (\eng{UK}, \eng{USA}), les \eng{opinion polls} rythment la vie politique (\eng{Gallup}, \eng{YouGov}, \eng{Ipsos MORI}). Le \eng{exit poll} (sondage à la sortie des urnes) prédit les résultats le soir de l'élection. Le terme \eng{pollster} est très courant dans la presse (ex. \eng{“Pollsters predict a tight race”}). Attention : en anglais britannique, on écrit \eng{favour}, \eng{analyse} ; en américain \eng{favor}, \eng{analyze}. Au Cameroun, l'orthographe britannique est la norme scolaire.
\end{savaistu}

\begin{popfigure}
\begin{tikzpicture}[
  mindmap,
  concept color=popblue,
  level 1 concept/.append style={level distance=3.5cm, sibling angle=90, concept color=popblue!70},
  level 2 concept/.append style={level distance=3cm, sibling angle=45, concept color=poporange!70},
  level 3 concept/.append style={level distance=2.5cm, sibling angle=30, concept color=popgreen!70},
  every node/.style={font=\small},
  text=popdark
]
\node[concept] {Polls \& Surveys}
  child[concept color=popblue] {node[concept] {Actors}
    child {node[concept] {pollster}}
    child {node[concept] {respondent}}
    child {node[concept] {interviewer / interviewee}}
    child {node[concept] {sample / population}}
  }
  child[concept color=poporange] {node[concept] {Tools}
    child {node[concept] {questionnaire}}
    child {node[concept] {survey / poll}}
    child {node[concept] {ballot box}}
    child {node[concept] {rating scale}}
  }
  child[concept color=popgreen] {node[concept] {Actions}
    child {node[concept] {conduct / carry out}}
    child {node[concept] {design / select}}
    child {node[concept] {interview / question}}
    child {node[concept] {take part in}}
    child {node[concept] {release findings}}
  }
  child[concept color=poppurple] {node[concept] {Results}
    child {node[concept] {findings / figures}}
    child {node[concept] {percentage / per cent}}
    child {node[concept] {majority / minority}}
    child {node[concept] {margin of error}}
    child {node[concept] {trend / turnout}}
  };
\end{tikzpicture}
\legende{Carte mentale récapitulative du lexique des sondages (Lesson 45).}
\end{popfigure}

\newpage

\coursec{Exercices}

\courssub{Je vérifie mes connaissances}

\begin{exercice}[QCM : vocabulaire de base]
Chaque phrase ne comporte qu'une seule réponse correcte. Choisissez le mot qui convient.
\begin{enumerate}
\item The \eng{\_\_\_\_\_\_\_\_\_\_} interviewed 1,200 voters across the country last week.
\begin{itemize}
\item[a)] pollster \item[b)] respondent \item[c)] census \item[d)] ballot
\end{itemize}
\item A \eng{\_\_\_\_\_\_\_\_\_\_} of 5\% means the true value could be 5 points higher or lower.
\begin{itemize}
\item[a)] trend \item[b)] margin of error \item[c)] turnout \item[d)] panel
\end{itemize}
\item The \eng{\_\_\_\_\_\_\_\_\_\_} was designed to be representative of the whole population.
\begin{itemize}
\item[a)] questionnaire \item[b)] sample \item[c)] finding \item[d)] statistic
\end{itemize}
\item Many young people remain \eng{\_\_\_\_\_\_\_\_\_\_} about which candidate to support.
\begin{itemize}
\item[a)] undecided \item[b)] biased \item[c)] random \item[d)] eligible
\end{itemize}
\item The \eng{\_\_\_\_\_\_\_\_\_\_} showed a clear \eng{trend} towards the opposition party.
\begin{itemize}
\item[a)] census \item[b)] poll \item[c)] ballot \item[d)] panel
\end{itemize}
\end{enumerate}
\end{exercice}

\begin{exercice}[Vrai ou Faux ? Justifiez en anglais]
Lisez les affirmations ci-dessous. Écrivez \eng{True} ou \eng{False} et justifiez votre réponse par une phrase complète en anglais en reprenant le vocabulaire de la leçon.
\begin{enumerate}
\item A \eng{census} questions every single member of a population, while a \eng{poll} questions only a \eng{sample}.
\item The \eng{margin of error} increases when the \eng{sample size} increases.
\item A \eng{respondent} is the person who conducts the interview.
\item \eng{Undecided} voters have already made their choice but refuse to reveal it.
\item A \eng{representative sample} accurately reflects the characteristics of the target population.
\end{enumerate}
\end{exercice}

\begin{exercice}[Vocabulaire : familles de mots et collocations]
Complétez le tableau suivant en donnant la forme manquante (nom, verbe, adjectif) et une collocation fréquente (association de mots naturelle). Le premier est fait en exemple.
\begin{center}
\begin{tabularx}{\linewidth}{|X|X|X|X|}
\hline
\rowcolor{popblueL}
\textbf{Verbe} & \textbf{Nom} & \textbf{Adjectif} & \textbf{Collocation} \\
\hline
to poll & a poll / polling & — & \eng{to conduct a poll} \\
\hline
to respond & & & \\
\hline
& representation & & \\
\hline
to sample & & & \\
\hline
& & random & \\
\hline
to predict & & & \\
\hline
\end{tabularx}
\end{center}
\end{exercice}

\begin{exercice}[Conjugaison et forme verbale : le langage des résultats]
Mettez les verbes entre parenthèses à la forme correcte (temps, voix, participe) pour commenter des résultats de sondage.
\begin{enumerate}
\item The latest poll \eng{\_\_\_\_\_\_\_\_\_\_} (show) that the incumbent \eng{\_\_\_\_\_\_\_\_\_\_} (lead) by 4 points.
\item The data \eng{\_\_\_\_\_\_\_\_\_\_} (collect) over a period of three days last month.
\item If the sample \eng{\_\_\_\_\_\_\_\_\_\_} (be) larger, the margin of error \eng{\_\_\_\_\_\_\_\_\_\_} (be) smaller.
\item The pollster recommended that the questionnaire \eng{\_\_\_\_\_\_\_\_\_\_} (revise) before the next wave.
\item By the time the results \eng{\_\_\_\_\_\_\_\_\_\_} (publish), the campaign \eng{\_\_\_\_\_\_\_\_\_\_} (already / end).
\end{enumerate}
\end{exercice}

\courssub{Je m'entraîne}

\begin{exercice}[Traduction : du français vers l'anglais]
Traduisez les phrases suivantes en anglais en utilisant le vocabulaire précis de la leçon (poll, sample, respondent, margin of error, undecided, turnout, findings, to conduct).
\begin{enumerate}
\item L'institut de sondage a réalisé une enquête auprès d'un échantillon représentatif de 2 000 électeurs.
\item La marge d'erreur est de 3 points de pourcentage.
\item Selon les résultats, 15 \% des personnes interrogées sont encore indécises.
\item Le taux de participation attendu est élevé pour cette élection.
\item Les sondeurs ont interrogé un panel de jeunes à Yaoundé et Douala.
\end{enumerate}
\end{exercice}

\begin{exercice}[Transformation : style direct / rapporté et reformulation]
Transformez les phrases suivantes en utilisant la structure indiquée entre parenthèses. Gardez le sens exact.
\begin{enumerate}
\item “Our sample is perfectly representative,” said the pollster. (Reported speech with \eng{that})
\item The margin of error is 4\%. The lead is not significant. (Combine with \eng{so ... that})
\item They interviewed 500 respondents. (Passive voice)
\item The findings indicate a clear trend. (Start with \eng{According to the findings, ...})
\item “Did you vote in the last election?” the interviewer asked. (Reported question with \eng{whether})
\end{enumerate}
\end{exercice}

\begin{exercice}[Texte à trous : article de presse]
Complétez le texte suivant avec les mots de la liste ci-dessous. Chaque mot s'utilise une seule fois. Attention à la forme (pluriel, adjectif, etc.).
\begin{center}
\textbf{Liste :} \eng{pollster, sample, respondent, questionnaire, representative, undecided, margin of error, findings, trend, turnout}
\end{center}

\begin{textebox}[Election Poll Shows Tight Race]
A new opinion \eng{\_\_\_\_\_\_\_\_\_\_} published yesterday by the National Institute of Statistics reveals a very tight race between the two main candidates. The \eng{\_\_\_\_\_\_\_\_\_\_} interviewed 1,500 eligible voters between March 10 and 15. The \eng{\_\_\_\_\_\_\_\_\_\_} was carefully selected to be \eng{\_\_\_\_\_\_\_\_\_\_} of the national voting population in terms of age, gender and region.

The \eng{\_\_\_\_\_\_\_\_\_\_} show that Candidate A leads with 48\% against 46\% for Candidate B. However, the \eng{\_\_\_\_\_\_\_\_\_\_} is 3\%, which means the lead is not statistically significant. A notable \eng{\_\_\_\_\_\_\_\_\_\_} is the high number of \eng{\_\_\_\_\_\_\_\_\_\_} voters, reaching 12\%. The \eng{\_\_\_\_\_\_\_\_\_\_} rate in the last election was 65\%, but analysts predict a higher \eng{\_\_\_\_\_\_\_\_\_\_} this year due to youth mobilisation.
\end{textebox}
\end{exercice}

\begin{exercice}[Appariement : mots et définitions]
Associez chaque terme de la colonne A à sa définition dans la colonne B. Écrivez les paires (ex. 1–e).
\begin{center}
\begin{tabularx}{\linewidth}{|X|X|}
\hline
\rowcolor{popblueL}
\textbf{A. Terme} & \textbf{B. Définition} \\
\hline
1. Pollster & a. The total number or percentage of eligible voters who cast a ballot. \\
2. Census & b. A person who answers questions in a survey. \\
3. Sample & c. A subset of a population selected for study. \\
4. Respondent & d. An organization or person that conducts polls. \\
5. Panel & e. An official count of every member of a population. \\
6. Turnout & f. A group of people who agree to participate in multiple surveys over time. \\
\hline
\end{tabularx}
\end{center}
\end{exercice}

\courssub{Je me prépare au Bac}

\begin{exercice}[Compréhension écrite : type Baccalauréat]
Lisez attentivement le texte ci-dessous, puis répondez aux questions qui suivent en anglais. Pour les questions Vrai/Faux, citez la phrase du texte qui justifie votre réponse.

\begin{textebox}[Youth Participation Surges in Cameroon’s Municipal Polls — Survey]
A recent survey conducted by the Cameroon Institute for Democratic Studies (CIDS) highlights a significant shift in youth electoral engagement ahead of the February municipal elections. The poll, which sampled 2,400 registered voters aged 18–35 across the ten regions, was carried out between January 5 and January 20. The sample was stratified to ensure it was representative of the urban and rural youth population.

According to the findings, 68\% of respondents declared they would “certainly vote,” compared to only 42\% in the previous municipal election cycle. This 26-point increase suggests a growing political awareness among young Cameroonians. The margin of error for the survey is ±2.1\%.

The questionnaire also explored the main motivations for voting. The top three issues cited were unemployment (34\%), access to quality education (22\%), and infrastructure development (18\%). Interestingly, 15\% of the sample remained undecided just two weeks before the vote, indicating that campaign messages could still sway a substantial bloc.

Dr. Amina Njoya, lead pollster at CIDS, commented: “This trend is encouraging. However, high declared intention does not always translate into actual turnout. Logistical barriers, such as distance to polling stations and voter card collection delays, often reduce the final figure.”

The survey also measured trust in institutions. Only 28\% of respondents said they trusted the electoral commission “a lot” or “somewhat,” while 55\% expressed little or no trust. This scepticism, analysts warn, could depress turnout if not addressed by civic education campaigns.
\end{textebox}

\begin{enumerate}
\item \textbf{True or False?} Justify by quoting the text.
\begin{itemize}
\item[a)] The survey interviewed all registered young voters in Cameroon.
\item[b)] The percentage of youth declaring they will certainly vote has increased since the last municipal elections.
\item[c)] Unemployment is the least cited motivation for voting.
\item[d)] The majority of respondents trust the electoral commission.
\end{itemize}
\item \textbf{Multiple Choice.} Choose the best answer.
\begin{itemize}
\item[1.] The phrase “stratified to ensure it was representative” (line 5) means the sample:
\begin{itemize}
\item[A] was chosen randomly without any criteria.
\item[B] was divided into subgroups reflecting the population structure.
\item[C] only included urban voters.
\item[D] was too small to be reliable.
\end{itemize}
\item[2.] The “26-point increase” (line 9) refers to:
\begin{itemize}
\item[A] the margin of error.
\item[B] the difference in declared voting intention between the two elections.
\item[C] the number of undecided voters.
\item[D] the trust level in the electoral commission.
\end{itemize}
\end{itemize}
\item \textbf{Answer in your own words (based on the text).}
\begin{itemize}
\item[a)] What two logistical barriers does Dr. Njoya mention that could lower actual turnout?
\item[b)] What percentage of respondents expressed little or no trust in the electoral commission?
\item[c)] According to the text, why could the 15\% of undecided voters be important?
\end{itemize}
\item \textbf{Vocabulary in context.} Find in the text words or expressions that match these definitions:
\begin{itemize}
\item[a)] A person who answers questions in a survey (para 2).
\item[b)] The range within which the true value likely lies (para 3).
\item[c)] A general direction in which something is developing (para 5).
\item[d)] To make something lower or weaker (para 6).
\end{itemize}
\end{enumerate}
\end{exercice}

\begin{exercice}[Traduction : passage d'un article vers le français]
Traduisez en français le passage suivant extrait du texte ci-dessus (lignes 1 à 10 : “A recent survey … ±2.1%.”). Votre traduction doit être naturelle et précise, en respectant le registre journalistique.
\begin{textebox}[Passage to translate]
A recent survey conducted by the Cameroon Institute for Democratic Studies (CIDS) highlights a significant shift in youth electoral engagement ahead of the February municipal elections. The poll, which sampled 2,400 registered voters aged 18–35 across the ten regions, was carried out between January 5 and January 20. The sample was stratified to ensure it was representative of the urban and rural youth population.

According to the findings, 68\% of respondents declared they would “certainly vote,” compared to only 42\% in the previous municipal election cycle. This 26-point increase suggests a growing political awareness among young Cameroonians. The margin of error for the survey is ±2.1\%.
\end{textebox}
\end{exercice}

\begin{exercice}[Expression écrite : sujet type Bac]
\textbf{Sujet :} \eng{“The results of a survey on young voters in Cameroon”}

Rédigez un paragraphe cohérent de \textbf{80 à 100 mots} en anglais. Vous imaginerez les résultats d'un sondage fictif sur l'engagement électoral des jeunes au Cameroun (intentions de vote, motivations, obstacles, confiance). Votre production doit impérativement inclure \textbf{au moins 6 mots ou expressions} du lexique étudié dans cette leçon (soulignez-les dans votre copie).

\textit{Indications :} Utilisez des structures pour commenter des données (\eng{According to the poll, ... ; The findings show that ... ; X\% of respondents ... ; The margin of error is ... ; A significant trend is ...}). Soignez la liaison entre les phrases (connecteurs logiques).
\end{exercice}

\newpage

\coursec{Corrigés des exercices}

\begin{corrige}[Exercice 1 — QCM : vocabulaire de base]
\begin{enumerate}
\item \textbf{a) pollster.} \eng{The pollster interviewed 1,200 voters across the country last week.} Le sujet est l'agent qui mène les interviews : le \eng{pollster} (le sondeur). Un \eng{respondent} est interrogé, il n'interroge pas ; un \eng{census} (recensement) et un \eng{ballot} (bulletin/scrutin) ne peuvent pas « interviewer ».
\item \textbf{b) margin of error.} \eng{A margin of error of 5\% means the true value could be 5 points higher or lower.} C'est la définition exacte de la marge d'erreur. Une \eng{trend} est une tendance, le \eng{turnout} le taux de participation, un \eng{panel} un groupe suivi dans le temps.
\item \textbf{b) sample.} \eng{The sample was designed to be representative of the whole population.} C'est l'échantillon qui doit être représentatif, pas le questionnaire (instrument de questions).
\item \textbf{a) undecided.} \eng{Many young people remain undecided about which candidate to support.} « Indécis » : ils n'ont pas encore choisi. \eng{Biased} = biaisé ; \eng{random} = aléatoire ; \eng{eligible} = éligible.
\item \textbf{b) poll.} \eng{The poll showed a clear trend towards the opposition party.} Un sondage révèle une tendance. Le \eng{census} porte sur toute la population ; le \eng{ballot} est le vote lui-même.
\end{enumerate}
\end{corrige}

\begin{corrige}[Exercice 2 — Vrai ou Faux ? Justifiez en anglais]
\begin{enumerate}
\item \textbf{True.} \eng{A census questions every single member of a population, whereas a poll only surveys a sample of that population.}
\item \textbf{False.} \eng{The margin of error decreases when the sample size increases: the larger the sample, the more accurate the results.}
\item \textbf{False.} \eng{A respondent is the person who answers the questions; the person who conducts the interview is the interviewer or the pollster.}
\item \textbf{False.} \eng{Undecided voters have not made their choice yet; they are still hesitating between the candidates.}
\item \textbf{True.} \eng{A representative sample accurately reflects the characteristics of the target population, for example in terms of age, gender and region.}
\end{enumerate}
\end{corrige}

\begin{corrige}[Exercice 3 — Vocabulaire : familles de mots et collocations]
\begin{center}
\begin{tabularx}{\linewidth}{|X|X|X|X|}
\hline
\rowcolor{popblueL}
\textbf{Verbe} & \textbf{Nom} & \textbf{Adjectif} & \textbf{Collocation} \\
\hline
to poll & a poll / polling & — & \eng{to conduct a poll} \\
\hline
to respond & a response / a respondent & responsive & \eng{to respond to a questionnaire} \\
\hline
to represent & representation & representative & \eng{a representative sample} \\
\hline
to sample & a sample / sampling & sampled & \eng{a random sample} \\
\hline
to randomise & randomness & random & \eng{selected at random} \\
\hline
to predict & a prediction & predictable & \eng{to predict the turnout} \\
\hline
\end{tabularx}
\end{center}
\textbf{Points de vigilance :} le nom de \eng{to respond} est \eng{a response} (une réponse) et non « a respond » ; attention à l'orthographe de \eng{representative} (un seul « e » après « repr ») ; on dit \eng{at random} et non « by random ».
\end{corrige}

\begin{corrige}[Exercice 4 — Conjugaison et forme verbale : le langage des résultats]
\begin{enumerate}
\item \eng{The latest poll \textbf{shows} that the incumbent \textbf{is leading} by 4 points.} Présent simple pour un fait établi par le sondage ; présent continu (ou \eng{leads}) pour une situation en cours.
\item \eng{The data \textbf{were collected} over a period of three days last month.} Passif au prétérit (\eng{were} + participe passé) ; \eng{data} est un pluriel ; \eng{last month} impose le prétérit.
\item \eng{If the sample \textbf{were} larger, the margin of error \textbf{would be} smaller.} Conditionnel de type 2 (hypothèse irréelle) : \eng{if} + prétérit, \eng{would} + base verbale.
\item \eng{The pollster recommended that the questionnaire \textbf{be revised} (ou \eng{should be revised}) before the next wave.} Subjonctif après \eng{recommend that} ; voix passive car le questionnaire subit l'action.
\item \eng{By the time the results \textbf{were published}, the campaign \textbf{had already ended}.} Prétérit passif pour le repère ; past perfect (\eng{had} + participe passé) pour l'action antérieure.
\end{enumerate}
\end{corrige}

\begin{corrige}[Exercice 5 — Traduction : du français vers l'anglais]
\begin{enumerate}
\item \eng{The polling institute conducted a survey on a representative sample of 2,000 voters.} (« réaliser une enquête » = \eng{to conduct / carry out a survey}, jamais « to realize a survey ».)
\item \eng{The margin of error is 3 percentage points.} (On dit \eng{percentage points} pour une marge, non « error margin ».)
\item \eng{According to the findings, 15\% of the respondents are still undecided.} (« personnes interrogées » = \eng{respondents} ; « indécis » = \eng{undecided}, pas « undecisive ».)
\item \eng{The expected turnout is high for this election.} (« taux de participation » = \eng{turnout}, faux ami avec « participation rate » qui reste acceptable mais moins idiomatique.)
\item \eng{The pollsters interviewed a panel of young people in Yaoundé and Douala.} (« sondeurs » = \eng{pollsters} ; prétérit simple car l'action est achevée.)
\end{enumerate}
\end{corrige}

\begin{corrige}[Exercice 6 — Transformation : style direct / rapporté et reformulation]
\begin{enumerate}
\item \eng{The pollster said (that) their sample was perfectly representative.} Recul du temps : présent simple → prétérit simple ; \eng{our} → \eng{their}.
\item \eng{The margin of error is so large that the lead is not significant.} (ou \eng{The margin of error is 4\%, so the lead is not significant.}) Structure \eng{so + adjectif + that}.
\item \eng{500 respondents were interviewed.} Passif au prétérit : sujet patient + \eng{were} + participe passé ; l'agent peut être omis.
\item \eng{According to the findings, there is a clear trend.} \eng{According to} + nom, sans « to » supplémentaire.
\item \eng{The interviewer asked whether I had voted in the last election.} Question rapportée : \eng{whether} + ordre sujet–verbe ; recul du prétérit au past perfect (\eng{had voted}).
\end{enumerate}
\end{corrige}

\begin{corrige}[Exercice 7 — Texte à trous : article de presse]
\begin{enumerate}
\item \eng{poll} — \eng{A new opinion poll published yesterday...}
\item \eng{pollster} — \eng{The pollster interviewed 1,500 eligible voters...}
\item \eng{sample} — \eng{The sample was carefully selected...}
\item \eng{representative} — \eng{...to be representative of the national voting population...}
\item \eng{findings} — \eng{The findings show that Candidate A leads with 48\%...} (pluriel, d'où \eng{show} sans « s ».)
\item \eng{margin of error} — \eng{However, the margin of error is 3\%...}
\item \eng{trend} — \eng{A notable trend is the high number of...}
\item \eng{undecided} — \eng{...the high number of undecided voters, reaching 12\%.}
\item \eng{turnout} — \eng{The turnout rate in the last election was 65\%...}
\item \eng{turnout} — attention : le mot \eng{turnout} figurant deux fois dans le texte, la deuxième occurrence attendue est \eng{turnout} (\eng{a higher turnout this year}). Le mot \eng{questionnaire} de la liste reste donc inutilisé dans cette version ; en copie d'examen, on acceptera \eng{turnout} aux deux emplacements, le contexte (« rate », « higher ») l'imposant.
\end{enumerate}
\textbf{Point de méthode :} toujours vérifier l'accord après le mot choisi (\eng{the findings show}, pas « shows ») et la cohérence du champ lexical (un \eng{pollster} interviewe, un \eng{sample} est sélectionné).
\end{corrige}

\begin{corrige}[Exercice 8 — Appariement : mots et définitions]
\begin{center}
\begin{tabularx}{\linewidth}{|X|X|}
\hline
\rowcolor{popblueL}
\textbf{Paire} & \textbf{Justification} \\
\hline
1–d & The \eng{pollster} is the person or organisation that conducts polls. \\
\hline
2–e & A \eng{census} is an official count of every member of a population. \\
\hline
3–c & A \eng{sample} is a subset of a population selected for study. \\
\hline
4–b & A \eng{respondent} is a person who answers questions in a survey. \\
\hline
5–f & A \eng{panel} is a group of people who participate in multiple surveys over time. \\
\hline
6–a & \eng{Turnout} is the number or percentage of eligible voters who cast a ballot. \\
\hline
\end{tabularx}
\end{center}
\end{corrige}

\begin{corrige}[Exercice 9 — Compréhension écrite : type Baccalauréat]
\textbf{Barème indicatif (sur 20) :} Q1 : 4 × 1,5 = 6 pts (0,5 pour True/False + 1 pour la citation) ; Q2 : 2 × 1 = 2 pts ; Q3 : 3 × 2 = 6 pts ; Q4 : 4 × 1,5 = 6 pts.

\textbf{1. True or False?}
\begin{itemize}
\item[a)] \textbf{False.} \eng{“The poll, which sampled 2,400 registered voters aged 18–35 across the ten regions...”} — seulement un échantillon, pas tous les jeunes électeurs.
\item[b)] \textbf{True.} \eng{“68\% of respondents declared they would ‘certainly vote,’ compared to only 42\% in the previous municipal election cycle.”}
\item[c)] \textbf{False.} \eng{“The top three issues cited were unemployment (34\%), access to quality education (22\%), and infrastructure development (18\%).”} — le chômage est la motivation la \emph{plus} citée.
\item[d)] \textbf{False.} \eng{“Only 28\% of respondents said they trusted the electoral commission ‘a lot’ or ‘somewhat,’ while 55\% expressed little or no trust.”}
\end{itemize}

\textbf{2. Multiple Choice.}
\begin{itemize}
\item[1.] \textbf{B} — \eng{stratified} signifie divisé en sous-groupes (strates) reflétant la structure de la population.
\item[2.] \textbf{B} — il s'agit de l'écart d'intention de vote déclarée entre les deux élections (68\% − 42\% = 26 points).
\end{itemize}

\textbf{3. Answer in your own words.}
\begin{itemize}
\item[a)] \eng{Dr. Njoya mentions the distance to polling stations and delays in collecting voter cards.}
\item[b)] \eng{55\% of respondents expressed little or no trust in the electoral commission.}
\item[c)] \eng{They are important because campaign messages could still sway them, so they may change the final result of the election.}
\end{itemize}

\textbf{4. Vocabulary in context.}
\begin{itemize}
\item[a)] \eng{respondent} \item[b)] \eng{margin of error} \item[c)] \eng{trend} \item[d)] \eng{to depress} (dans \eng{could depress turnout}).
\end{itemize}

\textbf{Points de vigilance :} citer la phrase \emph{exacte} du texte entre guillemets pour les Vrai/Faux ; pour les questions ouvertes, reformuler avec ses propres mots sans recopier de longs passages ; ne pas confondre \eng{turnout} (participation réelle) et intention déclarée de voter.
\end{corrige}

\begin{corrige}[Exercice 10 — Traduction : passage d'un article vers le français]
\textbf{Proposition de traduction :}

« Une enquête récente menée par le Cameroon Institute for Democratic Studies (CIDS) met en lumière une évolution significative de l'engagement électoral des jeunes à l'approche des élections municipales de février. Le sondage, qui a interrogé un échantillon de 2 400 électeurs inscrits âgés de 18 à 35 ans dans les dix régions, a été réalisé entre le 5 et le 20 janvier. L'échantillon a été stratifié afin d'être représentatif de la population jeune urbaine et rurale.

Selon les résultats, 68 \% des personnes interrogées ont déclaré qu'elles iraient “certainement voter”, contre seulement 42 \% lors du précédent cycle d'élections municipales. Cette hausse de 26 points traduit une conscience politique croissante chez les jeunes Camerounais. La marge d'erreur de l'enquête est de ± 2,1 \%. »

\textbf{Barème indicatif (sur 10) :} exactitude du sens (4 pts), lexique précis du sondage (3 pts), qualité du français et registre journalistique (2 pts), orthographe et ponctuation (1 pt).

\textbf{Points de vigilance :} \eng{to conduct a survey} = « mener/réaliser une enquête » (pas « conduire ») ; \eng{respondents} = « personnes interrogées / sondés » ; \eng{highlights} = « met en lumière / souligne » ; \eng{ahead of} = « à l'approche de » ; conserver les guillemets autour de « certainly vote ».
\end{corrige}

\begin{corrige}[Exercice 11 — Expression écrite : sujet type Bac]
\textbf{Proposition de rédaction modèle (environ 95 mots) :}

\begin{textebox}[Model answer]
According to a recent \eng{poll} conducted among 1,800 young voters in Yaoundé and Douala, youth electoral engagement is rising. The \eng{findings} show that 62\% of \eng{respondents} intend to vote in the next municipal elections, compared to 45\% five years ago. The main motivations are unemployment and education. However, 14\% of the \eng{sample} remain \eng{undecided}, and only 30\% trust the electoral commission. The \eng{margin of error} is ±2.5\%. \eng{Pollsters} warn that a high declared intention does not guarantee a high \eng{turnout}, as logistical barriers may still discourage young voters on election day.
\end{textebox}

\textbf{Barème indicatif (sur 10) :} respect de la consigne et de la longueur 80–100 mots (2 pts) ; emploi d'au moins 6 éléments du lexique de la leçon, soulignés (3 pts) ; correction grammaticale (2 pts) ; cohérence et connecteurs logiques (2 pts) ; richesse des structures de commentaire de données (1 pt).

\textbf{Points de vigilance :}
\begin{itemize}
\item Employer \eng{according to the poll / the findings show that / X\% of respondents} pour introduire les données.
\item Accorder le verbe avec le sujet : \eng{the findings show}, \eng{62\% of respondents intend}.
\item Éviter les calques : « taux de participation » = \eng{turnout} (pas « participation rate » si possible) ; « indécis » = \eng{undecided} ; « marge d'erreur » = \eng{margin of error} (pas « error margin »).
\item Utiliser des connecteurs : \eng{however, moreover, as a result, in addition}.
\item Ne pas dépasser 100 mots : compter les mots avant de rendre la copie.
\end{itemize}
\end{corrige}

\newpage

\coursec{Fiche bilan}

\begin{popfigure}
\centering
\begin{tikzpicture}[
    mindmap,
    grow cyclic,
    every node/.style={concept, execute at begin node=\hskip0pt},
    concept color=popblue,
    level 1/.append style={level distance=4.5cm, sibling angle=360/7},
    level 2/.append style={level distance=3cm, sibling angle=45},
    actors/.style={concept color=poporange},
    tools/.style={concept color=popgreen},
    verbs/.style={concept color=poppurple},
    results/.style={concept color=poppink},
    idioms/.style={concept color=popgold},
    falsefriends/.style={concept color=popdark},
    families/.style={concept color=popmuted},
    text=white,
    font=\small\bfseries,
    align=center
]
\node[concept, scale=1.2, align=center]{Polls \& Surveys\\Vocabulary}
    child[actors] { node[align=center]{Actors\\Pollster, Respondent,\\Sample, Target population} }
    child[tools] { node[align=center]{Tools \& Process\\Questionnaire, Ballot,\\Margin of error, Bias,\\Random sampling} }
    child[verbs] { node[align=center]{Verbs \& Actions\\Conduct, Carry out,\\Release, Publish,\\Predict, Forecast} }
    child[results] { node[align=center]{Results \& Analysis\\Outcome, Turnout,\\Landslide, Swing,\\Trend, Rating} }
    child[idioms] { node[align=center]{Expressions\\A straw poll,\\A snap poll,\\Neck and neck,\\Too close to call} }
    child[falsefriends] { node[align=center]{False Friends\\Enquête $\neq$ Enquiry\\Sondage $\neq$ Sounding\\Biais $\neq$ Bias} }
    child[families] { node[align=center]{Word Families\\Poll $\rightarrow$ Pollster,\\Survey $\rightarrow$ Surveyor,\\Predict $\rightarrow$ Prediction,\\Respond $\rightarrow$ Respondent} };
\end{tikzpicture}
\legende{Carte mentale : Vocabulaire des sondages et enquêtes (Lesson 45)}
\end{popfigure}

\begin{bilan}

\end{bilan}

\courssub{Lexique clé — English / Français}

\begin{center}
\begin{tabularx}{\linewidth}{|X|X|X|}
\hline
\rowcolor{popblueL}
\textbf{Actors / Les acteurs} & \textbf{Tools \& Process / Outils \& processus} & \textbf{Verbs / Verbes d'action} \\
\hline
Pollster (sondeur) & Questionnaire (questionnaire) & Conduct a poll / survey (mener un sondage) \\
Respondent (personne interrogée) & Ballot / Ballot paper (bulletin de vote) & Carry out a survey (réaliser une enquête) \\
Sample / Sample size (échantillon / taille) & Margin of error (marge d'erreur) & Release / Publish results (publier les résultats) \\
Target population (population cible) & Bias / Biased sample (biais / échantillon biaisé) & Predict / Forecast (prédire / prévoir) \\
Interviewer (enquêteur) & Random sampling (échantillonnage aléatoire) & Commission a poll (commanditer un sondage) \\
Exit pollster (sondeur à la sortie des urnes) & Turnout (taux de participation) & Weight the data (pondérer les données) \\
\hline
\end{tabularx}
\end{center}

\vspace{0.5em}

\begin{center}
\begin{tabularx}{\linewidth}{|X|X|X|}
\hline
\rowcolor{poporangeL}
\textbf{Results \& Analysis / Résultats \& analyse} & \textbf{Adjectives / Adjectifs utiles} & \textbf{Idioms \& Collocations / Expressions} \\
\hline
Outcome (résultat, issue) & Representative (représentatif) & A straw poll (sondage informel, test) \\
Landslide (victoire écrasante) & Reliable / Unreliable (fiable / non fiable) & A snap poll (sondage éclair, instantané) \\
Swing (variation, bascule) & Accurate / Inaccurate (précis / imprécis) & Neck and neck (au coude-à-coude) \\
Rating / Approval rating (cote de popularité) & Nationwide (national) & Too close to call (trop serré pour trancher) \\
Trend (tendance) & Biased / Unbiased (biaisé / impartial) & A landslide victory (victoire écrasante) \\
Projection (projection, estimation) & Volatile (volatil, instable) & To gauge public opinion (mesurer l'opinion) \\
\hline
\end{tabularx}
\end{center}

\courssub{Points de grammaire \& structures à maîtriser}

\begin{center}
\begin{tabularx}{\linewidth}{|X|X|X|}
\hline
\rowcolor{popgreenL}
\textbf{Structure} & \textbf{Règle / Emploi} & \textbf{Exemple traduit} \\
\hline
\eng{Poll / Survey \textbf{on} + topic} & Préposition \textbf{on} pour le sujet du sondage. & \eng{A poll on education.} « Un sondage sur l'éducation. » \\
\eng{Survey \textbf{of} + population} & Préposition \textbf{of} pour la population étudiée. & \eng{A survey of voters.} « Une enquête auprès des électeurs. » \\
\eng{Margin \textbf{of} error} & Nom composé figé : \textbf{of}. & \eng{A margin of error of 3\%.} « Une marge d'erreur de 3 \%. » \\
\eng{Predict / Forecast \textbf{that} + clause} & Verbes suivis de \textbf{that} (omissible). & \eng{They predict (that) he will win.} « Ils prédisent qu'il gagnera. » \\
\eng{Representative \textbf{of}} & Adjectif suivi de \textbf{of}. & \eng{The sample is representative of the youth.} « L'échantillon est représentatif des jeunes. » \\
\eng{According \textbf{to} + source} & Pour citer la source (jamais \eng{according to me}). & \eng{According to the latest poll...} « Selon le dernier sondage... » \\
\hline
\end{tabularx}
\end{center}

\courssub{Méthodes express}

\begin{methode}[Décrire les résultats d'un sondage]
\begin{enumerate}
\item Identifier la source : \eng{According to a poll conducted by [Institute]...}
\item Préciser l'échantillon : \eng{...based on a sample of [number] respondents...}
\item Citer la marge d'erreur : \eng{...with a margin of error of [X] percent.}
\item Annoncer la tendance : \eng{The poll shows / reveals / indicates that...}
\item Utiliser le vocabulaire de variation : \eng{A swing of [X] points towards...}, \eng{A landslide for...}, \eng{Neck and neck.}
\end{enumerate}
\end{methode}

\begin{methode}[Rédiger une question de questionnaire (Questionnaire Design)]
\begin{enumerate}
\item \textbf{Closed question (Question fermée) :} \eng{Do you approve of the government's action? (Yes / No / No opinion)}
\item \textbf{Open question (Question ouverte) :} \eng{What is the main issue facing your community?}
\item \textbf{Likert scale (Échelle de Likert) :} \eng{Strongly agree / Agree / Neutral / Disagree / Strongly disagree.}
\item \textbf{Attention au biais :} Éviter \eng{Leading questions} (questions orientées) : \eng{\textbf{Don't you think} the president is doing a good job?} $\rightarrow$ Préférer \eng{How would you rate the president's performance?}
\end{enumerate}
\end{methode}



\begin{aretenir}[Checklist avant le Bac]
$\square$ Je sais traduire \cle{pollster}, \cle{respondent}, \cle{sample}, \cle{margin of error}, \cle{turnout}, \cle{landslide}, \cle{swing}.
$\square$ Je distingue \cle{poll} (sondage d'opinion, vote) de \cle{survey} (enquête plus large, étude) et \cle{census} (recensement exhaustif).
$\square$ J'utilise les bonnes prépositions : \eng{a poll \textbf{on} a topic}, \eng{a survey \textbf{of} a population}, \eng{a margin \textbf{of} error}.
$\square$ Je connais les faux amis : \eng{Enquiry / Inquiry} (demande d'information, enquête officielle) $\neq$ \eng{Survey} ; \eng{Bias} (biais statistique) $\neq$ \eng{Prejudice} (préjugé).
$\square$ Je maîtrise les familles de mots : \eng{Poll $\rightarrow$ Pollster}, \eng{Respond $\rightarrow$ Respondent / Response}, \eng{Predict $\rightarrow$ Prediction / Predictable}, \eng{Analyse $\rightarrow$ Analysis / Analyst}.
$\square$ Je sais décrire une évolution : \eng{A sharp rise/fall}, \eng{A slight increase/decrease}, \eng{To remain stable / unchanged}.
$\square$ Je sais exprimer l'incertitude : \eng{Too close to call}, \eng{Within the margin of error}, \eng{Neck and neck}.
$\square$ Je sais citer une source correctement : \eng{According to [Source], ...} (pas \eng{According to me} pour donner mon avis $\rightarrow$ \eng{In my opinion}).
$\square$ Je prononce correctement : \eng{Questionnaire} [\eng{kwes-chuh-nair} « kwes-tchou-nère »], \eng{Margin} [\eng{mar-jin} « mar-djin »], \eng{Bias} [\eng{bai-uhs} « baï-èss »], \eng{Turnout} [\eng{ter-nout} « tèr-naout »].
$\square$ Je suis capable de rédiger 3 types de questions de questionnaire (fermée, ouverte, échelle) sans question orientée (\eng{leading question}).
$\square$ Je connais les expressions idiomatiques : \eng{A straw poll}, \eng{A snap poll}, \eng{Exit poll}.
$\square$ Je sais analyser un graphique de sondage en anglais (bar chart, pie chart, trend line) en utilisant le lexique de la leçon.
\end{aretenir}

\findecours

\end{document}
